% Leçon 27 — Expression écrite : gestion de l'argent de poche et caractères complexes du budget — Chinois Tle A — Cours signé OnBuch+
% Programme officiel MINESEC. Compiler avec : tectonic ZH27-ecrit-gestion-de-l-argent-de-poche-et-caracteres-c.tex
\newcommand{\DOCMATIERE}{Chinois}
\newcommand{\DOCNIVEAU}{Tle A}
\newcommand{\DOCMODULE}{Module 4 : Activités économiques et monde du travail}
\newcommand{\DOCLECON}{ZH27}
\newcommand{\DOCTITRE}{Leçon 27 — Expression écrite : gestion de l'argent de poche et caractères complexes du budget}
% OnBuch+ — préambule « cours signés » ENRICHI (compilé avec tectonic / XeLaTeX)
% Système visuel « Pop » d'OnBuch+ : fond crème (jamais blanc), bordures encre
% épaisses, ombres « dures » décalées, titres Archivo Black en capitales.
% Version MATHÉMATIQUES : graphes (pgfplots/tikz), boîtes pédagogiques
% (définition, propriété, méthode, attention, le savais-tu, exercice résolu,
% exercice, TP, bilan), page de garde et signature OnBuch+ ancrée en bas.
%
% Macros attendues dans main.tex AVANT \input{preamble} :
%   \DOCMATIERE  \DOCNIVEAU  \DOCMODULE  \DOCLECON (numéro)  \DOCTITRE
\documentclass[11pt]{article}

\usepackage{fontspec}
\usepackage[french]{babel} % français = langue principale ; le chinois via \zh{..} / textebox (xeCJK)
\usepackage{xeCJK}
\usepackage{pifont} % \ding{51} \ding{55} utilisés par les modèles
\usepackage[a4paper,margin=2cm,top=3.4cm,bottom=2.8cm,headheight=1.4cm,footskip=1.3cm]{geometry}
\usepackage{amsmath,amssymb,amsfonts}
\usepackage{mathtools} % \xmapsto, \coloneqq... souvent utilisés par les modèles
\usepackage{cancel} % simplifications barrées (bilans de forces)
% \abs{z} (module, valeur absolue) et \norm{u} (norme), utilisés par les modèles
\providecommand{\abs}{}\let\abs\relax\DeclarePairedDelimiter{\abs}{\lvert}{\rvert}
\providecommand{\norm}{}\let\norm\relax\DeclarePairedDelimiter{\norm}{\lVert}{\rVert}
\usepackage[table]{xcolor} % [table] : fournit aussi \rowcolors (alternance de lignes)
\usepackage{pagecolor}
\usepackage{tikz}
\usetikzlibrary{arrows.meta,positioning,calc,shapes.geometric,shapes.misc,shapes.arrows,decorations.pathmorphing,decorations.pathreplacing,decorations.markings,patterns,shadows,mindmap,trees,fit,backgrounds,matrix,intersections,angles,quotes}
\usepackage{pgfplots}
\usepackage[version=4]{mhchem} % équations nucléaires \ce{^{235}_{92}U}
\usepackage[european,siunitx]{circuitikz} % schémas électriques
\pgfplotsset{compat=1.18}
\usepgfplotslibrary{fillbetween,groupplots} % aires entre courbes (calcul intégral)
\usepackage{enumitem}
\usepackage{fancyhdr}
% [most] charge toutes les bibliothèques tcolorbox (listings, minted, raster,
% documentation...) et gonfle inutilement la mémoire TeX sur un document long
% (~300 boîtes) : on ne charge que ce qui est réellement utilisé.
\usepackage[skins,breakable]{tcolorbox}
\usepackage{graphicx}
\usepackage{colortbl}
\usepackage{booktabs}
\usepackage{tabularx}
\usepackage{diagbox} % case d'en-tête diagonale des tableaux à double entrée
% Colonnes extensibles centrée / gauche / droite (C, L, R), souvent utilisées par les modèles
\newcolumntype{C}{>{\centering\arraybackslash}X}
\newcolumntype{L}{>{\raggedright\arraybackslash}X}
\newcolumntype{R}{>{\raggedleft\arraybackslash}X}
\usepackage{array}
\usepackage{multicol}
\usepackage{multirow}
\usepackage{siunitx}
% parse-numbers=false : accepte aussi \SI{2,5 \times 10^{-3}}{...}
\sisetup{locale=FR,output-decimal-marker={,},group-minimum-digits=4,parse-numbers=false}
\DeclareSIUnit\ohm{\ensuremath{\symup{\Omega}}} % U+2126 absent de la police
\DeclareSIUnit\division{div} % réglages d'oscilloscope (ms/div, V/div)
\DeclareSIUnit\tr{tr} % tours (vitesse de rotation, ex. \SI{1500}{\tr\per\minute})
\DeclareSIUnit\molar{M} % concentration molaire (ex. \SI{3}{\molar}, \SI{1}{\micro\molar})
\DeclareSIUnit\an{an} % année (le modèle confond parfois avec \year, un compteur TeX) — ex. \SI{5}{\kilo\meter\per\an}
\DeclareSIUnit\FCFA{FCFA} % franc CFA (ex. \SI{500}{\FCFA\per\kilo\gram})
\DeclareSIUnit\degreeBrix{\textdegree Brix} % teneur en sucre d'un sirop/jus (réfractométrie)
\DeclareSIUnit\month{mois} % le modèle utilise parfois \month, un compteur TeX, comme unité de durée
\usepackage{qrcode}
% \degree : commande fréquemment utilisée par les modèles pour un angle ou une
% température en dehors de \SI{}{} (non fournie par les paquets chargés).
\providecommand{\degree}{^{\circ}}
% \qed (fin de démonstration), utilisé par les modèles sans amsthm
\providecommand{\qed}{\hfill$\square$}
% \barre : double barre des valeurs interdites dans un tableau de variations (tkz-tab)
\providecommand{\barre}{\ensuremath{\|}}
% environnement proof (amsthm non chargé) utilisé par les modèles
\ifdefined\proof\else\newenvironment{proof}[1][Démonstration]{\par\noindent\textit{#1.}\ }{\qed\par}\fi
\usepackage{lastpage}
\usepackage{hyperref}
\hypersetup{hidelinks}

% ============================================================
% Palette « Pop » OnBuch+ — mêmes hex que la classe P (plus_ui.dart)
% ============================================================
\definecolor{poporange}{HTML}{F59321}
\definecolor{popdark}{HTML}{E07A0C}
\definecolor{popcream}{HTML}{FFF6E8}
\definecolor{popink}{HTML}{1C1714}
\definecolor{poppurple}{HTML}{7A5AE0}
\definecolor{popgreen}{HTML}{2E9E6A}
\definecolor{poppink}{HTML}{E0457B}
\definecolor{popblue}{HTML}{2F7DE1}
\definecolor{popgold}{HTML}{C98A12}
\definecolor{popmuted}{HTML}{8A7F76}
\definecolor{popline}{HTML}{EADFD2}
\definecolor{popgray}{HTML}{8A7F76}
\colorlet{popgrey}{popgray}
% Alias tolérants (le modèle nomme parfois les couleurs différemment)
\colorlet{popwhite}{white}
\colorlet{popblack}{popink}
\colorlet{popred}{poppink}
\colorlet{popyellow}{popgold}
% Teintes claires (fonds de tableaux/encadrés) — le modèle nomme parfois
% les couleurs avec un suffixe L (ex. popblueL) sans les avoir définies.
\colorlet{poporangeL}{poporange!20}
\colorlet{popdarkL}{popdark!20}
\colorlet{poppurpleL}{poppurple!20}
\colorlet{popgreenL}{popgreen!20}
\colorlet{poppinkL}{poppink!20}
\colorlet{popblueL}{popblue!20}
\colorlet{popgoldL}{popgold!20}
\colorlet{popredL}{popred!20}
\colorlet{popgrayL}{popgray!20}
% Teintes claires pour fonds de tableaux / zones
\colorlet{popblueL}{popblue!10!white}
\colorlet{poporangeL}{poporange!14!white}
\colorlet{popgreenL}{popgreen!12!white}
\colorlet{poppurpleL}{poppurple!12!white}
\colorlet{poppinkL}{poppink!12!white}
\colorlet{popgoldL}{popgold!14!white}

\pagecolor{popcream}

% ============================================================
% Typographie
% ============================================================
\defaultfontfeatures{Ligatures=TeX}
\setmainfont{PlusJakartaSans-Medium.ttf}[Path=../../../fonts/,BoldFont=PlusJakartaSans-Bold.ttf]
% Symboles, API et emojis absents de la police principale : repli sur DejaVu Sans (caractères actifs)
\newfontface\symfont{DejaVuSans.ttf}[Path=../../../fonts/]
\makeatletter
\newcommand\symactive[1]{\catcode#1=13 \begingroup\lccode`\~=#1 \lowercase{\endgroup\def~}{{\symfont\char#1}}}
\newcommand\symblank[1]{\catcode#1=13 \begingroup\lccode`\~=#1 \lowercase{\endgroup\def~}{}}
\newcommand\symhyph[1]{\catcode#1=13 \begingroup\lccode`\~=#1 \lowercase{\endgroup\def~}{\mbox{-}}}
\newcount\symc
\newcommand\symrange[2]{\symc=#1 \@whilenum\symc<#2 \do{\expandafter\symactive\expandafter{\the\symc}\advance\symc by 1 }}
\newcommand\symblankrange[2]{\symc=#1 \@whilenum\symc<#2 \do{\expandafter\symblank\expandafter{\the\symc}\advance\symc by 1 }}
\makeatother
\symrange{"0250}{"02FF}\symactive{"2713}\symactive{"2714}\symactive{"2717}\symactive{"2718}\symactive{"2610}\symactive{"2611}\symactive{"2612}
\symrange{"2600}{"27BF}
\catcode"2705=13 \begingroup\lccode`\~="2705 \lowercase{\endgroup\def~}{{\symfont\char"2713}}\symblank{"26BD}\symblank{"26A1}
\symrange{"2460}{"257F}\symactive{"223C}\symactive{"2248}\symactive{"2260}
\symactive{"01F8}\symactive{"01F9}\symactive{"1E3E}\symactive{"1E3F}
\symhyph{"2011}\symhyph{"2E1A}\symblankrange{"1F300}{"1FAFF}\symblank{"FE0F}
% Caractères chinois : WenQuanYi Zen Hei (sous-ensemble GB2312 + pinyin), gras/italique simulés
\setCJKmainfont{WenQuanYiZenHei-subset.ttf}[Path=../../../fonts/,AutoFakeBold=3,AutoFakeSlant=0.2]
\xeCJKsetup{CJKspace=false}
% Pinyin : voyelles à caron (ǎ ǐ ǒ ǔ ǖ ǘ ǚ ǜ) absentes de la police principale -> caractères actifs
\newfontface\pyfont{WenQuanYiZenHei-subset.ttf}[Path=../../../fonts/]
\newcommand\pyactive[1]{\catcode#1=13 \begingroup\lccode`\~=#1 \lowercase{\endgroup\def~}{{\pyfont\char#1}}}
\pyactive{"01CD}\pyactive{"01CE}\pyactive{"01CF}\pyactive{"01D0}\pyactive{"01D1}\pyactive{"01D2}\pyactive{"01D3}\pyactive{"01D4}\pyactive{"01D5}\pyactive{"01D6}\pyactive{"01D7}\pyactive{"01D8}\pyactive{"01D9}\pyactive{"01DA}\pyactive{"01DB}\pyactive{"01DC}
% Maths en sans-serif (Fira Math) avec chiffres et lettres droites de Plus
% Jakarta, pour que les formules se fondent dans le texte.
\usepackage[math-style=ISO,bold-style=ISO]{unicode-math}
\setmathfont{FiraMath-Regular.otf}[Path=../../../fonts/]
\setmathfont{PlusJakartaSans-Medium.ttf}[Path=../../../fonts/,range={up/num,up/{Latin,latin},\mathup}]
\usepackage{icomma} % virgule décimale sans espace en mode math
% Caractères mathématiques Unicode absents des polices de texte (Plus Jakarta,
% Archivo Black) : rendus par leur équivalent mathématique, en texte comme en
% formule — sinon ils s'affichent en carré vide (ex. « Arithmétique dans ℤ »).
\usepackage{newunicodechar}
\newunicodechar{ℝ}{\ensuremath{\mathbb{R}}}
\newunicodechar{ℤ}{\ensuremath{\mathbb{Z}}}
\newunicodechar{ℂ}{\ensuremath{\mathbb{C}}}
\newunicodechar{ℕ}{\ensuremath{\mathbb{N}}}
\newunicodechar{ℚ}{\ensuremath{\mathbb{Q}}}
\newunicodechar{𝔻}{\ensuremath{\mathbb{D}}}
\newunicodechar{α}{\ensuremath{\alpha}}
\newunicodechar{β}{\ensuremath{\beta}}
\newunicodechar{γ}{\ensuremath{\gamma}}
\newunicodechar{δ}{\ensuremath{\delta}}
\newunicodechar{ε}{\ensuremath{\varepsilon}}
\newunicodechar{θ}{\ensuremath{\theta}}
\newunicodechar{λ}{\ensuremath{\lambda}}
\newunicodechar{μ}{\ensuremath{\mu}}
\newunicodechar{σ}{\ensuremath{\sigma}}
\newunicodechar{τ}{\ensuremath{\tau}}
\newunicodechar{φ}{\ensuremath{\varphi}}
\newunicodechar{ψ}{\ensuremath{\psi}}
\newunicodechar{ω}{\ensuremath{\omega}}
\newunicodechar{Σ}{\ensuremath{\Sigma}}
% majuscules produites par \MakeUppercase dans les titres (α-aminés -> Α)
\newunicodechar{Α}{\ensuremath{\alpha}}
\newunicodechar{Β}{\ensuremath{\beta}}
\newunicodechar{Γ}{\ensuremath{\Gamma}}
\newunicodechar{Δ}{\ensuremath{\Delta}}
\newunicodechar{Ε}{\ensuremath{\varepsilon}}
\newunicodechar{Θ}{\ensuremath{\Theta}}
\newunicodechar{Λ}{\ensuremath{\Lambda}}
\newunicodechar{Μ}{\ensuremath{\mu}}
\newunicodechar{Τ}{\ensuremath{\tau}}
\newunicodechar{Φ}{\ensuremath{\Phi}}
\newunicodechar{Ψ}{\ensuremath{\Psi}}
\newunicodechar{Ω}{\ensuremath{\Omega}}
\newunicodechar{↦}{\ensuremath{\mapsto}}
\newunicodechar{∈}{\ensuremath{\in}}
\newunicodechar{∉}{\ensuremath{\notin}}
\newunicodechar{∧}{\ensuremath{\wedge}}
\newunicodechar{⇔}{\ensuremath{\Leftrightarrow}}
\newunicodechar{⟺}{\ensuremath{\Longleftrightarrow}}
\newunicodechar{⇒}{\ensuremath{\Rightarrow}}
\newunicodechar{⟹}{\ensuremath{\Longrightarrow}}
\newunicodechar{∘}{\ensuremath{\circ}}
\newunicodechar{∪}{\ensuremath{\cup}}
\newunicodechar{∩}{\ensuremath{\cap}}
\newunicodechar{≡}{\ensuremath{\equiv}}
\newunicodechar{⊂}{\ensuremath{\subset}}
\newunicodechar{⊥}{\ensuremath{\perp}}
\newunicodechar{∃}{\ensuremath{\exists}}
\newunicodechar{∀}{\ensuremath{\forall}}
\newunicodechar{∛}{\ensuremath{\sqrt[3]{\,}}}
\newunicodechar{∜}{\ensuremath{\sqrt[4]{\,}}}
\newunicodechar{⃗}{\ensuremath{{}^{\to}}}
\newunicodechar{ⁿ}{\ifmmode^{n}\else\textsuperscript{\ensuremath{n}}\fi}
\newunicodechar{ˣ}{\ifmmode^{x}\else\textsuperscript{\ensuremath{x}}\fi}
\newunicodechar{ᵇ}{\ifmmode^{b}\else\textsuperscript{\ensuremath{b}}\fi}
\newunicodechar{ʳ}{\ifmmode^{r}\else\textsuperscript{\ensuremath{r}}\fi}
\newunicodechar{ᵘ}{\ifmmode^{u}\else\textsuperscript{\ensuremath{u}}\fi}
\newunicodechar{ʸ}{\ifmmode^{y}\else\textsuperscript{\ensuremath{y}}\fi}
\newunicodechar{ᵃ}{\ifmmode^{a}\else\textsuperscript{\ensuremath{a}}\fi}
\newunicodechar{ᵉ}{\ifmmode^{e}\else\textsuperscript{\ensuremath{e}}\fi}
\newunicodechar{ᵏ}{\ifmmode^{k}\else\textsuperscript{\ensuremath{k}}\fi}
\newunicodechar{ᵖ}{\ifmmode^{p}\else\textsuperscript{\ensuremath{p}}\fi}
\newunicodechar{ᴺ}{\ifmmode^{N}\else\textsuperscript{\ensuremath{N}}\fi}
\newunicodechar{ᶜ}{\ifmmode^{c}\else\textsuperscript{\ensuremath{c}}\fi}
\newunicodechar{ʰ}{\ifmmode^{h}\else\textsuperscript{\ensuremath{h}}\fi}
\newunicodechar{ᵠ}{\ifmmode^{q}\else\textsuperscript{\ensuremath{q}}\fi}
\newunicodechar{ₙ}{\ifmmode_{n}\else\textsubscript{\ensuremath{n}}\fi}
\newunicodechar{ₐ}{\ifmmode_{a}\else\textsubscript{\ensuremath{a}}\fi}
\newunicodechar{ᵢ}{\ifmmode_{i}\else\textsubscript{\ensuremath{i}}\fi}
\newunicodechar{ᵣ}{\ifmmode_{r}\else\textsubscript{\ensuremath{r}}\fi}
\newunicodechar{ₖ}{\ifmmode_{k}\else\textsubscript{\ensuremath{k}}\fi}
\newunicodechar{ᵦ}{\ifmmode_{\beta}\else\textsubscript{\ensuremath{\beta}}\fi}
\newunicodechar{ₓ}{\ifmmode_{x}\else\textsubscript{\ensuremath{x}}\fi}
\newfontfamily\popdisplay{ArchivoBlack-Regular.ttf}[Path=../../../fonts/]
\newfontfamily\poplogo{SpaceGrotesk-Bold.ttf}[Path=../../../fonts/]
\newfontfamily\popsemi{PlusJakartaSans-SemiBold.ttf}[Path=../../../fonts/]
\newfontfamily\popxbold{PlusJakartaSans-ExtraBold.ttf}[Path=../../../fonts/]
\newfontfamily\popxboldls{PlusJakartaSans-ExtraBold.ttf}[Path=../../../fonts/,LetterSpace=3.0]

\setlist[enumerate]{leftmargin=1.7em,itemsep=5pt,topsep=4pt}
\setlist[itemize]{leftmargin=1.7em,itemsep=4pt,topsep=4pt,label={\color{poporange}\textbullet}}
\setlength{\parindent}{0pt}
\setlength{\parskip}{5pt}
\renewcommand{\arraystretch}{1.25}

% pgfplots : style Pop par défaut
\pgfplotsset{
  every axis/.append style={axis line style={popink,line width=1pt},tick style={popink},
    grid style={popline,line width=0.5pt},label style={font=\small},tick label style={font=\footnotesize}},
  popcurve/.style={poporange,line width=1.8pt,no markers,smooth},
}
% Même style disponible dans un \draw TikZ simple (hors pgfplots)
\tikzset{popcurve/.style={poporange,line width=1.8pt}}
\tikzset{popgrid/.style={popline,very thin}}
\colorlet{popcurve}{poporange} % le nom du style est parfois utilisé comme couleur

% ============================================================
% Pill / squiggle
% ============================================================
\newcommand{\poppill}[3]{%
  \tikz[baseline=-0.55ex]\node[draw=popink,line width=1.2pt,rounded corners=2.6pt,%
    fill=#2,inner xsep=6.5pt,inner ysep=3pt]{%
    \popxboldls\fontsize{7}{7}\selectfont\color{#3}\MakeUppercase{#1}};%
}
\newcommand{\popsquiggle}[1]{%
  \tikz[baseline=-0.4ex]{
    \draw[#1,line width=2.6pt,line cap=round]
      plot [smooth] coordinates {(0,0.09) (0.34,0.01) (0.68,0.21) (1.02,0.01) (1.36,0.10)};
  }%
}
% Mot-clé surligné (marqueur orange sous le texte)
\newcommand{\cle}[1]{{\popxbold\color{popdark}#1}}

% ============================================================
% En-tête / pied
% ============================================================
\pagestyle{fancy}
\fancyhf{}
\renewcommand{\headrulewidth}{0pt}
\renewcommand{\footrulewidth}{0pt}
\lhead{\raisebox{-0.15ex}{\poplogo\fontsize{13}{13}\selectfont\color{popink}OnBuch\color{poporange}+}}
\rhead{\poppill{\DOCMATIERE\ \textbullet\ \DOCNIVEAU\ \textbullet\ Leçon \DOCLECON}{white}{popink}}
\lfoot{\popxbold\fontsize{7.6}{7.6}\selectfont\color{popmuted}\MakeUppercase{Cours signé OnBuch+}\ \textbullet\ {\popsemi\DOCTITRE}}
\rfoot{\tikz[baseline=-0.6ex]\node[rounded corners=8.5pt,fill=popink,minimum height=17pt,minimum width=17pt,inner xsep=5pt,inner sep=0]{\popxbold\fontsize{8}{8}\selectfont\color{popcream}\thepage\,/\,\pageref*{LastPage}};}

% ============================================================
% Titres
% ============================================================
\newcommand{\chaptitre}[2]{%
  \begin{tcolorbox}[enhanced,colback=poporange,colframe=popink,boxrule=2.6pt,arc=11pt,
      drop shadow=popink,left=15pt,right=15pt,top=12pt,bottom=11pt,before skip=3pt,after skip=18pt]
    \poppill{#2}{popink}{popcream}\par\vspace{9pt}
    {\raggedright\hyphenpenalty=10000\exhyphenpenalty=10000\popdisplay\fontsize{22}{24}\selectfont\color{popcream}\MakeUppercase{#1}\par}
  \end{tcolorbox}
}
% Section numérotée : pastille orange avec le numéro + titre Archivo
\newcounter{popsec}
\newcommand{\coursec}[1]{%
  \stepcounter{popsec}\setcounter{popsub}{0}%
  \par\vspace{16pt}\noindent
  \begin{minipage}{\linewidth}
  \tikz[baseline=-0.7ex]\node[draw=popink,line width=1.6pt,fill=poporange,rounded corners=4pt,minimum size=22pt,inner sep=0]{\popdisplay\fontsize{12}{12}\selectfont\color{popcream}\thepopsec};%
  \hspace{8pt}{\popdisplay\fontsize{15}{17}\selectfont\color{popink}\MakeUppercase{#1}}\par
  \vspace{4pt}\noindent{\color{poporange}\rule{36pt}{3.6pt}}
  \end{minipage}\par\nopagebreak\vspace{8pt}
}
\newcounter{popsub}
\newcommand{\courssub}[1]{%
  \stepcounter{popsub}%
  \par\vspace{9pt}\noindent{\popxbold\fontsize{11.5}{13}\selectfont\color{popink}{\color{poporange}\thepopsec.\thepopsub}\hspace{6pt}#1}\par\nopagebreak\vspace{4pt}
}

% ============================================================
% Boîtes pédagogiques — même recette : fond blanc, bordure épaisse colorée,
% ombre dure encre, badge plein. Titre optionnel [#1].
% ============================================================
\tcbset{popbase/.style={breakable,enhanced,colback=white,boxrule=2.3pt,arc=9pt,
  shadow={3pt}{-3pt}{0pt}{popink},left=14pt,right=14pt,top=11pt,bottom=11pt,before skip=11pt,after skip=15pt,
  fontupper=\popsemi\fontsize{10.6}{14.5}\selectfont\color{popink},
  fontlower=\popsemi\fontsize{10.6}{14.5}\selectfont\color{popink}}}
% Coche dessinée (le glyphe ✓ n'existe pas dans les polices du document)
\newcommand{\popcheck}[1]{\tikz[baseline=-0.5ex]\draw[#1,line width=1.6pt,line cap=round,line join=round](-0.13,0.0)--(-0.04,-0.09)--(0.13,0.1);}
\providecommand{\coche}{\popcheck{popgreen}}
\providecommand{\resultat}[1]{\par\smallskip\noindent\fcolorbox{popgreen}{popgreenL}{\parbox{\dimexpr\linewidth-2\fboxsep-2\fboxrule\relax}{\textbf{Résultat.} #1}}\par\smallskip}
\newcommand{\popboxhead}[3]{\poppill{#1}{#2}{white}\ifx\relax#3\relax\else\hspace{6pt}{\popxbold\fontsize{10.6}{12}\selectfont\color{popink}#3}\fi\par\vspace{7pt}}

\newtcolorbox{aretenir}[1][]{popbase,colframe=popgreen,before upper={\popboxhead{À retenir}{popgreen}{#1}}}
% Texte en chinois (document de lecture, dialogue, extrait) : cadre
% d'étude. Un titre optionnel (source, auteur) s'affiche dans l'en-tête.
\newtcolorbox{textebox}[1][]{popbase,colframe=popblue,before upper={\popboxhead{文本 · Texte}{popblue}{#1}},after upper={}}
% Marques ✓ / ✗ (forme correcte / fautive) souvent utilisées par les modèles
\providecommand{\cross}{\textcolor{poppink}{\ensuremath{\times}}}
\providecommand{\xmark}{\cross}\providecommand{\crossmark}{\cross}\providecommand{\cmark}{\ensuremath{\checkmark}}
% Ligne de réponse / justification à compléter (inventée par les modèles)
\providecommand{\justification}[1]{\par\noindent\textit{Justification~:} \dotfill\par}
% \textipa (package tipa non chargé) : la police ne couvre pas l'API -> simple passage
\providecommand{\textipa}[1]{#1}
% Soulignement (ulem non chargé) : \uline, \ul
\providecommand{\uline}[1]{\underline{#1}}\providecommand{\ul}[1]{\underline{#1}}
% \glossaire{\item ...} inventée par les modèles : liste de glossaire
\providecommand{\glossaire}[1]{\begin{itemize}#1\end{itemize}}
% \alert (beamer) inventée par les modèles : texte mis en évidence
\providecommand{\alert}[1]{\textbf{\textcolor{poppink}{#1}}}
% variantes ASCII d'ordinaux écrites en macro
\providecommand{\iere}{\textsuperscript{re}}\providecommand{\ieme}{\textsuperscript{e}}\providecommand{\ere}{\textsuperscript{re}}\providecommand{\eme}{\textsuperscript{e}}\providecommand{\er}{\textsuperscript{er}}\providecommand{\ie}{\textsuperscript{e}}
% Ordinaux français écrits en macro par les modèles : 1\ière, 3\ième
\newcommand{\ière}{\textsuperscript{re}}\newcommand{\ième}{\textsuperscript{e}}
% \exemple (inventée par les modèles) : étiquette « Exemple : »
\providecommand{\exemple}{\textbf{Exemple~:}\ }
% \square utilisable aussi hors mode math (cases à cocher)
\AtBeginDocument{\let\squareMath\square\renewcommand{\square}{\ensuremath{\squareMath}}}
% Mot ou phrase en chinois à l'intérieur d'un paragraphe français
\newcommand{\zh}[1]{\ifmmode\text{#1}\else{#1}\fi}
\let\eng\zh \let\esp\zh \let\all\zh % alias tolérés
% flèches d'intonation inventées par les modèles
\providecommand{\textsearrow}{}\renewcommand{\textsearrow}{\ensuremath{\searrow}}\providecommand{\textnearrow}{}\renewcommand{\textnearrow}{\ensuremath{\nearrow}}
% \fr{..} : mot français (inventé par les modèles)
\providecommand{\fr}[1]{\textit{#1}}
% zhbox : encadré de texte chinois inventé par les modèles
\newenvironment{zhbox}{\begin{quote}}{\end{quote}}
% \courssubsub : titre de niveau 3 inventé par les modèles
\providecommand{\courssubsub}[1]{\par\medskip\noindent\textbf{#1}\par\smallskip}
% \zhbf : texte chinois en gras inventé par les modèles
\providecommand{\zhbf}[1]{\textbf{#1}}
% \vs : « versus » inventé par les modèles
\providecommand{\vs}{\textit{vs}\ }
% \blank : case à compléter inventée par les modèles
\providecommand{\blank}[1][]{\underline{\hspace{1.6em}}}
\newtcolorbox{exemplebox}[1][]{popbase,colframe=poporange,before upper={\popboxhead{Exemple}{poporange}{#1}}}
\newtcolorbox{definition}[1][]{popbase,colframe=popblue,before upper={\popboxhead{Définition}{popblue}{#1}}}
\newtcolorbox{propriete}[1][]{popbase,colframe=poppurple,before upper={\popboxhead{Propriété}{poppurple}{#1}}}
\newtcolorbox{methode}[1][]{popbase,colframe=popgold,before upper={\popboxhead{Méthode}{popgold}{#1}}}
\newtcolorbox{attention}[1][]{popbase,colframe=poppink,before upper={\popboxhead{Attention}{poppink}{#1}}}
\newtcolorbox{experience}[1][]{popbase,colframe=popblue,colback=popblueL,before upper={\popboxhead{Expérience}{popblue}{#1}}}
% « Le savais-tu ? » : carte violette pleine (contexte, culture, Cameroun)
\newtcolorbox{savaistu}[1][]{popbase,colback=poppurple,colframe=popink,
  fontupper=\popsemi\fontsize{10.4}{14.2}\selectfont\color{white},
  before upper={\poppill{Le savais-tu\,?}{white}{poppurple}\ifx\relax#1\relax\else\hspace{6pt}{\popxbold\color{white}#1}\fi\par\vspace{7pt}}}
% Exercice résolu : énoncé en haut, \tcblower puis la solution sur fond crème
\newcounter{popexo}\newcounter{popexr}
\newtcolorbox{exoresolu}[1][]{popbase,colframe=poporange,colbacklower=popcream,
  segmentation style={popink,line width=1.2pt,dashed},
  before upper={\stepcounter{popexr}\popboxhead{Exercice résolu \thepopexr}{poporange}{#1}},
  before lower={\poppill{Solution}{popink}{popcream}\par\vspace{6pt}}}
% Exercice (non résolu en place) — la correction est dans la partie Corrigés
\newtcolorbox{exercice}[1][]{popbase,colframe=popink,
  before upper={\stepcounter{popexo}\popboxhead{Exercice \thepopexo}{popink}{#1}}}
% Corrigé
\newtcolorbox{corrige}[1][]{popbase,colframe=popgreen,colback=popgreenL,
  before upper={\popboxhead{Corrigé}{popgreen}{#1}}}
% Objectifs / prérequis / bilan
\newtcolorbox{objectifs}{popbase,colframe=popink,colback=poporangeL,
  before upper={\popboxhead{Objectifs de la leçon}{popink}{}}}
\newtcolorbox{prerequis}{popbase,colframe=popink,colback=popblueL,
  before upper={\popboxhead{Prérequis}{popblue}{}}}
\newtcolorbox{bilan}{popbase,colframe=popink,colback=white,boxrule=2.8pt,
  before upper={\popboxhead{Fiche bilan}{poporange}{L'essentiel à connaître pour l'examen}}}
% Figure encadrée avec légende
\newtcolorbox{popfigure}[1][]{enhanced,breakable=false,colback=white,colframe=popline,boxrule=1.4pt,arc=8pt,
  left=10pt,right=10pt,top=10pt,bottom=8pt,before skip=10pt,after skip=12pt,halign=center,
  lower separated=false,halign lower=center,
  fontlower=\popsemi\fontsize{9}{11.5}\selectfont\color{popmuted},
  #1}
\newcommand{\legende}[1]{\par\vspace{6pt}{\popsemi\fontsize{9}{11.5}\selectfont\color{popmuted}#1}}

% ============================================================
% Signature OnBuch+
% ============================================================
\newcommand{\poplogoinline}[1]{{\poplogo\fontsize{#1}{#1}\selectfont\color{popink}OnBuch\color{poporange}+}}
\newcommand{\signatureonbuch}{%
  \begin{tcolorbox}[enhanced,colback=popink,colframe=popink,boxrule=2.2pt,arc=10pt,
      drop shadow=poporange,left=14pt,right=14pt,top=10pt,bottom=10pt,before skip=0pt,after skip=0pt]
    \tikz[baseline=-0.7ex]\node[circle,fill=popgreen,draw=popcream,line width=1.3pt,minimum size=22pt,inner sep=0]{\popcheck{white}};%
    \hspace{9pt}\begin{minipage}[c]{0.62\linewidth}
      {\popxbold\fontsize{12}{13}\selectfont\color{popcream}Signé\ \textbullet\ {\poplogo OnBuch\color{poporange}+}}\par\vspace{2pt}
      {\raggedright\popsemi\fontsize{8}{10.5}\selectfont\color{popline}\MakeUppercase{Équipe pédagogique OnBuch+}\\\MakeUppercase{Conforme au programme officiel MINESEC}\par}
    \end{minipage}\hfill
    {\poplogo\fontsize{22}{22}\selectfont\color{popcream}OnBuch\color{poporange}+}
  \end{tcolorbox}%
}

% ============================================================
% Page de garde
% #1 titre · #2 matière · #3 niveau · #4 module · #5 n° leçon · #6 durée ·
% #7 description / objectifs (texte ou itemize)
% ============================================================
\newcommand{\pagegarde}[7]{%
  \thispagestyle{empty}
  \newgeometry{a4paper,margin=1.7cm,top=1.6cm,bottom=1.4cm}%
  \begin{tikzpicture}[remember picture,overlay]
    \fill[poporange!18] ([xshift=-3.2cm,yshift=-3.6cm]current page.north east) circle (4.6cm);
    \fill[poppurple!14] ([xshift=2.4cm,yshift=7.6cm]current page.south west) circle (3.8cm);
  \end{tikzpicture}%
  {\poplogo\fontsize{30}{30}\selectfont\color{popink}OnBuch\color{poporange}+}\hfill
  \raisebox{0.6ex}{\poppill{Cours signé \textbullet\ Premium}{popink}{popcream}}\par
  \vspace{1pt}\popsquiggle{poppurple}\par
  \vspace{8pt}
  \begin{tcolorbox}[enhanced,colback=poporange,colframe=popink,boxrule=2.8pt,arc=13pt,
      drop shadow=popink,left=18pt,right=18pt,top=12pt,bottom=13pt,after skip=10pt]
    \poppill{#2 \textbullet\ #3}{popink}{popcream}\hspace{5pt}\poppill{#4}{white}{popink}\par\vspace{8pt}
    {\popdisplay\fontsize{14}{14}\selectfont\color{popink}LEÇON #5}\par\vspace{4pt}
    {\raggedright\hyphenpenalty=10000\exhyphenpenalty=10000\popdisplay\fontsize{25}{27}\selectfont\color{popcream}\MakeUppercase{#1}\par}
    \vspace{8pt}
    {\popxbold\fontsize{9.5}{9.5}\selectfont\color{popink}\MakeUppercase{Durée conseillée : #6}}
  \end{tcolorbox}
  \begin{tcolorbox}[enhanced,colback=white,colframe=popink,boxrule=2.2pt,arc=10pt,
      drop shadow=popink,left=16pt,right=16pt,top=10pt,bottom=10pt,after skip=8pt]
    \poppill{Dans cette leçon}{poppurple}{white}\par\vspace{5pt}
    \popsemi\fontsize{9.3}{12.6}\selectfont\color{popink}#7
  \end{tcolorbox}
  \noindent\begin{minipage}[t]{0.487\linewidth}
    \begin{tcolorbox}[enhanced,colback=popgreen,colframe=popink,boxrule=2.2pt,arc=10pt,
        drop shadow=popink,left=11pt,right=11pt,top=10pt,bottom=10pt,height=3.1cm,valign=center]
      \tcbox[colback=white,colframe=popink,boxrule=1.3pt,arc=5pt,boxsep=0pt,nobeforeafter,
        left=3pt,right=3pt,top=3pt,bottom=3pt,valign=center]{\qrcode[height=1.9cm]{https://play.google.com/store/apps/details?id=com.luvvix.onbuch&hl=fr}}%
      \hspace{8pt}\begin{minipage}[c]{0.5\linewidth}
        {\popdisplay\fontsize{11}{12.5}\selectfont\color{white}\MakeUppercase{Télécharge OnBuch}}\par\vspace{4pt}
        {\raggedright\popsemi\fontsize{8.4}{11}\selectfont\color{white}Gratuit \textbullet\ Google Play\\Tuteur IA, annales, concours, résultats\par}
      \end{minipage}
    \end{tcolorbox}
  \end{minipage}\hfill
  \begin{minipage}[t]{0.487\linewidth}
    \begin{tcolorbox}[enhanced,colback=poppurple,colframe=popink,boxrule=2.2pt,arc=10pt,
        drop shadow=popink,left=11pt,right=11pt,top=10pt,bottom=10pt,height=3.1cm,valign=center]
      \tikz[baseline=-0.7ex]\node[circle,fill=white,draw=popink,line width=1.3pt,minimum size=1.6cm,inner sep=0]{\popdisplay\fontsize{24}{24}\selectfont\color{poppurple}+};%
      \hspace{8pt}\begin{minipage}[c]{0.6\linewidth}
        {\popdisplay\fontsize{11}{12.5}\selectfont\color{white}\MakeUppercase{Passe à OnBuch+}}\par\vspace{4pt}
        {\raggedright\popsemi\fontsize{8.4}{11}\selectfont\color{white}Tous les cours signés, Tuteur IA illimité, fiches PDF — onglet OnBuch+ de l'app\par}
      \end{minipage}
    \end{tcolorbox}
  \end{minipage}
  \vspace{8pt}
  \popcontenu
  \vfill
  \signatureonbuch
  \restoregeometry
  \newpage
}

% Rangée « Ce cours contient » (page de garde)
\newcommand{\poptile}[3]{% #1 couleur #2 gros texte #3 légende
  \begin{tcolorbox}[enhanced,colback=#1,colframe=popink,boxrule=2pt,arc=9pt,shadow={2.5pt}{-2.5pt}{0pt}{popink},
      left=6pt,right=6pt,top=9pt,bottom=9pt,halign=center,valign=center,height=1.75cm,width=\linewidth,nobeforeafter]
    {\popdisplay\fontsize{12.5}{13}\selectfont\color{white}\MakeUppercase{#2}}\par\vspace{3pt}
    {\centering\popsemi\fontsize{7.8}{9.5}\selectfont\color{white}#3\par}
  \end{tcolorbox}}
\newcommand{\popcontenu}{%
  \par\noindent{\popxboldls\fontsize{7.5}{7.5}\selectfont\color{popmuted}CE COURS SIGNÉ CONTIENT}\par\vspace{7pt}
  \noindent\begin{minipage}[t]{0.235\linewidth}\poptile{popblue}{Cours}{complet, illustré, pas à pas}\end{minipage}\hfill
  \begin{minipage}[t]{0.235\linewidth}\poptile{poporange}{Méthodes}{et exercices résolus}\end{minipage}\hfill
  \begin{minipage}[t]{0.235\linewidth}\poptile{poppink}{Exercices}{type examen + corrigés}\end{minipage}\hfill
  \begin{minipage}[t]{0.235\linewidth}\poptile{popgreen}{Fiche bilan}{l'essentiel à retenir}\end{minipage}\par}

% Sommaire
\newtcolorbox{sommaire}{enhanced,colback=white,colframe=popink,boxrule=2.2pt,arc=10pt,
  drop shadow=popink,left=18pt,right=18pt,top=13pt,bottom=13pt,before skip=6pt,after skip=20pt,
  before upper={\poppill{Sommaire}{poppurple}{white}\par\vspace{9pt}\popsemi\fontsize{10.6}{14.5}\selectfont\color{popink}}}

% Signature de fin de cours — poussée tout en bas de la dernière page
\newcommand{\findecours}{%
  \par\vfill\nopagebreak
  \begin{center}\popsquiggle{poporange}\end{center}\vspace{4pt}
  \signatureonbuch
}
\AtBeginDocument{\def\llbracket{\mathopen{[\mkern-3.5mu[}}\def\rrbracket{\mathclose{]\mkern-3.5mu]}}} % intervalles d'entiers (glyphe ⟦ absent de la police)
% Glyphes absents de Fira Math : redéfinis à partir de glyphes disponibles
\AtBeginDocument{%
  \def\setminus{\mathbin{\backslash}}\let\smallsetminus\setminus
  \def\vdots{\mathord{\vbox{\baselineskip3.2pt\lineskiplimit0pt\kern4pt\hbox{.}\hbox{.}\hbox{.}}}}%
  \def\ddots{\mathinner{\mkern1mu\raise6.5pt\hbox{.}\mkern2mu\raise3.6pt\hbox{.}\mkern2mu\raise0.7pt\hbox{.}\mkern1mu}}%
  \def\triangle{\mathord{\scalebox{0.9}{$\symup{\Delta}$}}}%
  \def\nabla{\mathord{\raisebox{\depth}{\scalebox{1}[-1]{$\symup{\Delta}$}}}}%
  \def\checkmark{\popcheck{popdark}}%
  \def\star{\mathbin{\ast}}\def\bigstar{\mathord{\ast}}%
  \def\diamond{\mathbin{\scalebox{0.6}{$\Diamond$}}}%
  \def\bigcup{\mathop{\mathchoice{\raisebox{-0.3em}{\scalebox{1.7}{$\cup$}}}{\scalebox{1.25}{$\cup$}}{\cup}{\cup}}\displaylimits}%
  \def\bigcap{\mathop{\mathchoice{\raisebox{-0.3em}{\scalebox{1.7}{$\cap$}}}{\scalebox{1.25}{$\cap$}}{\cap}{\cap}}\displaylimits}%
  \def\lesseqqgtr{\mathrel{\lessgtr}}\def\bigoplus{\mathop{\oplus}\displaylimits}%
}
\providecommand{\ssi}{si et seulement si} % abréviation fréquente
\AtBeginDocument{\providecommand{\wideparen}{\overparen}} % arc de cercle

%% FIGFIX-BEGIN v5 — correctifs globaux des figures (docs/onbuch-generation/tools/figfix)
\usepackage{adjustbox}\usepackage{etoolbox}\tcbuselibrary{raster}
% toute figure TikZ/pgfplots plus large que la ligne est réduite à la largeur disponible
\BeforeBeginEnvironment{tikzpicture}{\begin{adjustbox}{max width=\linewidth}}
\AfterEndEnvironment{tikzpicture}{\end{adjustbox}}
% légendes pgfplots lisibles par-dessus les courbes
\pgfplotsset{every axis/.append style={legend style={fill=white,fill opacity=0.92,text opacity=1,draw=black!15,inner sep=3pt}}}
% étiquettes d'arêtes (syntaxe edge["..."]) avec fond blanc
\tikzset{every edge quotes/.append style={fill=white,inner sep=1.2pt}}
%% FIGFIX-END
\begin{document}

\pagegarde{Leçon 27 — Expression écrite : gestion de l'argent de poche et caractères complexes du budget}{Chinois}{Tle A}{Module 4 : Activités économiques et monde du travail}{ZH27}{2 h}{Cette leçon d'expression écrite apprend à rédiger un texte simple en chinois sur la gestion de l'argent de poche et le budget. Elle entraîne l'écriture correcte des caractères complexes du thème (钱, 费, 算, 储蓄...) et la traduction thème/version de phrases courtes.
\vspace{4pt}
\begin{multicols}{2}\small
\begin{itemize}
\item À la fin de cette leçon, je sais écrire correctement les caractères complexes du thème de l'argent (钱, 费, 算, 储蓄, 预算) en respectant radicaux et ordre des traits
\item je sais distinguer les caractères proches (钱/线, 费/佛, 买/卖, 借/错)
\item je sais utiliser le vocabulaire du budget : 零花钱, 花钱, 省钱, 存钱, 预算, 花费
\item je sais structurer un texte simple avec les connecteurs 首先, 然后, 但是, 所以, 因为
\item je sais rédiger un texte de 6 à 10 phrases sur ma gestion de l'argent de poche
\item je sais traduire du français vers le chinois des phrases simples sur le budget (thème)
\end{itemize}\end{multicols}}

\begin{sommaire}
\begin{multicols}{2}
\begin{enumerate}
\item Modèle de texte commenté
\item Structure et connecteurs du texte
\item Écriture correcte des caractères du budget
\item Lexique actif et collocations du thème
\item Thème : français → chinois
\item Version : chinois → français et grille d'évaluation
\item Activité d'intégration
\item Méthodes et exercices résolus
\item Exercices
\item Corrigés des exercices
\item Fiche bilan
\end{enumerate}
\end{multicols}
\end{sommaire}

\begin{objectifs}
\begin{itemize}
\item À la fin de cette leçon, je sais écrire correctement les caractères complexes du thème de l'argent (钱, 费, 算, 储蓄, 预算) en respectant radicaux et ordre des traits
\item je sais distinguer les caractères proches (钱/线, 费/佛, 买/卖, 借/错)
\item je sais utiliser le vocabulaire du budget : 零花钱, 花钱, 省钱, 存钱, 预算, 花费
\item je sais structurer un texte simple avec les connecteurs 首先, 然后, 但是, 所以, 因为
\item je sais rédiger un texte de 6 à 10 phrases sur ma gestion de l'argent de poche
\item je sais traduire du français vers le chinois des phrases simples sur le budget (thème)
\item je sais traduire du chinois vers le français un court texte sur l'argent de poche (version)
\item je sais m'auto-évaluer avec une grille : caractères, grammaire, connecteurs, respect du sujet
\end{itemize}
\end{objectifs}

\begin{prerequis}
\begin{itemize}
\item Nombres, mesures et monnaie (块, 毛, 钱) vus en Première
\item Structure de base sujet-verbe-objet et particule 了
\item Verbes modaux 想, 要, 会, 可以
\item Comparatif simple avec 比 (leçons antérieures du module)
\item Connecteurs simples 和, 因为, 所以 déjà rencontrés
\end{itemize}
\end{prerequis}

\newpage

\coursec{Modèle de texte commenté}

À Yaoundé comme à Douala, beaucoup d'élèves de Terminale reçoivent chaque semaine un peu d'argent de poche pour le transport, le beignet du matin ou le crédit téléphonique. Mais à la fin de la semaine, certains n'ont plus rien, tandis que d'autres ont réussi à économiser pour s'acheter un cahier ou des chaussures. En Chine, les lycéens reçoivent aussi de l'argent de poche (\zh{零花钱} \emph{líng huā qián}) et apprennent très tôt à tenir un petit budget. Comment raconter, en chinois écrit, comment tu gères ton argent de poche ? C'est ce que cette leçon va t'apprendre, avec les caractères complexes du budget et les connecteurs pour structurer ton texte.

\courssub{Lecture et compréhension du texte modèle : 《我的零花钱》}

\begin{textebox}{Texte modèle : 《我的零花钱》}
我每个星期从父母那儿得到五千法郎的零花钱。\\
首先，我用一部分钱坐车去学校。\\
然后，我买早饭和学习用品。\\
但是，我不乱花钱。\\
因为我想存钱买一本汉语词典，所以我每个星期存一千法郎。\\
我觉得管理零花钱很重要。
\end{textebox}

\noindent \textbf{Pinyin intégral (ligne par ligne)} \\
\emph{Wǒ měi ge xīngqī cóng fùmǔ nàr dédào wǔqiān fǎláng de líng huā qián.} \\
\emph{Shǒuxiān, wǒ yòng yī bùfen qián zuò chē qù xuéxiào.} \\
\emph{Ránhòu, wǒ mǎi zǎofàn hé xuéxí yòngpǐn.} \\
\emph{Dànshì, wǒ bù luàn huā qián.} \\
\emph{Yīnwèi wǒ xiǎng cún qián mǎi yì běn Hànyǔ cídiǎn, suǒyǐ wǒ měi ge xīngqī cún yìqiān fǎláng.} \\
\emph{Wǒ juéde guǎnlǐ líng huā qián hěn zhòngyào.}

\noindent \textbf{Traduction littérale} \\
« Moi chaque semaine depuis parents là obtenir cinq mille francs \textsc{de} argent de poche. D'abord, moi utiliser une partie argent prendre voiture aller école. Ensuite, moi acheter petit-déjeuner et étude fournitures. Mais, moi pas désordonné dépenser argent. Parce que moi vouloir économiser argent acheter un \textsc{cl} chinois dictionnaire, donc moi chaque semaine économiser mille francs. Moi sentir gérer argent de poche très important. »

\noindent \textbf{Traduction naturelle} \\
« Chaque semaine, je reçois 5 000 FCFA d'argent de poche de mes parents. D'abord, j'utilise une partie de l'argent pour aller à l'école en voiture (taxi/bus). Ensuite, j'achète le petit-déjeuner et des fournitures scolaires. Mais je ne dépense pas mon argent n'importe comment. Parce que je veux économiser pour acheter un dictionnaire de chinois, j'économise 1 000 FCFA chaque semaine. Je trouve que gérer son argent de poche est très important. »

\begin{aretenir}{Repérage de la structure du texte}
Le texte suit une progression logique en \textbf{quatre étapes} typiques d'un récit de gestion budgétaire :
\begin{enumerate}
    \item \textbf{Situation de départ (Revenu)} : \zh{我每个星期从父母那儿得到五千法郎的零花钱} — Source et montant fixe.
    \item \textbf{Dépenses courantes (Sorties obligatoires)} : \zh{首先... 然后...} — Transport, nourriture, fournitures.
    \item \textbf{Épargne / Choix raisonné (Arbitrage)} : \zh{但是... 因为...所以...} — Renonciation aux dépenses superflues (\zh{不乱花钱}) pour un objectif précis (\zh{买汉语词典}).
    \item \textbf{Opinion finale (Réflexion)} : \zh{我觉得管理零花钱很重要} — Généralisation et valeur éducative.
\end{enumerate}
\end{aretenir}

\courssub{Analyse de la structure et des connecteurs logiques}

Les \cle{connecteurs} (连词 \emph{liáncí}) sont les « charnières » qui articulent les phrases et guident le lecteur. Dans ce texte court, cinq connecteurs majeurs assurent la cohérence.

\begin{center}
\begin{tabularx}{\linewidth}{|X|X|X|X|}
\hline
\rowcolor{popblueL}
\textbf{Connecteur} & \textbf{Pinyin} & \textbf{Fonction logique} & \textbf{Exemple du texte} \\
\hline
首先 & shǒuxiān & Énumération : premier élément (chronologique ou ordre d'importance) & \zh{首先，我用一部分钱坐车去学校。} \\
\hline
然后 & ránhòu & Succession : action suivante, étape ultérieure & \zh{然后，我买早饭和学习用品。} \\
\hline
但是 & dànshì & Opposition / Contraste : introduit une idée contraire à l'attente & \zh{但是，我不乱花钱。} \\
\hline
因为... 所以... & yīnwèi... suǒyǐ... & Cause $\rightarrow$ Conséquence (corrélatif inséparable en chinois écrit) & \zh{因为我想存钱买一本汉语词典，所以我每个星期存一千法郎。} \\
\hline
觉得 & juéde & Verbe d'opinion / sentiment subjectif : introduit le commentaire final & \zh{我觉得管理零花钱很重要。} \\
\hline
\end{tabularx}
\end{center}

\begin{propriete}{La corrélation 因为... 所以... (Yīnwèi... suǒyǐ...)}
En français, on peut dire « Parce que je veux..., je... » (subordonnée + principale) ou « Je... parce que... » (principale + subordonnée). En chinois standard écrit, la structure \textbf{因为... 所以...} est \textbf{quasi obligatoire} pour lier explicitement la cause et la conséquence dans une phrase complexe. On ne dit généralement pas *\zh{因为我想买词典，我存钱} sans \zh{所以}.
\begin{itemize}
    \item \textbf{Schéma} : \zh{因为} + Cause (Sujet + Verbe...) \zh{，所以} + Conséquence (Sujet + Verbe...).
    \item \textbf{Ponctuation} : Une virgule chinoise \zh{，} sépare généralement les deux propositions.
    \item \textbf{Exception} : À l'oral familier, on peut omettre \zh{所以} (\zh{因为下雨，我不去了。}), mais à l'examen écrit, exigez toujours la paire complète.
\end{itemize}
\end{propriete}

\begin{exemplebox}{Variations autour de 因为...所以...}
\begin{itemize}
    \item \zh{因为下雨，所以我不去踢球。} (Yīnwèi xià yǔ, suǒyǐ wǒ bù qù tī qiú.) — Parce qu'il pleut, je ne vais pas jouer au foot.
    \item \zh{他因为生病，所以没来上课。} (Tā yīnwèi shēng bìng, suǒyǐ méi lái shàng kè.) — \textbf{Attention} : \zh{因为} peut être placé après le sujet de la première proposition (Sujet + 因为 + Cause, 所以 + Conséquence).
    \item \textbf{Forme courte (sujet commun omis)} : \zh{因为想买词典，所以（我）存钱。} — Le sujet de la seconde proposition peut être omis s'il est identique au premier.
\end{itemize}
\end{exemplebox}

\begin{popfigure}
\begin{tikzpicture}[
    node distance=1.2cm and 1.5cm,
    box/.style={draw=popblue, thick, rounded corners, fill=popblueL, minimum width=3.2cm, minimum height=1.2cm, align=center, font=\small},
    arrow/.style={-Stealth, thick, popdark},
    labelnode/.style={font=\tiny\itshape, text=popmuted}
]
% Nodes
\node[box, align=center] (rev){Revenu\\ \zh{得到零花钱}\\(5 000 FCFA)};
\node[box, right=of rev, align=center] (dep){Dépenses\\ \zh{坐车、买饭、买用品}\\(\zh{首先、然后})};
\node[box, below=of dep, align=center] (sav){Épargne\\ \zh{存钱买词典}\\(\zh{因为...所以...})};
\node[box, left=of sav, align=center] (op){Opinion\\ \zh{管理很重要}\\(\zh{我觉得})};

% Arrows
\draw[arrow] (rev) -- (dep) node[midway, above, labelnode] {utilise une partie};
\draw[arrow] (dep) -- (sav) node[midway, right, labelnode] {mais contrôle};
\draw[arrow] (sav) -- (op) node[midway, below, labelnode] {réflexion};
\draw[arrow] (op) -- (rev) node[midway, left, labelnode] {cycle hebdo};
\end{tikzpicture}
\legende{Schéma du flux logique du texte : Revenu $\rightarrow$ Dépenses $\rightarrow$ Épargne $\rightarrow$ Opinion (cycle hebdomadaire).}
\end{popfigure}

\begin{popfigure}
\begin{tikzpicture}[
    grow=right,
    level 1/.style={sibling distance=2.6cm, level distance=3.2cm},
    level 2/.style={sibling distance=1.1cm, level distance=3.4cm},
    every node/.style={align=center, font=\small},
    edge from parent/.style={draw=popmuted, thick}
]
\node[draw=popblue, thick, rounded corners, fill=popblueL, font=\small\bfseries, align=center]{零花钱\\ \emph{líng huā qián}}
    child { node[draw=poporange, rounded corners, fill=poporangeL, align=center]{得到 \emph{dédào}\\ obtenir}
        child { node {父母 \emph{fùmǔ}} }
        child { node {长辈 \emph{zhǎngbèi}} }
    }
    child { node[draw=poppurple, rounded corners, fill=poppurpleL, align=center]{花 \emph{huā}\\ dépenser}
        child { node {交通 \emph{jiāotōng}} }
        child { node {饮食 \emph{yǐnshí}} }
        child { node {学习用品 \emph{xuéxí yòngpǐn}} }
    }
    child { node[draw=popgreen, rounded corners, fill=popgreenL, align=center]{存 \emph{cún}\\ épargner}
        child { node {不乱花钱 \emph{bù luàn huā qián}} }
        child { node {目标 \emph{mùbiāo}：买词典} }
    }
    child { node[draw=poppink, rounded corners, fill=poppinkL, align=center]{管理 \emph{guǎnlǐ}\\ gérer}
        child { node {计划 \emph{jìhuà}} }
        child { node {习惯 \emph{xíguàn}} }
    };
\end{tikzpicture}
\legende{Arbre lexical autour de \zh{零花钱} : quatre piliers (Obtenir, Dépenser, Économiser, Gérer).}
\end{popfigure}

\courssub{Focus sur les phrases clés et structures utiles}

\begin{exemplebox}{Phrases modèles à mémoriser (Patterns productifs)}
\begin{enumerate}
    \item \textbf{Origine de l'argent : \zh{从 ... 那儿 得到 ...} (Cóng... nàr dédào...)}
    \zh{我每个星期从父母那儿得到五千法郎的零花钱。} \\
    \emph{Wǒ měi ge xīngqī cóng fùmǔ nàr dédào wǔqiān fǎláng de líng huā qián.} \\
    $\rightarrow$ Je reçois 5 000 FCFA d'argent de poche de mes parents chaque semaine.
    \item \textbf{Utilisation d'une ressource : \zh{用 ... 做 ... / 用 ... 去 ...} (Yòng... zuò... / qù...)}
    \zh{我用一部分钱坐车去学校。} \\
    \emph{Wǒ yòng yī bùfen qián zuò chē qù xuéxiào.} \\
    $\rightarrow$ J'utilise une partie de l'argent pour prendre le bus pour l'école.
    \item \textbf{Discipline financière : \zh{不 乱 动词} (Bù luàn V) — Ne pas V n'importe comment}
    \zh{我不乱花钱。} \\
    \emph{Wǒ bù luàn huā qián.} \\
    $\rightarrow$ Je ne dépense pas l'argent n'importe comment / Je ne gaspille pas.
    \item \textbf{Objectif d'épargne : \zh{想 存钱 买 ...} (Xiǎng cún qián mǎi...)}
    \zh{我想存钱买一本汉语词典。} \\
    \emph{Wǒ xiǎng cún qián mǎi yì běn Hànyǔ cídiǎn.} \\
    $\rightarrow$ Je veux économiser pour acheter un dictionnaire de chinois.
    \item \textbf{Jugement de valeur : \zh{觉得 ... 很 形容词} (Juéde... hěn Adj)}
    \zh{我觉得管理零花钱很重要。} \\
    \emph{Wǒ juéde guǎnlǐ líng huā qián hěn zhòngyào.} \\
    $\rightarrow$ Je trouve que gérer l'argent de poche est très important.
\end{enumerate}
\end{exemplebox}

\begin{attention}{Pièges fréquents des francophones}
\begin{itemize}
    \item \textbf{Ordre des mots (Cause/Conséquence)} : En français « Je veux acheter un dictionnaire, \emph{donc} j'économise ». En chinois : \zh{因为想买词典，\textbf{所以}存钱}. Ne jamais oublier \zh{所以} à l'écrit.
    \item \textbf{\zh{那儿} (nàr) vs \zh{那里} (nàlǐ)} : \zh{那儿} (avec \emph{érhuà}) est plus oral, courant au Nord (Pékin) ; \zh{那里} est plus standard à l'écrit. Le texte modèle utilise \zh{那儿} car c'est un texte « parlé-écrit » naturel. Acceptez les deux.
    \item \textbf{Classificateur \zh{本} (běn)} : Obligatoire pour les livres (\zh{词典}、\zh{书}、\zh{笔记本}). Erreur : *\zh{买一汉语词典}.
    \item \textbf{\zh{乱} (luàn) position} : \zh{不乱花钱} (ne pas dépenser n'importe comment) $\neq$ *\zh{不花乱钱} (incorrect). \zh{乱} est un adverbe de manière placé \textbf{avant} le verbe.
    \item \textbf{Monnaie} : \zh{法郎} (fǎláng) pour FCFA. En Chine : \zh{元} (yuán) à l'écrit ou \zh{块} (kuài) à l'oral. Ne pas confondre.
    \item \textbf{\zh{一部分} (yī bùfen)} : \zh{分} se prononce ici au ton neutre (\emph{fen}), pas \emph{fēn}.
\end{itemize}
\end{attention}

\courssub{Méthode : Analyser un texte modèle pour s'en inspirer}

\begin{methode}{Comment analyser un texte modèle (3 étapes)}
\begin{enumerate}
    \item \textbf{Souligner les connecteurs} (au fluo) : \zh{首先}、\zh{然后}、\zh{但是}、\zh{因为...所以...}、\zh{觉得}. Ils sont le squelette du texte.
    \item \textbf{Encadrer les verbes d'action principaux} (au stylo) : \zh{得到}、\zh{用}、\zh{坐}、\zh{买}、\zh{乱花}、\zh{想}、\zh{存}、\zh{觉得}. Ils portent le sens.
    \item \textbf{Entourer le vocabulaire thématique « Budget »} (crayon couleur) : \zh{零花钱}、\zh{法郎}、\zh{一部分钱}、\zh{存钱}、\zh{管理}. Ce sont les mots à réutiliser dans votre production.
\end{enumerate}
\emph{Application} : Reprenez le texte modèle, appliquez les 3 codes couleurs, puis réécrivez le texte en changeant : le montant (ex : 3 000 FCFA), le moyen de transport (moto-taxi \zh{坐摩托车}), l'objectif d'épargne (téléphone \zh{手机}).
\end{methode}

\begin{savaistu}{Le \zh{红包} (hóngbāo) : l'argent de poche rituel en Chine}
En Chine, le moment le plus important pour l'argent de poche des jeunes est le \textbf{Nouvel An chinois} (春节 \emph{Chūnjié}). Les aînés offrent aux enfants et aux lycéens des \textbf{enveloppes rouges} (\zh{红包} \emph{hóngbāo}) contenant de l'argent neuf. C'est une coutume ancienne (\zh{压岁钱} \emph{yā suì qián}, « argent pour écarter le mauvais sort de l'année »). Aujourd'hui, les \zh{红包} numériques (via WeChat \zh{微信红包}) sont devenus très courants. Les parents chinois profitent souvent de cet apport financier annuel pour apprendre à leurs enfants à \textbf{tenir un carnet de comptes} (\zh{记账} \emph{jì zhàng}) : noter les revenus, classer les dépenses (études, loisirs, snacks), et fixer un objectif d'épargne. C'est une éducation financière précoce, très valorisée socialement.
\end{savaistu}

\coursec{Écriture correcte des caractères du budget}

Les caractères du thème « argent et budget » comptent parmi les plus complexes du niveau Terminale. Ils sont fréquents à l'examen (dictée, thème, expression écrite) : une erreur de trait peut coûter des points. Analysons-les par \cle{radical} (部首 \emph{bùshǒu}) et par ordre des traits.

\begin{center}
\begin{tabularx}{\linewidth}{|X|X|X|X|}
\hline
\rowcolor{popblueL}
\textbf{Caractère} & \textbf{Pinyin} & \textbf{Radical (clé)} & \textbf{Traduction} \\
\hline
钱 & qián & 钅 (métal, à gauche) & argent \\
\hline
零 & líng & 雨 (pluie, en haut) & zéro ; (dans 零花钱) \\
\hline
花 & huā & 艹 (herbe, en haut) & fleur ; dépenser \\
\hline
存 & cún & 子 (enfant) & économiser, déposer \\
\hline
费 & fèi & 贝 (coquillage = argent) & frais, dépense \\
\hline
管 & guǎn & 竹 (bambou, en haut) & gérer, administrer \\
\hline
理 & lǐ & 王 (jade, à gauche) & raison, ordre ; (dans 管理) \\
\hline
贵 & guì & 贝 (coquillage = argent) & cher, coûteux \\
\hline
便 & biàn & 亻(homme, à gauche) & (dans 便宜 piányi : bon marché) \\
\hline
省 & shěng & 目 (œil) & économiser ; province \\
\hline
\end{tabularx}
\end{center}

\begin{propriete}{Le radical 贝 (bèi) : la clé de l'argent}
Le \cle{coquillage} \zh{贝} servait de monnaie dans la Chine antique. C'est pourquoi on le retrouve dans presque tous les caractères liés à l'argent : \zh{费} (frais), \zh{贵} (cher), \zh{财} (richesse), \zh{购} (acheter), \zh{账} (compte), \zh{赚} (gagner), \zh{赔} (rembourser/perdre). Repérer \zh{贝} dans un caractère inconnu vous donne un indice de sens : « cela concerne probablement l'argent ».
\end{propriete}

\textbf{Ordre des traits et pièges graphiques :}
\begin{itemize}
    \item \zh{钱} (10 traits) : d'abord la clé \zh{钅} à gauche (5 traits :撇、横、横、横、竖提), puis \zh{戋} à droite (横、横、斜钩、撇、点). \textbf{Piège} : ne pas oublier le point final en haut à droite ; ne pas écrire \zh{线} (fil) à la place.
    \item \zh{存} (6 traits) : 横、撇、竖、横钩、竖钩、横. La partie droite contient \zh{子} (enfant). \textbf{Piège} : ne pas confondre avec \zh{在} (être à) ni \zh{荐}.
    \item \zh{费} (9 traits) : d'abord \zh{弗} en haut, puis \zh{贝} en bas (竖、横折、撇、点). Règle générale : \textbf{haut $\rightarrow$ bas}.
    \item \zh{管} (14 traits) : clé \zh{竹} en haut (deux fois ⺮ simplifié), puis \zh{官} en bas. \textbf{Piège} : le bas est \zh{官} (fonctionnaire), pas \zh{宫} (palais).
    \item \zh{零} (13 traits) : clé \zh{雨} en haut, puis \zh{令} en bas. \textbf{Piège} : \zh{令} (4 traits) n'a qu'un seul point en bas, contrairement à \zh{今} + point.
    \item \zh{花} (7 traits) : clé \zh{艹} en haut, puis \zh{化} en bas (亻+ 匕). \textbf{Piège} : le dernier trait de \zh{化} est un 竖弯钩, pas un 撇.
\end{itemize}

\begin{attention}{Caractères proches à ne pas confondre}
\begin{itemize}
    \item \zh{钱} (qián, argent) $\neq$ \zh{线} (xiàn, fil/ligne) $\neq$ \zh{浅} (qiǎn, peu profond) — même phonétique \zh{戋}, clés différentes (钅/ 纟/ 氵).
    \item \zh{存} (cún, épargner) $\neq$ \zh{在} (zài, être à) — un trait de plus et une structure différente.
    \item \zh{费} (fèi, frais) $\neq$ \zh{废} (fèi, inutile/gaspiller) — homophones ! \zh{浪费} (làngfèi) = gaspiller.
    \item \zh{贵} (guì, cher) $\neq$ \zh{遗} (yí, perdre/laisser) — tous deux contiennent une forme proche en haut.
    \item \zh{花} (huā, dépenser/fleur) : attention au \textbf{faux ami interne} — le même caractère signifie « fleur » et « dépenser ». Le contexte décide : \zh{花钱} (dépenser de l'argent) vs \zh{一朵花} (une fleur).
\end{itemize}
\end{attention}

\coursec{Lexique actif et collocations du thème}

\begin{center}
\begin{tabularx}{\linewidth}{|X|X|X|X|}
\hline
\rowcolor{popblueL}
\textbf{Caractères} & \textbf{Pinyin} & \textbf{Nature} & \textbf{Traduction} \\
\hline
零花钱 & líng huā qián & n. & argent de poche \\
\hline
得到 & dédào & v. & obtenir, recevoir \\
\hline
一部分 & yī bùfen & n. & une partie \\
\hline
花（钱） & huā (qián) & v. & dépenser (de l'argent) \\
\hline
乱花 & luàn huā & v. & dépenser n'importe comment \\
\hline
存（钱） & cún (qián) & v. & épargner, mettre de côté \\
\hline
省钱 & shěng qián & loc. v. & économiser de l'argent \\
\hline
浪费 & làngfèi & v./adj. & gaspiller ; gaspillage \\
\hline
管理 & guǎnlǐ & v. & gérer, administrer \\
\hline
计划 & jìhuà & n./v. & plan ; planifier \\
\hline
预算 & yùsuàn & n. & budget (prévisionnel) \\
\hline
手机费 & shǒujī fèi & n. & frais de téléphone (crédit) \\
\hline
交通费 & jiāotōng fèi & n. & frais de transport \\
\hline
学习用品 & xuéxí yòngpǐn & n. & fournitures scolaires \\
\hline
课外书 & kèwài shū & n. & livres extrascolaires \\
\hline
旅游 & lǚyóu & v./n. & voyager ; tourisme \\
\hline
决定 & juédìng & v. & décider \\
\hline
有意义 & yǒu yìyì & loc. adj. & avoir du sens, être significatif \\
\hline
习惯 & xíguàn & n./v. & habitude ; avoir l'habitude de \\
\hline
记账 & jì zhàng & loc. v. & tenir ses comptes \\
\hline
\end{tabularx}
\end{center}

\begin{exemplebox}{Collocations utiles (verbe + complément)}
\begin{itemize}
    \item \zh{交手机费} (jiāo shǒujī fèi) — payer le crédit téléphonique.
    \item \zh{付交通费} (fù jiāotōng fèi) — payer les frais de transport.
    \item \zh{定一个计划} (dìng yí ge jìhuà) — établir un plan.
    \item \zh{养成好习惯} (yǎngchéng hǎo xíguàn) — prendre une bonne habitude.
    \item \zh{省钱买词典} (shěng qián mǎi cídiǎn) — économiser pour acheter un dictionnaire.
\end{itemize}
\end{exemplebox}

\coursec{Thème : français $\rightarrow$ chinois}

\begin{methode}{Méthode de traduction (thème) en 4 étapes}
\begin{enumerate}
    \item \textbf{Identifier le squelette} : Sujet + (complément de temps/lieu) + Verbe + Complément. En chinois, le temps et le lieu se placent \textbf{avant} le verbe.
    \item \textbf{Choisir les mots-outils} : \zh{的} pour la possession, classificateur adapté (\zh{本} pour livres, \zh{双} pour paires, \zh{个} par défaut), \zh{了} seulement pour une action accomplie ponctuelle.
    \item \textbf{Placer les connecteurs} : \zh{因为} en tête de cause, \zh{所以} en tête de conséquence ; virgule chinoise entre les deux.
    \item \textbf{Relire} : ordre des mots, tons du pinyin, ponctuation pleine chasse \zh{，。}.
\end{enumerate}
\end{methode}

\begin{exoresolu}{Thème guidé}
\textbf{Énoncé :} Traduisez en chinois : « Chaque mois, je reçois dix mille francs de mon père. Parce que je veux acheter un ballon de football, j'économise deux mille francs chaque semaine. »

\tcblower
\textbf{Solution détaillée}

Étape 1 — Squelette : \zh{我} + temps \zh{每个月} + source \zh{从爸爸那儿} + verbe \zh{得到} + montant.

Étape 2 — Montants : dix mille = \zh{一万} (attention : le chinois compte par \zh{万} « dix mille », pas par milliers au-delà de 9 000) ; deux mille = \zh{两千} (devant un classificateur ou une unité monétaire, on préfère \zh{两} à \zh{二}).

Étape 3 — Cause/conséquence : \zh{因为...所以...} ; classificateur du ballon : \zh{个}.

\textbf{Traduction :} \\
\zh{我每个月从爸爸那儿得到一万法郎的零花钱。因为我想买一个足球，所以我每个星期存两千法郎。} \\
\emph{Wǒ měi ge yuè cóng bàba nàr dédào yí wàn fǎláng de líng huā qián. Yīnwèi wǒ xiǎng mǎi yí ge zúqiú, suǒyǐ wǒ měi ge xīngqī cún liǎng qiān fǎláng.}
\end{exoresolu}

\begin{exercice}{Thème : à vous de traduire}
Traduisez en chinois (caractères + pinyin) :
\begin{enumerate}
    \item Je ne dépense pas mon argent de poche n'importe comment.
    \item D'abord je paie le transport, ensuite j'achète le petit-déjeuner.
    \item Bien que ces chaussures soient chères, je veux les acheter.
    \item Parce qu'il veut économiser de l'argent, il ne gaspille pas.
    \item Je trouve que tenir ses comptes est une bonne habitude.
\end{enumerate}
\end{exercice}

\begin{corrige}{Thème : à vous de traduire}
\begin{enumerate}
    \item \zh{我不乱花零花钱。} (Wǒ bù luàn huā líng huā qián.)
    \item \zh{首先我付交通费，然后我买早饭。} (Shǒuxiān wǒ fù jiāotōng fèi, ránhòu wǒ mǎi zǎofàn.)
    \item \zh{虽然这双鞋很贵，但是我想买。} (Suīrán zhè shuāng xié hěn guì, dànshì wǒ xiǎng mǎi.) — Classificateur \zh{双} pour une paire de chaussures.
    \item \zh{因为他想省钱，所以他不浪费。} (Yīnwèi tā xiǎng shěng qián, suǒyǐ tā bù làngfèi.)
    \item \zh{我觉得记账是个好习惯。} (Wǒ juéde jì zhàng shì ge hǎo xíguàn.)
\end{enumerate}
\end{corrige}

\coursec{Version : chinois $\rightarrow$ français et grille d'évaluation}

\begin{exoresolu}{Version guidée}
\textbf{Énoncé :} Traduisez en français le texte suivant.

\zh{我弟弟每个月有两千法郎的零花钱。虽然他很喜欢吃零食，但是他决定不乱花钱。因为他想买一本汉语词典，所以他每个星期存五百法郎。我觉得他做得对。}

\tcblower
\textbf{Solution détaillée}

Repérage : \zh{弟弟} = petit frère ; \zh{零食} = snacks, friandises ; \zh{做得对} = bien faire (structure verbe + \zh{得} + adjectif).

\textbf{Traduction :} « Mon petit frère reçoit deux mille francs d'argent de poche chaque mois. Bien qu'il aime beaucoup manger des friandises, il a décidé de ne pas dépenser n'importe comment. Parce qu'il veut acheter un dictionnaire de chinois, il économise cinq cents francs chaque semaine. Je trouve qu'il fait bien. »
\end{exoresolu}

\begin{exercice}{Version : à vous de traduire}
Traduisez en français :

\zh{我的同学玛丽每个星期从妈妈那儿得到三千法郎。首先，她交手机费。然后，她买学习用品。虽然她想买新衣服，但是她决定省钱。因为她想去北京旅游，所以她每个月存五千法郎。最后，她说：“管理零花钱是很有用的习惯。”}
\end{exercice}

\begin{corrige}{Version : à vous de traduire}
« Ma camarade de classe Marie reçoit trois mille francs de sa mère chaque semaine. D'abord, elle paie son crédit téléphonique. Ensuite, elle achète des fournitures scolaires. Bien qu'elle veuille acheter de nouveaux vêtements, elle décide d'économiser. Parce qu'elle veut aller en voyage à Pékin, elle met de côté cinq mille francs chaque mois. Enfin, elle dit : “Gérer son argent de poche est une habitude très utile.” »

\textbf{Points de vigilance :} \zh{同学} = camarade de classe ; \zh{新衣服} = nouveaux vêtements (\zh{件} serait le classificateur, omis ici car générique) ; \zh{北京} = Pékin ; la citation directe s'ouvre avec \zh{说：“...”}.
\end{corrige}

\begin{aretenir}{Grille d'évaluation de l'expression écrite (type examen)}
Votre production écrite sera notée selon quatre critères :
\begin{enumerate}
    \item \textbf{Respect de la tâche} (25 \%) : le sujet est traité, toutes les contraintes (connecteurs, montants, conclusion) sont présentes.
    \item \textbf{Correction grammaticale} (25 \%) : ordre des mots, \zh{因为...所以...} complet, classificateurs corrects, place de \zh{乱} avant le verbe.
    \item \textbf{Exactitude des caractères} (25 \%) : traits complets, pas de confusion entre caractères proches (\zh{钱/线}, \zh{存/在}, \zh{费/废}).
    \item \textbf{Richesse du lexique et cohérence} (25 \%) : vocabulaire du budget varié, connecteurs variés (\zh{首先、然后、虽然...但是...、最后}), ponctuation chinoise correcte.
\end{enumerate}
\end{aretenir}

\courssub{Exercice résolu : Analyse d'un paragraphe court}

\begin{exoresolu}{Analyse structurelle et lexicale}
\textbf{Énoncé :} Lisez le paragraphe suivant écrit par un élève chinois de Terminale. Appliquez la méthode en 3 étapes (connecteurs, verbes, vocabulaire budget) et répondez aux questions.

\zh{我每个月有一千元的零花钱。首先，我交手机费二百元。然后，我买课外书。虽然我想买新游戏，但是我决定存钱。因为我想去北京旅游，所以我每个月存五百元。最后，我觉得这样做很有意义。}

\begin{enumerate}
    \item Identifiez les 5 connecteurs et leur fonction.
    \item Quel est le montant mensuel ? Combien est épargné ?
    \item Traduisez la phrase avec \zh{虽然...但是...}.
\end{enumerate}

\tcblower
\textbf{Solution détaillée}

\textbf{1. Connecteurs (soulignés) :}
\begin{itemize}
    \item \zh{首先} (Shǒuxiān) — Énumération 1 (Dépense fixe).
    \item \zh{然后} (Ránhòu) — Succession (Dépense variable).
    \item \zh{虽然...但是...} (Suīrán... dànshì...) — \textbf{Concession / Opposition} : « Bien que je veuille..., je décide de... ». Structure plus forte que \zh{但是} seul.
    \item \zh{因为...所以...} (Yīnwèi... suǒyǐ...) — Cause $\rightarrow$ Conséquence (Objectif voyage $\rightarrow$ Épargne).
    \item \zh{最后} (Zuìhòu) — Conclusion temporelle / logique : « Enfin / En définitive ».
\end{itemize}

\textbf{2. Verbes encadrés :} \zh{有} (avoir), \zh{交} (payer), \zh{买} (acheter), \zh{想} (vouloir), \zh{决定} (décider), \zh{存} (épargner), \zh{去} (aller), \zh{觉得} (trouver).

\textbf{3. Vocabulaire budget (entouré) :} \zh{零花钱}、\zh{元} (yuán)、\zh{手机费}、\zh{存钱}、\zh{旅游} (lǚyóu).

\textbf{Réponses aux questions :}
\begin{itemize}
    \item Revenu : 1 000 \zh{元}/mois. Épargne : 500 \zh{元}/mois (taux d'épargne de 50 \% !).
    \item \zh{虽然我想买新游戏，但是我决定存钱。} (Suīrán wǒ xiǎng mǎi xīn yóuxì, dànshì wǒ juédìng cún qián.) — « Bien que je veuille acheter un nouveau jeu (vidéo), je décide d'économiser. »
\end{itemize}
\end{exoresolu}

\courssub{Entraînement : À votre tour !}

\begin{exercice}{Production guidée : Mon budget de la semaine}
Rédigez un paragraphe de 6 à 8 phrases en chinois (caractères + pinyin) pour décrire \textbf{votre} gestion de l'argent de poche cette semaine. Inspirez-vous du texte modèle et de l'exercice résolu.

\textbf{Contraintes obligatoires :}
\begin{enumerate}
    \item Indiquer la source et le montant (FCFA ou Yuan).
    \item Utiliser \zh{首先} et \zh{然后} pour deux dépenses réelles (transport, manger, crédit, photocopies...).
    \item Utiliser \zh{但是} ou \zh{虽然...但是...} pour montrer un arbitrage (renoncer à un snack, un jeu, une sortie).
    \item Utiliser \zh{因为...所以...} pour justifier votre épargne (objectif : chaussures, livre, cadeau, voyage...).
    \item Conclure avec \zh{我觉得...很重要/有意义/有用}.
    \item Écrire les caractères soigneusement (attention aux radicaux vus en section 2).
\end{enumerate}

\textbf{Aide au démarrage (à compléter) :} \\
\zh{我每个星期从 \_\_\_\_\_\_ 那儿得到 \_\_\_\_\_\_ 法郎的零花钱。}\\
\zh{首先，我用一部分钱 \_\_\_\_\_\_ 。}\\
\zh{然后，我 \_\_\_\_\_\_ 。}\\
\zh{虽然我想 \_\_\_\_\_\_ ，但是我 \_\_\_\_\_\_ 。}\\
\zh{因为我想 \_\_\_\_\_\_ ，所以我每个星期存 \_\_\_\_\_\_ 。}\\
\zh{我觉得管理零花钱 \_\_\_\_\_\_ 。}
\end{exercice}

\begin{corrige}{Exercice : Mon budget de la semaine}
\textbf{Exemple de production attendue (niveau HSK 3/4 Terminale) :}

\zh{我每个星期从妈妈那儿得到三千法郎的零花钱。} (Wǒ měi ge xīngqī cóng māma nàr dédào sānqiān fǎláng de líng huā qián.)\\
\zh{首先，我用一千法郎坐摩托车去学校。} (Shǒuxiān, wǒ yòng yìqiān fǎláng zuò mótuōchē qù xuéxiào.)\\
\zh{然后，我买早饭，花五百法郎。} (Ránhòu, wǒ mǎi zǎofàn, huā wǔbǎi fǎláng.)\\
\zh{虽然我想买冰淇淋，但是我决定不买。} (Suīrán wǒ xiǎng mǎi bīngqílín, dànshì wǒ juédìng bù mǎi.)\\
\zh{因为我想买一双新篮球鞋，所以我每个星期存一千五百法郎。} (Yīnwèi wǒ xiǎng mǎi yì shuāng xīn lánqiúxié, suǒyǐ wǒ měi ge xīngqī cún yìqiān wǔbǎi fǎláng.)\\
\zh{我觉得管理零花钱很有意义。} (Wǒ juéde guǎnlǐ líng huā qián hěn yǒu yìyì.)

\textbf{Points de vigilance pour l'auto-correction :}
\begin{itemize}
    \item \zh{坐摩托车} (zuò mótuōchē) — Contexte camerounais adapté (moto-taxi).
    \item \zh{花} (huā) utilisé comme verbe transitif : \zh{花五百法郎}.
    \item \zh{一双} (yì shuāng) — Classificateur pour les chaussures/chaussettes (paires).
    \item \zh{很有意义} (hěn yǒu yìyì) — Structure adjectivale « avoir du sens / être significatif », plus riche que \zh{很重要}.
    \item Ponctuation chinoise \zh{，} \zh{。}, pas d'espaces entre caractères.
\end{itemize}
\end{corrige}

\coursec{Structure et connecteurs du texte}

Dans la section précédente, nous avons lu et commenté un texte modèle sur la gestion de l'argent de poche. Nous avons observé \emph{ce que dit} l'auteur ; il s'agit maintenant de comprendre \emph{comment il le dit} : dans quel ordre les idées s'enchaînent, quels mots relient les phrases entre elles, et comment une opinion vient conclure le texte. C'est exactement cette architecture que vous devrez reproduire dans vos propres productions écrites. Un bon texte chinois, même court, n'est jamais une simple juxtaposition de phrases : c'est un parcours guidé par des connecteurs.

\courssub{Le plan type en quatre parties}

Un texte simple sur la gestion de l'argent de poche suit presque toujours le même schéma en quatre étapes. Retenez ce plan : il vous servira de canevas pour toutes vos rédactions de ce module.

\begin{definition}[Le plan type en quatre parties]
\begin{enumerate}
\item \textbf{Introduction} : présenter la source de l'argent (qui me donne de l'argent, combien, quand).
\item \textbf{Développement} : décrire les dépenses (qu'est-ce que j'achète, où, pourquoi).
\item \textbf{Économies et buts} : expliquer ce que l'on économise et pour quel projet.
\item \textbf{Conclusion} : donner son opinion sur la gestion de l'argent.
\end{enumerate}
\end{definition}

\begin{popfigure}
\begin{center}
\begin{tikzpicture}[node distance=0.9cm, >=Stealth]
\node[draw=popblue, fill=popblueL, rounded corners, minimum width=2.6cm, minimum height=1.1cm, align=center] (a) {\textbf{引言} yǐnyán\\ Introduction : argent reçu};
\node[draw=poporange, fill=poporangeL, rounded corners, minimum width=2.6cm, minimum height=1.1cm, align=center, right=of a] (b) {\textbf{花钱} huā qián\\ Dépenses};
\node[draw=popgreen, fill=popgreenL, rounded corners, minimum width=2.6cm, minimum height=1.1cm, align=center, right=of b] (c) {\textbf{存钱} cún qián\\ Économies et buts};
\node[draw=poppurple, fill=poppurpleL, rounded corners, minimum width=2.6cm, minimum height=1.1cm, align=center, right=of c] (d) {\textbf{看法} kànfǎ\\ Opinion};
\draw[->, thick, popblue] (a) -- node[above, font=\small]{首先} (b);
\draw[->, thick, poporange] (b) -- node[above, font=\small]{然后} (c);
\draw[->, thick, popgreen] (c) -- node[above, font=\small]{最后} (d);
\end{tikzpicture}
\end{center}
\legende{Organigramme du plan type : quatre parties reliées par les connecteurs de séquence.}
\end{popfigure}

\courssub{Les connecteurs de séquence : 首先, 然后, 最后}

Avant d'énoncer la règle, observons deux phrases tirées de la vie quotidienne d'un élève de Yaoundé :

\begin{exemplebox}[Observation]
\zh{首先，我坐车去学校。然后，我买早饭。}\\
\emph{Shǒuxiān, wǒ zuò chē qù xuéxiào. Ránhòu, wǒ mǎi zǎofàn.}\\
« D'abord, je prends le transport pour aller à l'école. Ensuite, j'achète le petit-déjeuner. »
\end{exemplebox}

\begin{propriete}[Connecteurs de séquence]
\begin{itemize}
\item \zh{首先} \emph{shǒuxiān} : « d'abord, premièrement ». Il ouvre la séquence.
\item \zh{然后} \emph{ránhòu} : « ensuite, puis ». Il peut se répéter plusieurs fois dans le même texte.
\item \zh{最后} \emph{zuìhòu} : « enfin, pour finir ». Il ferme la séquence.
\end{itemize}
Structure : \textbf{connecteur + ，+ phrase complète}. Le connecteur se place en tête de phrase, suivi d'une virgule chinoise \zh{，}.
\end{propriete}

Exemple supplémentaire :

\zh{首先，妈妈给我五千法郎。然后，我买书和笔。最后，我把剩下的钱存起来。}\\
\emph{Shǒuxiān, māma gěi wǒ wǔ qiān fǎláng. Ránhòu, wǒ mǎi shū hé bǐ. Zuìhòu, wǒ bǎ shèngxià de qián cún qǐlai.}\\
« D'abord, maman me donne cinq mille francs. Ensuite, j'achète des livres et des stylos. Enfin, je mets de côté l'argent qui reste. »

Notez au passage la structure \zh{把 + complément d'objet + verbe + 起来} : \zh{把剩下的钱存起来} « mettre de côté l'argent restant ». La construction en \zh{把} place l'objet \emph{avant} le verbe quand l'action produit un résultat sur cet objet ; elle est très fréquente à l'écrit.

\begin{attention}[Ne pas confondre 然后 et 以后]
\zh{然后} \emph{ránhòu} signifie « ensuite » (enchaînement d'actions) et se place \emph{en tête} de la proposition ; \zh{以后} \emph{yǐhòu} signifie « après » et se place \emph{après} le verbe ou le groupe verbal : \zh{下课以后} « après la fin des cours ». Dire \zh{我吃饭然后} est une erreur : on dit \zh{我吃饭以后} « après avoir mangé » ou \zh{然后，我……} « ensuite, je... ».
\end{attention}

\courssub{Les connecteurs logiques : 因为...所以..., 但是, 如果...就...}

Passons maintenant aux connecteurs qui expriment non plus l'ordre des actions, mais les \emph{relations logiques} entre les idées : la cause, l'opposition, la condition. Ce sont eux qui donnent à votre texte sa valeur argumentative — et ils sont très attendus à l'examen.

\textbf{La cause et la conséquence : 因为...所以...}\par

\begin{propriete}[因为...所以... : une différence majeure avec le français]
En chinois, \zh{因为} \emph{yīnwèi} « parce que » et \zh{所以} \emph{suǒyǐ} « donc » peuvent — et même aiment — se combiner dans la \textbf{même phrase} :\\
\textbf{因为 + cause，所以 + conséquence。}\\
En français, « parce que » et « donc » ne se cumulent jamais : on dit « parce que je veux économiser, je dépense avec prudence » ou « je veux économiser, donc je dépense avec prudence », mais jamais les deux ensemble. Le chinois, lui, accepte et préfère souvent la structure complète.
\end{propriete}

\begin{exemplebox}[Exemple justificatif]
\zh{因为我想存钱，所以我很会省钱。}\\
\emph{Yīnwèi wǒ xiǎng cún qián, suǒyǐ wǒ hěn huì shěng qián.}\\
« Parce que je veux mettre de l'argent de côté, je sais dépenser avec prudence. »
\end{exemplebox}

Autres exemples :

\zh{因为车费很贵，所以我走路去学校。} \emph{Yīnwèi chēfèi hěn guì, suǒyǐ wǒ zǒulù qù xuéxiào.} « Comme le transport coûte cher, je vais à l'école à pied. »

\zh{因为我想买一台电脑，所以我每个月存两千法郎。} \emph{Yīnwèi wǒ xiǎng mǎi yì tái diànnǎo, suǒyǐ wǒ měi ge yuè cún liǎng qiān fǎláng.} « Parce que je veux acheter un ordinateur, j'économise deux mille francs chaque mois. »

Remarquez le classificateur \zh{台} \emph{tái} pour les machines et appareils (ordinateur, télévision), et l'emploi de \zh{两} \emph{liǎng} « deux » devant un classificateur ou une unité de mesure, jamais \zh{二}.

\begin{attention}[Erreur fréquente des francophones]
Traduire « donc » seul par \zh{所以} en début de phrase, sans exprimer la cause : \zh{所以，我不买零食。} est une phrase orpheline qui surprend le lecteur chinois. En chinois, la cause est presque toujours exprimée explicitement : dites plutôt \zh{因为我想省钱，所以我不买零食。} « Parce que je veux économiser, je n'achète pas de friandises. »
\end{attention}

\textbf{L'opposition : 但是}\par

\zh{但是} \emph{dànshì} « mais » introduit une concession. Il se place en tête de la seconde proposition, après la virgule : \textbf{phrase 1，但是 + phrase 2。}

\zh{我的钱不多，但是我很会安排。} \emph{Wǒ de qián bù duō, dànshì wǒ hěn huì ānpái.} « Je n'ai pas beaucoup d'argent, mais je sais bien l'organiser. »

\zh{零食很好吃，但是太贵了。} \emph{Língshí hěn hǎochī, dànshì tài guì le.} « Les friandises sont bonnes, mais elles coûtent trop cher. »

On peut renforcer l'opposition avec \zh{虽然...但是...} « bien que... (mais)... », qui, comme \zh{因为...所以...}, s'emploient ensemble en chinois : \zh{虽然我的钱不多，但是我很开心。} \emph{Suīrán wǒ de qián bù duō, dànshì wǒ hěn kāixīn.} « Bien que je n'aie pas beaucoup d'argent, je suis content. » Là encore, le français refuse le cumul (« bien que... mais... » est incorrect en français), le chinois l'exige presque.

\textbf{La condition : 如果...就...}\par

Structure : \textbf{如果 + condition，就 + résultat。} « si... alors... ». Le mot \zh{就} \emph{jiù} se place juste avant le verbe de la seconde proposition, après le sujet si celui-ci est exprimé.

\zh{如果我有足够的钱，我就买一双新鞋。} \emph{Rúguǒ wǒ yǒu zúgòu de qián, wǒ jiù mǎi yì shuāng xīn xié.} « Si j'ai assez d'argent, j'achète une nouvelle paire de chaussures. »

\zh{如果这个月我省钱，下个月就可以买书。} \emph{Rúguǒ zhè ge yuè wǒ shěng qián, xià ge yuè jiù kěyǐ mǎi shū.} « Si j'économise ce mois-ci, le mois prochain je pourrai acheter des livres. »

Notez le classificateur \zh{双} \emph{shuāng} pour les paires (chaussures, chaussettes, baguettes) : \zh{一双新鞋} « une paire de chaussures neuves ».

\courssub{Exprimer la fréquence et les habitudes}

Pour décrire ses dépenses régulières, on utilise les marqueurs de fréquence. Attention à leur place : les adverbes de fréquence se placent \textbf{avant le verbe}.

\begin{center}
\begin{tabularx}{\linewidth}{|X|X|X|X|}
\hline
\rowcolor{popblueL} Caractères & Pinyin & Nature & Traduction \\
\hline
每个星期 & měi ge xīngqī & complément de temps & chaque semaine \\
\hline
常常 & chángcháng & adverbe & souvent \\
\hline
有时候 & yǒu shíhou & adverbe & parfois \\
\hline
从不 & cóng bù & adverbe & ne... jamais \\
\hline
\end{tabularx}
\end{center}

\zh{我每个星期买一次面包。} \emph{Wǒ měi ge xīngqī mǎi yí cì miànbāo.} « J'achète du pain une fois par semaine. »

\zh{我常常在学校门口买早饭。} \emph{Wǒ chángcháng zài xuéxiào ménkǒu mǎi zǎofàn.} « J'achète souvent le petit-déjeuner devant l'école. »

\zh{有时候我买零食，有时候不买。} \emph{Yǒu shíhou wǒ mǎi língshí, yǒu shíhou bù mǎi.} « Parfois j'achète des friandises, parfois non. »

\zh{我从不乱花钱。} \emph{Wǒ cóng bù luàn huā qián.} « Je ne dépense jamais mon argent n'importe comment. »

\begin{attention}[Place des adverbes de fréquence]
En français, « souvent » peut se placer en fin de phrase (« j'achète du pain souvent »). En chinois, c'est impossible : \zh{常常}, \zh{有时候} et \zh{从不} se placent \emph{avant le verbe}. Ne dites jamais \zh{我买面包常常}. Rappel : le complément de lieu avec \zh{在} se place lui aussi avant le verbe : \zh{我在学校门口买早饭} (et non après, comme en français).
\end{attention}

\courssub{Conclure par une opinion : 我觉得, 我认为, ...很重要}

La dernière partie du plan donne votre avis. Trois outils simples :

\begin{itemize}
\item \zh{我觉得 + phrase} : « je trouve que... », opinion personnelle et courante.
\item \zh{我认为 + phrase} : « je pense que... », un peu plus formel.
\item \zh{...很重要} : « ... est très important », pour évaluer une chose ou une habitude.
\end{itemize}

\zh{我觉得存钱是一个好习惯。} \emph{Wǒ juéde cún qián shì yí ge hǎo xíguàn.} « Je trouve qu'économiser est une bonne habitude. »

\zh{我认为学生应该学会管理自己的钱。} \emph{Wǒ rènwéi xuésheng yīnggāi xuéhuì guǎnlǐ zìjǐ de qián.} « Je pense que les élèves doivent apprendre à gérer leur propre argent. »

\zh{合理安排零花钱很重要。} \emph{Hélǐ ānpái línghuāqián hěn zhòngyào.} « Bien organiser son argent de poche est très important. »

Observez dans le dernier exemple qu'un groupe verbal entier (\zh{合理安排零花钱} « organiser raisonnablement son argent de poche ») peut occuper la place du sujet, sans aucun mot supplémentaire : le chinois n'a pas besoin d'un équivalent de « c'est... que ».

\begin{center}
\begin{tabularx}{\linewidth}{|X|X|X|}
\hline
\rowcolor{popblueL} Fonction & Structure chinoise & Équivalent français \\
\hline
Séquence & 首先...然后...最后... & d'abord... ensuite... enfin... \\
\hline
Cause-conséquence & 因为...所以... (cumul possible) & parce que... / donc... (jamais cumulés) \\
\hline
Opposition & ...，但是... & ..., mais... \\
\hline
Condition & 如果...就... & si... alors... \\
\hline
Opinion & 我觉得 / 我认为 / ...很重要 & je trouve que / je pense que / c'est important \\
\hline
\end{tabularx}
\end{center}

\begin{savaistu}[L'argent de poche en Chine et au Cameroun]
En Chine, l'argent de poche se dit \zh{零花钱} \emph{línghuāqián}, littéralement « argent pour les petites dépenses ». Lors de la fête du Printemps (Nouvel An lunaire), les enfants reçoivent aussi de l'argent dans des enveloppes rouges, les \zh{红包} \emph{hóngbāo}, que beaucoup de parents les encouragent à économiser. Au Cameroun, l'« argent de poche » des élèves sert souvent à payer le transport et le petit-déjeuner : c'est un excellent sujet de comparaison interculturelle à l'oral comme à l'écrit.
\end{savaistu}

\courssub{Texte d'application : un paragraphe bien construit}

\begin{textebox}[Mon argent de poche — texte original]
\zh{首先，我每个星期从爸爸那里得到三千法郎的零花钱。然后，我用这些钱买早饭和笔。有时候我也买零食，但是我知道零食太贵，所以我不常常买。因为我想买一个新书包，所以我每个月存一千法郎。如果我省钱，三个月以后就可以买书包了。最后，我觉得管理零花钱很重要：会安排钱的学生，将来也会安排自己的生活。}\\
\emph{Shǒuxiān, wǒ měi ge xīngqī cóng bàba nàli dédào sān qiān fǎláng de línghuāqián. Ránhòu, wǒ yòng zhèxiē qián mǎi zǎofàn hé bǐ. Yǒu shíhou wǒ yě mǎi língshí, dànshì wǒ zhīdào língshí tài guì, suǒyǐ wǒ bù chángcháng mǎi. Yīnwèi wǒ xiǎng mǎi yí ge xīn shūbāo, suǒyǐ wǒ měi ge yuè cún yì qiān fǎláng. Rúguǒ wǒ shěng qián, sān ge yuè yǐhòu jiù kěyǐ mǎi shūbāo le. Zuìhòu, wǒ juéde guǎnlǐ línghuāqián hěn zhòngyào : huì ānpái qián de xuésheng, jiānglái yě huì ānpái zìjǐ de shēnghuó.}\\
« D'abord, je reçois chaque semaine trois mille francs d'argent de poche de mon père. Ensuite, j'utilise cet argent pour acheter le petit-déjeuner et des stylos. Parfois j'achète aussi des friandises, mais je sais qu'elles coûtent cher, donc je n'en achète pas souvent. Parce que je veux acheter un nouveau cartable, j'économise mille francs chaque mois. Si j'économise, dans trois mois je pourrai acheter le cartable. Enfin, je trouve que gérer son argent de poche est très important : un élève qui sait organiser son argent saura aussi organiser sa vie plus tard. »
\end{textebox}

Observez la structure \zh{从 + personne + 那里 + verbe} : \zh{从爸爸那里得到} « recevoir de la part de papa ». Le mot \zh{那里} « là » s'ajoute après la personne pour marquer l'origine : \zh{从妈妈那里} « de la part de maman », \zh{从朋友那里} « de la part d'un ami ».

\textbf{Questions de compréhension :} 1) Combien d'argent l'élève reçoit-il chaque semaine ? 2) Pourquoi n'achète-t-il pas souvent de friandises ? 3) Quel est son projet d'économie ? 4) Quelle est son opinion finale ?

\courssub{Méthode et entraînement}

\begin{methode}[Construire un paragraphe en cinq phrases]
\begin{enumerate}
\item \textbf{Introduction} (1 phrase) : qui me donne de l'argent, combien, quand — commencez par \zh{首先}.
\item \textbf{Dépenses} (2 phrases) : ce que j'achète et où — reliez avec \zh{然后}, ajoutez une nuance avec \zh{但是} ou \zh{有时候}.
\item \textbf{Économie} (1 phrase) : mon projet — utilisez \zh{因为...所以...} ou \zh{如果...就...}.
\item \textbf{Opinion} (1 phrase) : concluez avec \zh{最后，我觉得...} ou \zh{...很重要}.
\end{enumerate}
Vérifiez ensuite : chaque connecteur est-il en tête de proposition ? Chaque adverbe de fréquence est-il avant le verbe ? Chaque classificateur est-il le bon (\zh{台} pour les appareils, \zh{双} pour les paires, \zh{个} par défaut) ?
\end{methode}

\begin{exoresolu}[Remettre le paragraphe dans l'ordre]
Énoncé — Remettez ces quatre phrases dans l'ordre logique du plan type et ajoutez les connecteurs manquants :\\
a) \zh{我觉得省钱很重要。} b) \zh{我买早饭和本子。} c) \zh{妈妈每个星期给我两千法郎。} d) \zh{我把剩下的钱存起来，因为我想买一个足球。}
\tcblower
Solution détaillée — L'ordre suit le plan : introduction (argent reçu) → dépenses → économies → opinion. On obtient :\\
\zh{首先，妈妈每个星期给我两千法郎。然后，我买早饭和本子。最后，我把剩下的钱存起来，因为我想买一个足球。我觉得省钱很重要。}\\
\emph{Shǒuxiān, māma měi ge xīngqī gěi wǒ liǎng qiān fǎláng. Ránhòu, wǒ mǎi zǎofàn hé běnzi. Zuìhòu, wǒ bǎ shèngxià de qián cún qǐlai, yīnwèi wǒ xiǎng mǎi yí ge zúqiú. Wǒ juéde shěng qián hěn zhòngyào.}\\
« D'abord, maman me donne deux mille francs chaque semaine. Ensuite, j'achète le petit-déjeuner et des cahiers. Enfin, je mets de côté l'argent restant, parce que je veux acheter un ballon de football. Je trouve qu'économiser est très important. »\\
Remarque : la phrase d'opinion peut s'ouvrir directement par \zh{我觉得}, ou être introduite par \zh{最后} si la phrase d'économie n'en a pas besoin. Notez aussi \zh{两千} avec \zh{两} devant la mesure \zh{千}, et non \zh{二千}.
\end{exoresolu}

\begin{exercice}[À vous de jouer]
1) Traduisez en chinois : « Parce que le transport coûte cher, je vais à l'école à pied. »\\
2) Complétez avec les connecteurs qui conviennent : \zh{......我想买电脑，......我每个月存钱。}\\
3) Traduisez en chinois : « Bien que je n'aie pas beaucoup d'argent, je suis content. »\\
4) Rédigez un paragraphe de cinq phrases sur votre argent de poche en suivant la méthode ci-dessus, avec au moins trois connecteurs différents.
\end{exercice}

\begin{corrige}[Exercice]
1) \zh{因为车费很贵，所以我走路去学校。} \emph{Yīnwèi chēfèi hěn guì, suǒyǐ wǒ zǒulù qù xuéxiào.}\\
2) \zh{因为我想买电脑，所以我每个月存钱。} — le cumul \zh{因为...所以...} est ici naturel et recommandé.\\
3) \zh{虽然我的钱不多，但是我很开心。} \emph{Suīrán wǒ de qián bù duō, dànshì wǒ hěn kāixīn.} — en chinois, \zh{虽然} et \zh{但是} se cumulent, contrairement au français.\\
4) Modèle possible : \zh{首先，我每个星期得到两千法郎。然后，我买早饭和书。有时候我买零食，但是不常常买。因为我想买一双鞋，所以我每个月存五百法郎。最后，我觉得管理零花钱很重要。} \emph{Shǒuxiān, wǒ měi ge xīngqī dédào liǎng qiān fǎláng. Ránhòu, wǒ mǎi zǎofàn hé shū. Yǒu shíhou wǒ mǎi língshí, dànshì bù chángcháng mǎi. Yīnwèi wǒ xiǎng mǎi yì shuāng xié, suǒyǐ wǒ měi ge yuè cún wǔ bǎi fǎláng. Zuìhòu, wǒ juéde guǎnlǐ línghuāqián hěn zhòngyào.}
\end{corrige}

\begin{aretenir}[L'essentiel de la section]
Un texte sur l'argent de poche suit quatre étapes : argent reçu → dépenses → économies → opinion. Les connecteurs de séquence \zh{首先}, \zh{然后}, \zh{最后} marquent ce parcours ; \zh{因为...所以...} (cumulés, contrairement au français), \zh{虽然...但是...} et \zh{如果...就...} articulent la logique ; \zh{我觉得} et \zh{...很重要} concluent. Les adverbes de fréquence et les compléments de lieu en \zh{在} se placent toujours avant le verbe, et les classificateurs (\zh{台}, \zh{双}, \zh{个}) accompagnent les noms comptés.
\end{aretenir}

\coursec{Écriture correcte des caractères du budget}

Après avoir analysé la structure d'un texte argumentatif sur la gestion de l'argent de poche et identifié les connecteurs logiques qui en assurent la cohérence, nous devons maintenant nous assurer que le vocabulaire spécialisé du budget s'écrit correctement. En chinois, l'écriture n'est pas une simple transcription phonétique : chaque caractère porte en lui une logique interne (clé sémantique + composant phonétique) et un ordre des traits immuable. Maîtriser l'écriture des caractères \cle{钱}, \cle{费}, \cle{算}, \cle{储蓄}, \cle{预算} et savoir distinguer les paires visuellement proches (\cle{买}/\cle{卖}, \cle{钱}/\cle{线}, \cle{借}/\cle{错}, \cle{费}/\cle{佛}) est la condition \emph{sine qua non} pour produire un texte lisible, noté au maximum à l'examen, et pour éviter les contresens embarrassants.

\courssub{Les caractères fondamentaux du thème budgétaire : analyse structurale et ordre des traits}

Observons d'abord chaque caractère comme un architecte observe une pierre de taille : par sa clé (radical), son nombre de traits, sa structure interne et son étymologie. Cette approche \emph{analytique} remplace l'apprentissage par cœur pur et dur, source d'oublis rapides.

\begin{definition}[La clé \zh{贝} (bèi) : coquillage-monnaie de la Chine antique]
Le caractère \zh{贝} (bèi) représente un coquillage \emph{cauri} (牛螺), utilisé comme monnaie d'échange dans la Chine des Shang et des Zhou (environ 1600–256 av. J.-C.). Par extension, \zh{贝} est devenu la \textbf{clé sémantique} (部首, bùshǒu) de nombreux caractères liés à l'argent, à la valeur, aux transactions : \zh{费} (fèi, dépense), \zh{贵} (guì, cher/précieux), \zh{贫} (pín, pauvre), \zh{购} (gòu, acheter), \zh{赚} (zhuàn, gagner de l'argent), \zh{账} (zhàng, compte), \zh{资} (zī, capital/fonds). Reconnaître \zh{贝} à gauche ou en bas d'un caractère vous donne immédiatement son champ sémantique : « argent, valeur, échange ».
\end{definition}

\begin{exemplebox}[Mots à clé \zh{贝} — lecture et traduction]
\begin{itemize}
\item \zh{贵} \emph{guì} — cher, précieux, honorifique (\zh{贵姓} \emph{guìxìng} « votre nom honorable »).
\item \zh{贫} \emph{pín} — pauvre, démuni (\zh{贫穷} \emph{pínqióng} « pauvreté »).
\item \zh{购} \emph{gòu} — acheter, achat (\zh{购物} \emph{gòuwù} « faire les courses / shopping »).
\item \zh{赚} \emph{zhuàn} — gagner (argent), profiter (\zh{赚钱} \emph{zhuànqián} « gagner de l'argent »).
\item \zh{账} \emph{zhàng} — compte, facture (\zh{账单} \zh{zhàngdān} « facture, relevé » ; \zh{记账} \zh{jìzhàng} « tenir les comptes »).
\item \zh{资} \emph{zī} — capital, fonds, ressources (\zh{资金} \emph{zījīn} « fonds, capital »).
\end{itemize}
\end{exemplebox}

\begin{center}
\begin{tabularx}{\linewidth}{|X|X|X|X|}
\hline
\rowcolor{popblueL} \textbf{Caractères} & \textbf{Pinyin} & \textbf{Nature} & \textbf{Traduction} \\
\hline
贵 & guì & adj. & cher, précieux, honorifique \\
贫 & pín & adj. & pauvre \\
购 & gòu & v. & acheter \\
赚 & zhuàn & v. & gagner (argent) \\
账 & zhàng & n. & compte, facture \\
资 & zī & n. & capital, fonds, ressources \\
\hline
\end{tabularx}
\end{center}

\medskip

\textbf{1. \zh{钱} (qián) — argent, monnaie}\par
\begin{itemize}
\item \textbf{Clé} : \zh{钅} (jīn, métal/or), forme compressée de \zh{金} à gauche.
\item \textbf{Structure} : clé \zh{钅} (5 traits) + composant phonétique \zh{戋} (jiān, 7 traits) = \textbf{12 traits au total}.
\item \textbf{Ordre des traits} (respectez-le scrupuleusement) :
\begin{enumerate}
\item \zh{钅} : trait horizontal court (一), trait vertical (丨), trait horizontal court (一), trait horizontal court (一), trait pointé vers le bas à droite (㇏) — \emph{la clé métal s'écrit toujours en premier, de haut en bas}.
\item \zh{戋} : trait horizontal long (一), trait vertical (丨), trait horizontal (一), trait horizontal (一), trait vertical (丨), trait horizontal (一), trait horizontal final (一).
\end{enumerate}
\item \textbf{Astuce mnémotechnique} : « L'argent (钱) est du métal (钅) qu'on compte en petites quantités (戋 suggère la petitesse, la minutie) ».
\end{itemize}

\textbf{2. \zh{费} (fèi) — dépenser, frais, coûter}\par
\begin{itemize}
\item \textbf{Clé} : \zh{贝} (bèi, coquillage/monnaie) en bas.
\item \textbf{Structure} : partie haute \zh{弗} (fú, phonétique, 5 traits) + \zh{贝} (4 traits) = \textbf{9 traits}.
\item \textbf{Ordre des traits} :
\begin{enumerate}
\item \zh{弗} : horizontal (一), vertical (丨), horizontal (一), horizontal (一), crochet vertical (㇆).
\item \zh{贝} : horizontal (一), vertical (丨), horizontal (一), point final (丶) — \emph{attention : le dernier trait de 贝 est un point, pas un trait horizontal !}
\end{enumerate}
\item \textbf{Collocations} : \zh{学费} (xuéfèi, frais de scolarité), \zh{生活费} (shēnghuófèi, argent de poche / frais de vie), \zh{手续费} (shǒuxùfèi, frais de dossier), \zh{费心} (fèixīn, se donner du mal).
\end{itemize}

\textbf{3. \zh{算} (suàn) — calculer, compter, estimer, planifier}\par
C'est un caractère \emph{à étages} (上下结构) en trois niveaux, clé \zh{⺮} (zhú, bambou) en haut.

\begin{popfigure}
\legende{Décomposition du caractère \zh{算} (suàn) en trois étages avec ordre des traits}
\begin{tikzpicture}[
    node distance=0.8cm and 1.2cm,
    box/.style={draw=popdark, thick, rounded corners=4pt, fill=popcream, minimum width=2.8cm, minimum height=1.2cm, align=center, font=\small},
    arrow/.style={->, >=Stealth, thick, popblue},
    labelbox/.style={draw=poporange, thick, rounded corners=2pt, fill=poporangeL, minimum width=2.2cm, minimum height=0.9cm, align=center, font=\small},
]
% Étage 1 : ⺮ (bambou)
\node[box, align=center] (zhu){\zh{⺮}\\\emph{zhú} (bambou)\\6 traits};
% Flèche
\draw[arrow] (zhu.south) -- ++(0,-0.4) node[labelbox, anchor=north, align=center] (mu){\zh{目}\\\emph{mù} (œil)\\5 traits};
% Flèche
\draw[arrow] (mu.south) -- ++(0,-0.4) node[labelbox, anchor=north, align=center] (gong){\zh{廾}\\\emph{gǒng} (mains jointes)\\3 traits};
% Annotations ordre
\node[left=0.3cm of zhu, font=\tiny, popmuted, align=right] {1. 横\\(横)\\
2. 竖\\(垂直)\\
3. 撇\\(pié)\\
4. 点\\(diǎn)\\
5. 横\\(横)\\
6. 点\\(diǎn)};
\node[left=0.3cm of mu, font=\tiny, popmuted, align=right] {7. 竖\\(垂直)\\
8. 横折\\(pli)\\
9. 横\\(横)\\
10. 横\\(横)\\
11. 横\\(横)};
\node[left=0.3cm of gong, font=\tiny, popmuted, align=right] {12. 撇\\(pié)\\
13. 横折弯钩\\(crochet)\\
14. 撇\\(pié)};
% Total
\node[below=0.3cm of gong, draw=popgreen, thick, fill=popgreenL, rounded corners=4pt, minimum width=3cm, align=center, font=\small] {\textbf{Total : 14 traits}};
\end{tikzpicture}
\end{popfigure}

\begin{itemize}
\item \textbf{Clé} : \zh{⺮} (bambou, 6 traits) — les anciens calculs se faisaient sur des baguettes de bambou.
\item \textbf{Étage 2} : \zh{目} (mù, œil, 5 traits) — « l'œil qui examine les comptes ».
\item \textbf{Étage 3} : \zh{廾} (gǒng, deux mains jointes, 3 traits) — « les mains qui manipulent les baguettes de calcul ».
\item \textbf{Total} : \textbf{14 traits}. Ordre : \emph{bambou (haut) → œil (milieu) → mains jointes (bas)}.
\item \textbf{Mots composés} : \zh{预算} (yùsuàn, budget), \zh{打算} (dǎsuàn, planifier, avoir l'intention), \zh{算数} (suànshù, compter, tenir parole), \zh{心算} (xīnsuàn, calcul mental).
\end{itemize}

\textbf{4. \zh{储蓄} (chǔxù) — épargner, épargne}\par
Deux caractères, deux clés différentes.

\begin{itemize}
\item \zh{储} (chǔ) : clé \zh{亻} (rén, homme) à gauche (2 traits) + phonétique \zh{宁} (níng, 10 traits dans la composition) = \textbf{12 traits}. Sens : « l'homme qui met de côté, qui stocke ».
\item \zh{蓄} (xù) : clé \zh{艹} (cǎo, herbe/plante) en haut (3 traits) + phonétique \zh{畜} (chù, 10 traits) = \textbf{13 traits}. Sens originel : « faire pousser/accumuler du bétail » → « accumuler, stocker ».
\item \textbf{Mots} : \zh{储蓄所} (chǔxùsuǒ, caisse d'épargne), \zh{储蓄罐} (chǔxùguàn, tirelire), \zh{存储} (cúnchǔ, stocker / stockage informatique).
\end{itemize}

\textbf{5. \zh{预算} (yùsuàn) — budget (nom), budgétiser (verbe)}\par
\begin{itemize}
\item \zh{预} (yù) : clé \zh{页} (yè, page/tête) à droite (6 traits) + phonétique \zh{予} (yǔ, 4 traits) = \textbf{10 traits}. Sens : « anticiper, prévoir » (l'image d'une page qu'on tourne à l'avance).
\item \zh{算} (suàn) : vu ci-dessus, \textbf{14 traits}.
\item \textbf{Total du mot} : 24 traits.
\item \textbf{Collocations} : \zh{预算表} (yùsuànbiǎo, tableau de budget), \zh{编制预算} (biānzhì yùsuàn, établir un budget), \zh{超支} (chāozhī, dépassement de budget), \zh{结余} (jiéyú, solde positif).
\end{itemize}

\begin{aretenir}[Récapitulatif des caractères du budget]
\begin{center}
\begin{tabularx}{\linewidth}{|X|X|X|X|X|}
\hline
\rowcolor{popblueL} \textbf{Caractère} & \textbf{Pinyin} & \textbf{Clé} & \textbf{Nb traits} & \textbf{Structure} \\
\hline
钱 & qián & 钅 (métal) & 12 & 左右 (gauche-droite) \\
费 & fèi & 贝 (coquillage) & 9 & 上下 (haut-bas) \\
算 & suàn & ⺮ (bambou) & 14 & 上中下 (trois étages) \\
储 & chǔ & 亻 (homme) & 12 & 左右 \\
蓄 & xù & 艹 (herbe) & 13 & 上下 \\
预 & yù & 页 (page/tête) & 10 & 左右 (clé à droite) \\
\hline
\end{tabularx}
\end{center}
\end{aretenir}

\courssub{Caractères proches à distinguer : pièges visuels et inversions de sens}

Les francophones, habitués à l'alphabet où l'ordre des lettres change peu la forme globale du mot, sous-estiment souvent la différence qu'un seul trait peut faire en chinois. Voici les quatre paires les plus dangereuses pour le thème « argent de poche / budget ».

\begin{attention}[买 (mǎi) vs 卖 (mài) : un trait, un sens inverse]
\begin{itemize}
\item \zh{买} (mǎi, \textbf{acheter}) : \textbf{6 traits}. Clé \zh{乛} (yǐ, second trait). Structure : \zh{乛} + \zh{卖} sans le 十 du haut. \emph{Mnémotechnique} : « J'achète, je \textbf{reçois} l'objet, je n'ai pas de croix (十) sur la tête. »
\item \zh{卖} (mài, \textbf{vendre}) : \textbf{8 traits}. Même base que \zh{买} \textbf{mais avec un 十 (shí, dix/croix) en plus au sommet}. \emph{Mnémotechnique} : « Je vends, je \textbf{barre} l'objet d'une croix, il ne m'appartient plus. »
\item \textbf{Piège tonal} : même syllabe \emph{mai}, mais \zh{买} est 3\up{e} ton (bas-remontant), \zh{卖} est 4\up{e} ton (chutant). Confondre les deux inverse complètement la transaction !
\end{itemize}
\end{attention}

\begin{attention}[钱 (qián) vs 线 (xiàn) : la clé change tout]
\begin{itemize}
\item \zh{钱} (qián) : clé \zh{钅} (métal) à gauche → \emph{argent, monnaie}.
\item \zh{线} (xiàn) : clé \zh{纟} (sī, fil/soie) à gauche → \emph{fil, ligne, câble, indice}.
\item \textbf{Différence} : \zh{钅} a un trait pointé (㇏) en 5\up{e} position ; \zh{纟} a un trait relevé (㇀) en 3\up{e} position puis un point. Ne confondez pas « \zh{花钱} \emph{huā qián} (dépenser de l'argent) » et « \zh{花线} \emph{huā xiàn} (fil à broder) » !
\end{itemize}
\end{attention}

\begin{attention}[借 (jiè) vs 错 (cuò) : même phonétique, clés différentes]
\begin{itemize}
\item \zh{借} (jiè) : clé \zh{亻} (homme) → \emph{emprunter, prêter} (action humaine). \zh{借钱} (jiè qián) = emprunter de l'argent ; \zh{借给} (jiè gěi) = prêter à quelqu'un.
\item \zh{错} (cuò) : clé \zh{钅} (métal) → \emph{erreur, faux, incorrect}. \zh{错字} (cuòzì) = faute de caractère.
\item \textbf{Piège} : même composant phonétique \zh{昔} (xī) à droite, mais la clé change le sens radicalement. « \zh{我借了错书} » (Wǒ jiè le cuò shū) ne veut rien dire ; on dirait « \zh{我借错了书} \emph{Wǒ jiè cuò le shū} » (J'ai emprunté le mauvais livre).
\end{itemize}
\end{attention}

\begin{attention}[费 (fèi) vs 佛 (fó) : même phonétique \zh{弗}, clés opposées]
\begin{itemize}
\item \zh{费} (fèi) : clé \zh{贝} (argent) en bas → \emph{dépense, frais}.
\item \zh{佛} (fó) : clé \zh{亻} (homme) à gauche → \emph{Bouddha, bouddhisme}.
\item \textbf{Exemple} : \zh{学费} (frais de scolarité) vs \zh{佛教} (bouddhisme). Ne dites pas « \zh{交佛} » pour « payer les frais » !
\end{itemize}
\end{attention}

\begin{popfigure}
\legende{Comparatif visuel des paires \zh{买}/\zh{卖} et \zh{钱}/\zh{线} — traits différenciants surlignés}
\begin{tikzpicture}[
    node distance=0.6cm and 1cm,
    char/.style={font=\fontsize{28}{32}\selectfont, text=popdark},
    diff/.style={font=\fontsize{28}{32}\selectfont, text=popred},
    label/.style={font=\small, text=popmuted, align=center},
    box/.style={draw=popline, thick, rounded corners=4pt, fill=popcream, minimum width=3.2cm, minimum height=2.2cm, align=center},
]
% Paire 买/卖
\node[box, align=center] (mai){
    \begin{tabular}{c}
    \zh{买} \\ \emph{mǎi} (acheter) \\ 6 traits
    \end{tabular}
};
\node[box, right=of mai, align=center] (mai2){
    \begin{tabular}{c}
    \zh{卖} \\ \emph{mài} (vendre) \\ 8 traits
    \end{tabular}
};
\node[label, below=0.2cm of mai] {Trait manquant : \textcolor{popred}{十} en haut};
\node[label, below=0.2cm of mai2] {Trait supplémentaire : \textcolor{popred}{十} en haut};

% Paire 钱/线
\node[box, below=1.2cm of mai, align=center] (qian){
    \begin{tabular}{c}
    \zh{钱} \\ \emph{qián} (argent) \\ clé \zh{钅}
    \end{tabular}
};
\node[box, right=of qian, align=center] (xian){
    \begin{tabular}{c}
    \zh{线} \\ \emph{xiàn} (fil/ligne) \\ clé \zh{纟}
    \end{tabular}
};
\node[label, below=0.2cm of qian] {5\up{e} trait : \textcolor{popred}{㇏} (pointé)};
\node[label, below=0.2cm of xian] {3\up{e} trait : \textcolor{popred}{㇀} (relevé) + point};
\end{tikzpicture}
\end{popfigure}

\courssub{Méthodologie d'apprentissage et entraînement guidé}

\begin{methode}[Comment mémoriser un caractère complexe de façon durable]
Suivez ces \textbf{5 étapes} pour chaque nouveau caractère du thème (et au-delà) :
\begin{enumerate}
\item \textbf{Décomposer} : identifiez la \cle{clé sémantique} (indice de sens) et le \cle{composant phonétique} (indice de son). Ex. : \zh{算} = \zh{⺮} (bambou/calcul) + \zh{目} (œil) + \zh{廾} (mains).
\item \textbf{Compter et nommer} : écrivez le caractère \textbf{trait par trait} en \emph{comptant à voix haute} en chinois (ou en français) : « 一 (yī), 丨 (èr), 丿 (sān), 丶 (sì)… ». Cela ancre l'ordre moteur.
\item \textbf{Visualiser l'étymologie} : imaginez la scène originelle (baguettes de bambou, œil qui vérifie, mains qui comptent). Dessinez un petit schéma si besoin.
\item \textbf{Écrire en grille} : recopiez le caractère \textbf{5 fois minimum} dans une grille \emph{田字格} (tiánzìgé), en respectant la proportion : clé plus petite, composant phonétique centré.
\item \textbf{Contextualiser} : écrivez \textbf{2 mots composés} et \textbf{1 phrase complète} contenant le caractère. Ex. : \zh{预算} → \zh{编制预算} → \zh{我们要编制下个月的预算。}
\end{enumerate}
\end{methode}

\begin{savaistu}[Le caractère \zh{银} (yín) et le mot \zh{银行} (yínháng)]
Le caractère \zh{银} (yín, argent métal / argent monétaire) contient lui aussi la clé \zh{钅} (métal) à gauche + phonétique \zh{艮} (gèn) à droite = \textbf{11 traits}. \zh{银行} (yínháng) signifie littéralement « magasin/établissement (行) de l'argent (银) » → \textbf{banque}. En Chine, les quatre grandes banques d'État s'appellent : \zh{中国工商银行} (ICBC), \zh{中国农业银行} (ABC), \zh{中国银行} (BOC), \zh{中国建设银行} (CCB). Au Cameroun, on trouve des succursales de \zh{中国银行} à Yaoundé et Douala.
\end{savaistu}

\courssub{Exemple résolu : dictée guidée et auto-correction}

\begin{exoresolu}[Écrire correctement une phrase de budget]
\textbf{Énoncé} : Écrivez en caractères chinois (avec pinyin et traduction) la phrase : « Mon budget mensuel est de 20\,000 francs CFA, je dois bien calculer mes dépenses. »

\textbf{Solution détaillée} :
\begin{enumerate}
\item \textbf{Analyse lexicale} :
\begin{itemize}
\item « budget » → \zh{预算} (yùsuàn)
\item « mensuel » → \zh{每月} (měiyuè) ou \zh{月度} (yuèdù)
\item « 20\,000 francs CFA » → \zh{两万中非法郎} (liǎng wàn Zhōngfēi fǎláng) — \emph{deux myriades de francs d'Afrique centrale}
\item « calculer » → \zh{计算} (jìsuàn) ou \zh{打算} (dǎsuàn) — ici \zh{计算} plus formel
\item « dépenses » → \zh{费用} (fèiyòng) ou \zh{开销} (kāixiāo)
\item « bien » → \zh{好好} (hǎohao) ou \zh{认真} (rènzhēn)
\end{itemize}
\item \textbf{Construction de la phrase} (ordre chinois : Sujet + Temps + Lieu + Verbe + Objet) :
\zh{我的每月预算是两万中非法郎，我要认真计算费用。}
\item \textbf{Pinyin} : Wǒ de měiyuè yùsuàn shì liǎng wàn Zhōngfēi fǎláng, wǒ yào rènzhēn jìsuàn fèiyòng.
\item \textbf{Vérification des caractères clés} :
\begin{itemize}
\item \zh{预} (10 traits, clé 页 à droite) ✓
\item \zh{算} (14 traits, ⺮ + 目 + 廾) ✓
\item \zh{费} (9 traits, 弗 + 贝, point final) ✓
\item \zh{用} (5 traits, clé 用 lui-même) ✓
\end{itemize}
\end{enumerate}

\textbf{Résultat final} :
\begin{center}
\zh{我的每月预算是两万中非法郎，我要认真计算费用。} \\
\emph{Wǒ de měiyuè yùsuàn shì liǎng wàn Zhōngfēi fǎláng, wǒ yào rènzhēn jìsuàn fèiyòng.} \\
« Mon budget mensuel est de 20\,000 francs CFA, je dois calculer sérieusement mes dépenses. »
\end{center}
\end{exoresolu}

\courssub{Entraînement à l'écriture : grille de révision}

\begin{exercice}[Entraînement à l'écriture des caractères du budget]
Recopiez \textbf{5 fois} chaque caractère ci-dessous dans votre cahier (idéalement en grille \emph{田字格}), en \textbf{respectant l'ordre des traits} et en \textbf{comptant les traits à voix haute}. Cochez chaque case après chaque écriture.

\begin{center}
\begin{tabularx}{\linewidth}{|X|c|c|c|c|c|}
\hline
\rowcolor{popblueL} \textbf{Caractère} & \textbf{1} & \textbf{2} & \textbf{3} & \textbf{4} & \textbf{5} \\
\hline
\zh{钱} (qián) — 12 traits & \square & \square & \square & \square & \square \\
\zh{费} (fèi) — 9 traits & \square & \square & \square & \square & \square \\
\zh{算} (suàn) — 14 traits & \square & \square & \square & \square & \square \\
\zh{储} (chǔ) — 12 traits & \square & \square & \square & \square & \square \\
\zh{蓄} (xù) — 13 traits & \square & \square & \square & \square & \square \\
\zh{预} (yù) — 10 traits & \square & \square & \square & \square & \square \\
\zh{买} (mǎi) — 6 traits & \square & \square & \square & \square & \square \\
\zh{卖} (mài) — 8 traits & \square & \square & \square & \square & \square \\
\hline
\end{tabularx}
\end{center}

\textbf{Consigne supplémentaire} : Pour chaque caractère, écrivez \textbf{un mot composé} et \textbf{une phrase courte} (3 à 5 mots) en caractères + pinyin + traduction.
\end{exercice}

\begin{corrige}[Exercice — Entraînement à l'écriture]
\textbf{Mots composés et phrases types} (d'autres réponses correctes sont possibles) :
\begin{itemize}
\item \zh{钱} → \zh{零钱} (língqián, monnaie) — \zh{我没有零钱。} (Wǒ méiyǒu língqián. — Je n'ai pas de monnaie.)
\item \zh{费} → \zh{学费} (xuéfèi, frais de scolarité) — \zh{学费很贵。} (Xuéfèi hěn guì. — Les frais de scolarité sont chers.)
\item \zh{算} → \zh{打算} (dǎsuàn, planifier) — \zh{我打算存钱。} (Wǒ dǎsuàn cún qián. — Je prévois d'épargner.)
\item \zh{储} → \zh{储蓄} (chǔxù, épargne) — \zh{储蓄很重要。} (Chǔxù hěn zhòngyào. — L'épargne est importante.)
\item \zh{蓄} → \zh{蓄水池} (xùshuǐchí, réservoir d'eau) — \zh{蓄水池很大。} (Xùshuǐchí hěn dà. — Le réservoir est grand.)
\item \zh{预} → \zh{预习} (yùxí, réviser/préparer la leçon) — \zh{请预习课文。} (Qǐng yùxí kèwén. — Préparez le texte.)
\item \zh{买} → \zh{买菜} (mǎicài, faire le marché) — \zh{妈妈去买菜。} (Māma qù mǎicài. — Maman va au marché.)
\item \zh{卖} → \zh{卖水果} (mài shuǐguǒ, vendre des fruits) — \zh{他卖水果。} (Tā mài shuǐguǒ. — Il vend des fruits.)
\end{itemize}
\textbf{Auto-évaluation} : Vérifiez que votre \zh{费} finit par un \textbf{point} (丶) et non un trait horizontal, que votre \zh{算} a bien \textbf{14 traits} (pas 13 !), que \zh{买} n'a \textbf{pas de 十} en haut et que \zh{卖} \textbf{en a un}.
\end{corrige}

\courssub{Texte d'entraînement intégré : « Mon carnet de budget »}

Pour consolider l'écriture et le vocabulaire dans un contexte authentique, voici un texte original d'environ 250 mots, rédigé par un élève de Terminale à Yaoundé. Lisez-le, notez les caractères ciblés, puis répondez aux questions.

\begin{textebox}[Mon carnet de budget — 我的预算本]
我叫保罗，住在雅温得。每个月初，爸爸给我两万中非法郎生活费。我有一个小本子，专门记预算。第一页写“收入”，第二页写“支出”。收入只有零花钱，支出有饭费、车费、买文具、买饮料、有时请同学吃烤鱼。\\
上个月，我预算两万，结果超支了三千。因为我买了一双贵的球鞋，还借给朋友五千块，他到现在还没还。这个月，我下定决心要省钱。我每天带饭盒去学校，不买饮料，走路回家不坐摩托车。我还准备开一个储蓄账户，把剩下的钱存进银行。\\
昨天，我算了一下：如果每天省五百，一个月能省一万五。算数真有用！我把数字写在本子上，贴在床头，提醒自己：钱是辛苦赚来的，花钱要算计，存钱要尽早。\\
\emph{Wǒ jiào Bǎoluó, zhù zài Yǎwēndé. Měi gè yuè chū, bàba gěi wǒ liǎng wàn Zhōngfēi fǎláng shēnghuófèi. Wǒ yǒu yí gè xiǎo běnzi, zhuānmén jì yùsuàn. Dì yī yè xiě "shōurù", dì èr yè xiě "zhīchū". Shōurù zhǐyǒu línghuáqián, zhīchū yǒu fànfèi, chēfèi, mǎi wénjù, mǎi yǐnliào, yǒushí qǐng tóngxué chī kǎoyú.}\\\emph{Shàng gè yuè, wǒ yùsuàn liǎng wàn, jiéguǒ chāozhī le sān qiān. Yīnwèi wǒ mǎi le yī shuāng guì de qiúxié, hái jiè gěi péngyou wǔ qiān kuài, tā dào xiànzài hái méi huán. Zhège yuè, wǒ xiàdìngjuéxīn yào shěngqián. Wǒ měitiān dài fànhé qù xuéxiào, bù mǎi yǐnliào, zǒulù huíjiā bù zuò mótuōchē. Wǒ hái zhǔnbèi kāi yī gè chǔxù zhànghù, bǎ shèngxià de qián cún jìn yínháng.}\\\emph{Zuótiān, wǒ suàn le yī xià: rúguǒ měitiān shěng wǔ bǎi, yí gè yuè néng shěng yí wàn wǔ. Suànshù zhēn yǒuyòng! Wǒ bǎ shùzì xiě zài běnzi shàng, tiē zài chuángtóu, tíxǐng zìjǐ: qián shì xīnkǔ zhuàn lái de, huāqián yào suànjì, cúnqián yào jìnzǎo.}
\end{textebox}

\begin{center}
\begin{tabularx}{\linewidth}{|X|X|X|X|}
\hline
\rowcolor{popblueL} \textbf{Caractères} & \textbf{Pinyin} & \textbf{Nature} & \textbf{Traduction} \\
\hline
收入 & shōurù & n. & revenus, rentrée d'argent \\
支出 & zhīchū & n. & dépenses, sortie d'argent \\
零花钱 & línghuáqián & n. & argent de poche \\
饭费 & fànfèi & n. & frais de repas / cantine \\
车费 & chēfèi & n. & frais de transport \\
文具 & wénjù & n. & fournitures scolaires \\
饮料 & yǐnliào & n. & boisson \\
烤鱼 & kǎoyú & n. & poisson grillé (spécialité camerounaise) \\
超支 & chāozhī & v. & dépasser le budget \\
球鞋 & qiúxié & n. & baskets, chaussures de sport \\
借 & jiè & v. & prêter (à qqn) / emprunter (à qqn) \\
还 & huán & v. & rembourser, rendre \\
省钱 & shěngqián & v. & économiser de l'argent \\
饭盒 & fànhé & n. & boîte à repas, gamelle \\
摩托车 & mótuōchē & n. & moto-taxi (benskin) \\
储蓄账户 & chǔxù zhànghù & n. & compte d'épargne \\
存 & cún & v. & déposer (argent), épargner \\
算 & suàn & v. & calculer \\
辛苦 & xīnkǔ & adj. & pénible, dur (travail) \\
算计 & suànjì & v. & calculer soigneusement, planifier \\
尽早 & jìnzǎo & adv. & le plus tôt possible \\
\hline
\end{tabularx}
\end{center}

\begin{exercice}[Compréhension et réemploi du texte]
\begin{enumerate}
\item Relevez dans le texte \textbf{tous les caractères contenant la clé 贝} et donnez leur sens.
\item Trouvez \textbf{deux paires de caractères proches} étudiées dans cette leçon (ex. : 买/卖, 钱/线…) qui apparaissent ou sont suggérées par le contexte.
\item Traduisez en français la phrase : \zh{钱是辛苦赚来的，花钱要算计，存钱要尽早。}
\item \textbf{Production écrite} : Sur le modèle de Paul, rédigez \textbf{6 à 8 lignes} en caractères chinois (avec pinyin) pour présenter \emph{votre} gestion de l'argent de poche du mois dernier : revenus, dépenses principales, dépassement ou épargne, résolution pour ce mois-ci. Utilisez au moins \textbf{8 caractères du vocabulaire ci-dessus}.
\end{enumerate}
\end{exercice}

\begin{corrige}[Exercice — Compréhension et réemploi]
\begin{enumerate}
\item \textbf{Caractères à clé 贝 dans le texte} :
\begin{itemize}
\item \zh{费} (fèi, frais) — présent dans \zh{饭费} (fànfèi), \zh{车费} (chēfèi), \zh{生活费} (shēnghuófèi).
\item \zh{贵} (guì, cher) — présent dans \zh{贵的球鞋} (guì de qiúxié).
\item \zh{账} (zhàng, compte) — présent dans \zh{储蓄账户} (chǔxù zhànghù).
\item \zh{赚} (zhuàn, gagner) — présent dans \zh{赚来的} (zhuàn lái de).
\end{itemize}
\emph{Note : \zh{支出} ne contient pas la clé 贝 ; \zh{贫}、\zh{购}、\zh{资} n'apparaissent pas dans le texte.}

\item \textbf{Paires proches présentes ou suggérées} :
\begin{itemize}
\item \zh{买} (mǎi, acheter) / \zh{卖} (mài, vendre) : \zh{买} apparaît 3 fois (\zh{买文具}、\zh{买饮料}、\zh{买了一双}) ; \zh{卖} est le contraire sémantique étudié en classe.
\item \zh{借} (jiè, emprunter/prêter) / \zh{错} (cuò, erreur) : \zh{借给} (jiè gěi, prêter à) apparaît dans le texte (\zh{借给朋友}) ; \zh{错} est le piège visuel associé.
\item \zh{钱} (qián, argent) / \zh{线} (xiàn, fil) : \zh{钱} apparaît 3 fois ; \zh{线} est le piège visuel associé (clé \zh{纟} vs \zh{钅}).
\item \zh{费} (fèi, frais) / \zh{佛} (fó, Bouddha) : \zh{费} apparaît 3 fois ; \zh{佛} est le piège visuel associé (clé \zh{贝} vs \zh{亻}).
\end{itemize}

\item \textbf{Traduction} : « L'argent est durement gagné, il faut réfléchir avant de dépenser, et épargner le plus tôt possible. »
\begin{itemize}
\item \zh{钱是辛苦赚来的} : L'argent est gagné par le travail pénible / à la sueur de son front.
\item \zh{花钱要算计} : Dépenser de l'argent demande de calculer / planifier / faire attention.
\item \zh{存钱要尽早} : Épargner de l'argent doit se faire le plus tôt possible.
\end{itemize}

\item \textbf{Production écrite — Modèle de réponse} :
\begin{center}
\zh{我叫玛丽，住在杜阿拉。上个月，妈妈给我一万五中非法郎零花钱。我买了文具、饮料，还请朋友吃烤鱼，超支了两千。这个月，我决定省钱。我不买饮料，带饭盒去学校，走路不坐摩托车。我想开储蓄账户，存钱买球鞋。钱是辛苦赚来的，要算计着花。} \\
\emph{Wǒ jiào Mǎlì, zhù zài Dùālā. Shàng gè yuè, māma gěi wǒ yí wàn wǔ Zhōngfēi fǎláng línghuáqián. Wǒ mǎi le wénjù, yǐnliào, hái qǐng péngyou chī kǎoyú, chāozhī le liǎng qiān. Zhège yuè, wǒ juédìng shěngqián. Wǒ bù mǎi yǐnliào, dài fànhé qù xuéxiào, zǒulù bù zuò mótuōchē. Wǒ xiǎng kāi chǔxù zhànghù, cún qián mǎi qiúxié. Qián shì xīnkǔ zhuàn lái de, yào suànjì zhe huā.} \\
« Je m'appelle Marie, j'habite à Douala. Le mois dernier, maman m'a donné 15\,000 FCFA de poche. J'ai acheté fournitures, boissons, et invité un ami pour du poisson grillé, dépassant de 2000. Ce mois-ci, je décide d'économiser. Pas de boissons, j'apporte ma gamelle, je marche pas de moto. Je veux ouvrir un compte épargne, économiser pour des baskets. L'argent est durement gagné, il faut le dépenser avec réflexion. »
\end{center}
\textbf{Vocabulaire utilisé (10 items)} : 零花钱、文具、饮料、烤鱼、超支、省钱、饭盒、摩托车、储蓄账户、存钱、算计.
\end{enumerate}
\end{corrige}

\coursec{Lexique actif et collocations du thème}

Dans la section précédente, nous avons appris à écrire correctement les caractères du budget — \zh{钱}, \zh{花}, \zh{省}, \zh{存}, \zh{预}, \zh{算} — en respectant les traits et les radicaux. Savoir écrire un caractère, c'est bien ; savoir l'\emph{employer} dans une phrase, avec les bons partenaires de mots, c'est mieux. C'est précisément l'objectif de cette section : transformer votre vocabulaire passif en \cle{lexique actif}, c'est-à-dire en mots que vous pouvez mobiliser spontanément dans vos productions écrites sur la gestion de l'argent de poche.

\courssub{Observer d'abord : un texte pour faire émerger le lexique}

\begin{textebox}[我不乱花钱，因为我有一个预算]
我不乱花钱，因为我有一个预算。\\
Wǒ bù luàn huā qián, yīnwèi wǒ yǒu yí ge yùsuàn.\\
Je ne dépense pas n'importe comment, parce que j'ai un budget.\\[4pt]
每个月，妈妈给我一千法郎的零花钱。\\
Měi ge yuè, māma gěi wǒ yì qiān fǎláng de línghuāqián.\\
Chaque mois, maman me donne mille francs d'argent de poche.\\[4pt]
我先算一算这个月的费用：坐车、买书、打电话。\\
Wǒ xiān suàn yi suàn zhè ge yuè de fèiyòng : zuò chē, mǎi shū, dǎ diànhuà.\\
Je calcule d'abord les frais du mois : le transport, les livres, le téléphone.\\[4pt]
然后，我存一点儿钱。储蓄很重要！\\
Ránhòu, wǒ cún yìdiǎnr qián. Chǔxù hěn zhòngyào !\\
Ensuite, je mets un peu d'argent de côté. L'épargne, c'est très important !\\[4pt]
我觉得管理零花钱很有用。\\
Wǒ juéde guǎnlǐ línghuāqián hěn yǒuyòng.\\
Je pense que gérer son argent de poche est très utile.
\end{textebox}

\textbf{Questions d'observation.} 1) Quels mots du texte désignent l'argent ? 2) Quels verbes accompagnent le mot \zh{钱} ? 3) Quelle phrase exprime une opinion ? Soulignez-la : elle servira de modèle dans vos rédactions.

\textbf{Réponses attendues.} 1) \zh{钱}, \zh{零花钱}, \zh{储蓄}. 2) \zh{花} (dépenser), \zh{存} (mettre de côté). 3) \zh{我觉得管理零花钱很有用。}

\courssub{Les noms du thème : le vocabulaire de base}

\begin{definition}[Les cinq noms essentiels du budget]
\zh{零花钱} (línghuāqián) : l'argent de poche, littéralement « argent pour les petites dépenses ». \zh{预算} (yùsuàn) : le budget. \zh{费用} (fèiyòng) : les frais, les dépenses prévues. \zh{价格} (jiàgé) : le prix (affiché, d'un produit). \zh{储蓄} (chǔxù) : l'épargne (nom abstrait, l'argent mis de côté).
\end{definition}

\begin{center}
\begin{tabularx}{\linewidth}{|X|X|X|X|}
\hline
\rowcolor{popblueL} Caractères & Pinyin & Nature & Traduction \\
\hline
零花钱 & línghuāqián & nom & argent de poche \\
\hline
预算 & yùsuàn & nom & budget \\
\hline
费用 & fèiyòng & nom & frais, dépenses \\
\hline
价格 & jiàgé & nom & prix \\
\hline
储蓄 & chǔxù & nom & épargne \\
\hline
贵 & guì & adjectif & cher \\
\hline
便宜 & piányi & adjectif & bon marché \\
\hline
重要 & zhòngyào & adjectif & important \\
\hline
有用 & yǒuyòng & adjectif & utile \\
\hline
乱 & luàn & adverbe & n'importe comment, en désordre \\
\hline
\end{tabularx}
\end{center}

\begin{attention}[价格 ou 钱 ?]
Le français « prix » a deux traductions selon le contexte. \zh{价格} désigne le prix comme concept : \zh{米的价格很高。} \emph{Mǐ de jiàgé hěn gāo.} « Le prix du riz est élevé. » Mais pour demander un prix concret, on dit \zh{多少钱？} \emph{Duōshao qián ?} « Combien (ça coûte) ? ». Ne dites jamais \zh{多少价格} ! De même, « ce livre coûte cher » se dit \zh{这本书很贵。} et non \zh{这本书的价格很贵} dans le langage courant (on dirait alors \zh{价格很高}).
\end{attention}

\courssub{Les verbes et leurs collocations : le cœur du lexique actif}

En chinois, un verbe s'apprend \emph{toujours avec son complément habituel}. Cette combinaison stable verbe + nom s'appelle une \cle{collocation}. C'est elle qui fait la différence entre un texte d'élève moyen et un texte excellent.

\begin{center}
\begin{tabularx}{\linewidth}{|X|X|X|X|}
\hline
\rowcolor{popblueL} Caractères & Pinyin & Nature & Traduction \\
\hline
花 & huā & verbe & dépenser, passer (du temps) \\
\hline
省 & shěng & verbe & économiser \\
\hline
存 & cún & verbe & mettre de côté, déposer \\
\hline
借 & jiè & verbe & emprunter, prêter \\
\hline
算 & suàn & verbe & calculer, compter \\
\hline
管理 & guǎnlǐ & verbe & gérer, administrer \\
\hline
\end{tabularx}
\end{center}

\begin{definition}[Les cinq collocations à connaître par cœur]
\zh{花钱} (huā qián) : dépenser de l'argent. \zh{省钱} (shěng qián) : économiser de l'argent. \zh{存钱} (cún qián) : épargner, mettre de l'argent de côté. \zh{乱花钱} (luàn huā qián) : dépenser n'importe comment (\zh{乱} = en désordre, n'importe comment). \zh{管理零花钱} (guǎnlǐ línghuāqián) : gérer son argent de poche.
\end{definition}

\begin{exemplebox}[Les collocations en contexte]
\zh{我不乱花钱。} \emph{Wǒ bù luàn huā qián.} « Je ne dépense pas n'importe comment. »\\[3pt]
\zh{为了省钱，我走路去学校。} \emph{Wèile shěng qián, wǒ zǒulù qù xuéxiào.} « Pour économiser de l'argent, je vais à l'école à pied. »\\[3pt]
\zh{我每个月存一点儿钱。} \emph{Wǒ měi ge yuè cún yìdiǎnr qián.} « Je mets chaque mois un peu d'argent de côté. »\\[3pt]
\zh{中学生应该学会管理零花钱。} \emph{Zhōngxuéshēng yīnggāi xuéhuì guǎnlǐ línghuāqián.} « Les collégiens et lycéens devraient apprendre à gérer leur argent de poche. »
\end{exemplebox}

\begin{attention}[花 : un seul verbe pour l'argent ET le temps]
Les francophones cherchent souvent deux verbes différents pour « dépenser » (argent) et « passer » (temps). En chinois, c'est le même : \zh{花钱} (huā qián) = dépenser de l'argent ; \zh{花时间} (huā shíjiān) = passer du temps. Exemple : \zh{做作业花很多时间。} \emph{Zuò zuòyè huā hěn duō shíjiān.} « Faire les devoirs prend beaucoup de temps. » Une seule structure : \cle{花 + durée ou somme (+ verbe)}.
\end{attention}

\begin{attention}[借 : emprunter ou prêter ?]
\zh{借} (jiè) signifie à la fois « emprunter » et « prêter » ; c'est la préposition qui tranche : \zh{我向朋友借钱。} \emph{Wǒ xiàng péngyou jiè qián.} « J'emprunte de l'argent à un ami » (\zh{向} = auprès de) ; \zh{我借给朋友一百法郎。} \emph{Wǒ jiè gěi péngyou yì bǎi fǎláng.} « Je prête cent francs à un ami » (\zh{给} marque le bénéficiaire). Retenez : \cle{向 + personne + 借} = emprunter ; \cle{借 + 给 + personne} = prêter.
\end{attention}

\courssub{Carte mentale : organiser le lexique autour de 钱}

\begin{popfigure}
\begin{center}
\begin{center}\tcbox[colback=popink,colframe=popink,coltext=white,arc=9pt,boxsep=4pt,left=6pt,right=6pt]{\bfseries\small 钱 qián l'argent}\end{center}
\vspace{-2pt}
\begin{tcbraster}[raster columns=3,raster equal height=rows,raster column skip=6pt,raster row skip=6pt,enhanced,blankest]
\begin{tcolorbox}[enhanced,colback=white,colframe=popblue,boxrule=0.9pt,arc=3pt,title={花 huā dépenser},fonttitle=\bfseries\scriptsize,coltitle=white,colbacktitle=popblue,left=4pt,right=4pt,top=2pt,bottom=2pt,boxsep=1pt,toptitle=1.5pt,bottomtitle=1.5pt,halign=left]
\begin{itemize}[nosep,leftmargin=*,itemsep=1pt,font=\scriptsize]
\item 花钱买书 huā qián mǎi shū dépenser pour des livres
\end{itemize}
\end{tcolorbox}
\begin{tcolorbox}[enhanced,colback=white,colframe=poporange,boxrule=0.9pt,arc=3pt,title={省 shěng économiser},fonttitle=\bfseries\scriptsize,coltitle=white,colbacktitle=poporange,left=4pt,right=4pt,top=2pt,bottom=2pt,boxsep=1pt,toptitle=1.5pt,bottomtitle=1.5pt,halign=left]
\begin{itemize}[nosep,leftmargin=*,itemsep=1pt,font=\scriptsize]
\item 走路省钱 zǒulù shěng qián marcher pour économiser
\end{itemize}
\end{tcolorbox}
\begin{tcolorbox}[enhanced,colback=white,colframe=popgreen,boxrule=0.9pt,arc=3pt,title={存 cún mettre de côté},fonttitle=\bfseries\scriptsize,coltitle=white,colbacktitle=popgreen,left=4pt,right=4pt,top=2pt,bottom=2pt,boxsep=1pt,toptitle=1.5pt,bottomtitle=1.5pt,halign=left]
\begin{itemize}[nosep,leftmargin=*,itemsep=1pt,font=\scriptsize]
\item 每月存钱 měi yuè cún qián épargner chaque mois
\end{itemize}
\end{tcolorbox}
\begin{tcolorbox}[enhanced,colback=white,colframe=poppurple,boxrule=0.9pt,arc=3pt,title={借 jiè emprunter},fonttitle=\bfseries\scriptsize,coltitle=white,colbacktitle=poppurple,left=4pt,right=4pt,top=2pt,bottom=2pt,boxsep=1pt,toptitle=1.5pt,bottomtitle=1.5pt,halign=left]
\begin{itemize}[nosep,leftmargin=*,itemsep=1pt,font=\scriptsize]
\item 向朋友借钱 xiàng péngyou jiè qián emprunter à un ami
\end{itemize}
\end{tcolorbox}
\begin{tcolorbox}[enhanced,colback=white,colframe=popgold,boxrule=0.9pt,arc=3pt,title={算 suàn calculer},fonttitle=\bfseries\scriptsize,coltitle=white,colbacktitle=popgold,left=4pt,right=4pt,top=2pt,bottom=2pt,boxsep=1pt,toptitle=1.5pt,bottomtitle=1.5pt,halign=left]
\begin{itemize}[nosep,leftmargin=*,itemsep=1pt,font=\scriptsize]
\item 算费用 suàn fèiyòng calculer les frais
\end{itemize}
\end{tcolorbox}
\end{tcbraster}

\end{center}
\legende{Carte mentale du lexique de l'argent : un verbe, une collocation, un exemple.}
\end{popfigure}

\courssub{Exprimer les quantités d'argent}

\begin{propriete}[Structure de la quantité d'argent]
\cle{Nombre + (unité monétaire) + 钱 / 法郎}. À l'écrit et dans les contextes camerounais : \zh{一千法郎} (yì qiān fǎláng) « mille francs ». À l'oral chinois, l'unité familière est \zh{块} (kuài) : \zh{五块钱} (wǔ kuài qián) « cinq yuans ». Pour les quantités vagues : \zh{很多钱} (hěn duō qián) « beaucoup d'argent » ; \zh{一点儿钱} (yìdiǎnr qián) « un peu d'argent ». Attention à la mutation de ton de \zh{一} : il se prononce \emph{yì} devant un 1\ier{}, 2\ieme{} ou 3\ieme{} ton (\zh{一千} yì qiān, \zh{一百} yì bǎi) et \emph{yí} devant un 4\ieme{} ton (\zh{一个} yí ge).
\end{propriete}

\begin{exemplebox}[Quantités en contexte]
\zh{我的零花钱不多，只有一千法郎。} \emph{Wǒ de línghuāqián bù duō, zhǐyǒu yì qiān fǎláng.} « Mon argent de poche n'est pas beaucoup, seulement mille francs. »\\[3pt]
\zh{这本书很便宜，只要五块钱。} \emph{Zhè běn shū hěn piányi, zhǐ yào wǔ kuài qián.} « Ce livre est bon marché, il ne coûte que cinq yuans. »\\[3pt]
\zh{坐车太贵了，所以我走路去学校。} \emph{Zuò chē tài guì le, suǒyǐ wǒ zǒulù qù xuéxiào.} « Le transport est trop cher, donc je vais à l'école à pied. »
\end{exemplebox}

Notez au passage la structure d'opinion \cle{太 + adjectif + 了} (trop...) et le connecteur de conséquence \zh{所以} (suǒyǐ, donc) : deux outils précieux pour vos rédactions.

\courssub{La phrase-type à réutiliser et la méthode d'enrichissement}

\begin{aretenir}[Phrase-type d'opinion]
\zh{我觉得管理零花钱很重要。} \emph{Wǒ juéde guǎnlǐ línghuāqián hěn zhòngyào.} « Je pense que gérer son argent de poche est très important. » Structure : \cle{我觉得 + proposition (verbe nominalisé) + 很 + adjectif}. Le groupe verbal \zh{管理零花钱} fonctionne comme un nom : c'est le sujet de \zh{重要}. Variantes d'adjectifs possibles : \zh{有用} (utile), \zh{有意思} (intéressant), \zh{必要} (nécessaire).
\end{aretenir}

\begin{methode}[Enrichir un texte simple en quatre gestes]
\begin{enumerate}
\item Repérez les adjectifs faibles : \zh{很好} (hěn hǎo, très bien) est banal.
\item Remplacez par un adjectif précis du thème : \zh{很重要} (très important) ou \zh{很有用} (très utile).
\item Ajoutez une collocation du thème : \zh{乱花钱}, \zh{省钱}, \zh{存钱}, \zh{管理零花钱}.
\item Ajoutez une justification avec \zh{因为} (parce que) ou une conséquence avec \zh{所以} (donc).
\end{enumerate}
Avant : \zh{零花钱很好。} Après : \zh{我觉得管理零花钱很重要，因为我可以省钱，不乱花钱。} \emph{Wǒ juéde guǎnlǐ línghuāqián hěn zhòngyào, yīnwèi wǒ kěyǐ shěng qián, bù luàn huā qián.} « Je pense que gérer son argent de poche est très important, parce que je peux économiser et ne pas dépenser n'importe comment. »
\end{methode}

\begin{exoresolu}[Transformer une phrase pauvre en phrase riche]
Énoncé : améliorez la phrase \zh{预算很好。} (Yùsuàn hěn hǎo. — Le budget, c'est bien.) en utilisant un adjectif du thème, une collocation et un connecteur.
\tcblower
Solution détaillée : étape 1, on remplace \zh{很好} par \zh{很有用} (très utile). Étape 2, on introduit la collocation \zh{存钱} ou \zh{乱花钱}. Étape 3, on relie avec \zh{因为}. Résultat : \zh{预算很有用，因为有了预算，我就不乱花钱。} \emph{Yùsuàn hěn yǒuyòng, yīnwèi yǒu le yùsuàn, wǒ jiù bù luàn huā qián.} « Un budget est très utile, parce que quand on a un budget, on ne dépense plus n'importe comment. » On remarque la corrélation \zh{因为……就……} (parce que... alors...), typique du chinois : \zh{就} annonce la conséquence logique dans la seconde proposition.
\end{exoresolu}

\begin{exercice}[Lexique actif : à vous de jouer]
1) Traduisez en chinois : « J'économise de l'argent parce que le transport est trop cher. » 2) Complétez avec la bonne collocation : \zh{我不……，因为我有一个预算。} 3) Construisez une phrase d'opinion avec \zh{我觉得……很重要} sur le thème de l'épargne. 4) Traduisez : « Pour économiser de l'argent, je vais au marché à pied. »
\end{exercice}

\begin{corrige}[Exercice « Lexique actif »]
1) \zh{我省钱，因为坐车太贵了。} \emph{Wǒ shěng qián, yīnwèi zuò chē tài guì le.} 2) \zh{乱花钱} : \zh{我不乱花钱，因为我有一个预算。} \emph{Wǒ bù luàn huā qián, yīnwèi wǒ yǒu yí ge yùsuàn.} 3) Modèle possible : \zh{我觉得存钱很重要。} \emph{Wǒ juéde cún qián hěn zhòngyào.} « Je pense qu'épargner est très important. » 4) \zh{为了省钱，我走路去市场。} \emph{Wèile shěng qián, wǒ zǒulù qù shìchǎng.}
\end{corrige}

Vous disposez désormais d'un lexique actif complet : cinq noms, six verbes, cinq collocations, les quantités d'argent et une phrase-type d'opinion. La section suivante mettra tout cela en pratique dans le thème (français → chinois).

\coursec{Leçon 27 — Expression écrite : Gestion de l'argent de poche et caractères complexes du budget}

\courssub{Modèle de texte commenté : « Mon budget de lycéen à Yaoundé »}

\begin{textebox}[Texte original -- Niveau HSK 3 / Tle A]
\zh{我叫保罗，是雅温得一所高中的终班学生。每个星期一，爸爸给我两千法郎零花钱。首先，我存五百法郎，因为我想买一本新汉法词典。然后，我用八百法郎买早饭，坐车去学校也要花钱。我不乱买零食，也不乱玩手机游戏。上个星期五，我存了两千法郎，现在我已经有一万法郎了。买词典还差三千法郎，所以我下个星期要继续努力存钱。} \\
\\
Wǒ jiào Bǎoluó, shì Yǎwēndé yī suǒ gāozhōng de zhōngbān xuéshēng. Měi gè xīngqīyī, bàba gěi wǒ liǎngqiān fǎláng línghuáqián. Shǒuxiān, wǒ cún wǔbǎi fǎláng, yīnwèi wǒ xiǎng mǎi yī běn xīn Hàn-Fǎ cídiǎn. Ránhòu, wǒ yòng bābǎi fǎláng mǎi zǎofàn, zuò chē qù xuéxiào yě yào huā qián. Wǒ bù luàn mǎi língshí, yě bù luàn wán shǒujī yóuxì. Shàng gè xīngqīwǔ, wǒ cún le liǎngqiān fǎláng, xiànzài wǒ yǐjīng yǒu yī wàn fǎláng le. Mǎi cídiǎn hái chà sānqiān fǎláng, suǒyǐ wǒ xià gè xīngqī yào jìxù nǔlì cún qián. \\
\\
« Je m'appelle Paul, je suis élève de Terminale dans un lycée de Yaoundé. Chaque lundi, papa me donne 2000 F d'argent de poche. D'abord, j'en économise 500 F, car je veux acheter un nouveau dictionnaire chinois-français. Ensuite, j'utilise 800 F pour le petit-déjeuner, et le transport pour aller à l'école coûte aussi de l'argent. Je n'achète pas de snacks n'importe comment, et je ne joue pas n'importe comment aux jeux sur le téléphone. Vendredi dernier, j'ai économisé 2000 F, maintenant j'ai déjà 10 000 F. Il me manque encore 3000 F pour le dictionnaire, donc la semaine prochaine je dois continuer à m'efforcer d'économiser. »
\end{textebox}

\begin{center}
\begin{tabularx}{\linewidth}{|X|X|X|X|}\hline
\rowcolor{popblueL} \textbf{Caractères} & \textbf{Pinyin} & \textbf{Nature} & \textbf{Traduction} \\ \hline
雅温得 & Yǎwēndé & n.pr. & Yaoundé \\ \hline
终班 & zhōngbān & n. & Terminale (classe terminale) \\ \hline
汉法词典 & Hàn-Fǎ cídiǎn & n. & Dictionnaire chinois-français \\ \hline
用 & yòng & v. & Utiliser, employer \\ \hline
零食 & língshí & n. & Snacks, grignotines \\ \hline
手机游戏 & shǒujī yóuxì & n. & Jeux mobiles / sur téléphone \\ \hline
上个星期五 & shàng gè xīngqīwǔ & t. & Vendredi dernier \\ \hline
已经 & yǐjīng & adv. & Déjà (souvent avec 了) \\ \hline
一万 & yī wàn & num. & Dix mille (10 000) \\ \hline
差 & chà & v. & Manquer (quantité), il manque \\ \hline
继续 & jìxù & v. & Continuer \\ \hline
努力 & nǔlì & adv./v. & Faire des efforts, avec effort \\ \hline
\end{tabularx}
\end{center}

\textbf{Analyse du texte (Structure globale) :}
\begin{enumerate}
    \item \textbf{Présentation} : Sujet (\zh{我}) + Identité (\zh{叫保罗，是...学生}) + Lieu (\zh{雅温得}).
    \item \textbf{Revenu récurrent} : Temps (\zh{每个星期一}) + Source (\zh{爸爸}) + Verbe \zh{给} (donner) + Bénéficiaire \zh{我} + Objet \zh{两千法郎零花钱}.
    \item \textbf{Allocation (Épargne)} : Connecteur \zh{首先} + Verbe \zh{存} + Quantité + Cause \zh{因为...所以...} (implicitement lié à la phrase suivante).
    \item \textbf{Dépenses courantes} : Connecteur \zh{然后} + Verbe \zh{用} (utiliser) + Montant + Objet \zh{买早饭} + Dépense supplémentaire \zh{坐车...也要花钱}.
    \item \textbf{Discipline financière} : Négation de manière \zh{不乱} + Verbes \zh{买} / \zh{玩}.
    \item \textbf{Bilan passé / État présent} : Temps passé \zh{上个星期五} + Accompli \zh{存了} + Situation actuelle \zh{现在...已经...了}.
    \item \textbf{Objectif futur} : Manque \zh{差} + Conséquence \zh{所以} + Futur \zh{要继续努力}.
\end{enumerate}

\courssub{Structure et connecteurs du texte}

\begin{propriete}[Règle d'or : Place du complément de temps]
En français, le complément de temps se place souvent en fin de phrase. En chinois, il \textbf{doit précéder le verbe} (ou le sujet s'il est thématisé).
\medskip
\noindent\textbf{Schéma :} \cle{Sujet + Temps + Lieu + Manière/But + Verbe + Objet}
\medskip
\begin{center}
\begin{tabularx}{\linewidth}{|X|X|X|}\hline
\rowcolor{popblueL} \textbf{Français} & \textbf{Chinois (Structure)} & \textbf{Exemple} \\ \hline
Je mange \emph{tous les jours}. & 我 \textbf{每天} 吃饭。 & Wǒ \textbf{měitiān} chīfàn. \\ \hline
Il étudie \emph{chaque soir}. & 他 \textbf{每个晚上} 学习。 & Tā \textbf{měi gè wǎnshàng} xuéxí. \\ \hline
Nous recevons \emph{chaque semaine}. & 我们 \textbf{每个星期} 得到... & Wǒmen \textbf{měi gè xīngqī} dédào... \\ \hline
\end{tabularx}
\end{center}
\end{propriete}

\begin{attention}[Piège n°1 : \emph{每个星期} en fin de phrase]
\emph{Jamais} : \zh{*我得到零花钱每个星期。} (Erreur d'ordre des mots calquée sur le français).
\emph{Toujours} : \zh{我每个星期得到零花钱。} ou \zh{每个星期我都得到零花钱。} (Thématisation du temps avec \zh{都}).
\end{attention}

\begin{propriete}[Connecteurs de séquence : \zh{首先}... \zh{然后}... \zh{接着}... \zh{最后}...]
Pour énumérer des actions chronologiques, le chinois utilise des adverbes de liaison placés en \textbf{début de proposition}, suivis d'une virgule.
\medskip
\begin{center}
\begin{tabularx}{\linewidth}{|X|X|X|}\hline
\rowcolor{poporangeL} \textbf{Français} & \textbf{Chinois} & \textbf{Pinyin} \\ \hline
D'abord / Premièrement & \zh{首先} & shǒuxiān \\ \hline
Ensuite / Puis & \zh{然后} & ránhòu \\ \hline
Après / Par la suite & \zh{接着} & jiēzhe \\ \hline
Finalement / À la fin & \zh{最后} & zuìhòu \\ \hline
\end{tabularx}
\end{center}
\emph{Note :} Le sujet \zh{我} est répété pour la clarté, mais peut être omis à la 2e occurrence si identique : \zh{首先我买早饭，然后坐车。}
\end{propriete}

\begin{propriete}[Structure cause-conséquence : \zh{因为}... \zh{所以}...]
Structure standard pour la causalité logique (écrit/formel). \zh{因为} introduit la cause (souvent en tête de phrase), \zh{所以} la conséquence.
\medskip
\noindent\textbf{Schéma :} \zh{因为} + Cause , \zh{所以} + Conséquence.
\medskip
\begin{center}
\begin{tabularx}{\linewidth}{|X|X|}\hline
\rowcolor{poppurpleL} \textbf{Français} & \textbf{Chinois} \\ \hline
Comme il pleut, je ne sors pas. & \zh{因为下雨，所以我不出去。} \\ \hline
Puisque tu as faim, mange. & \zh{因为你饿了，所以吃饭吧。} \\ \hline
\end{tabularx}
\end{center}
\emph{Variante orale :} \zh{所以}... \zh{是因为}... (C'est parce que...).
\end{propriete}

\begin{exemplebox}[Vocabulaire clé : \zh{零花钱} (línghuáqián) -- Argent de poche]
Litt. « argent menu/frais pour fleurs ».
\begin{itemize}
    \item \zh{给零花钱} (gěi línghuáqián) : Donner de l'argent de poche.
    \item \zh{花零花钱} (huā línghuáqián) : Dépenser son argent de poche.
    \item \zh{存零花钱} (cún línghuáqián) : Économiser son argent de poche.
    \item \zh{省零花钱} (shěng línghuáqián) : Faire des économies sur son argent de poche (dépenser moins).
\end{itemize}
\end{exemplebox}

\courssub{Écriture correcte des caractères du budget : traits, radicaux, caractères proches}

\begin{methode}[Méthode d'écriture et mémorisation des caractères clés]
\begin{enumerate}
    \item \textbf{Identifier le radical (clé)} : Il donne souvent l'indice sémantique (ex: \zh{贝} bèi = coquillage/argent pour \zh{费} \zh{贷} \zh{赚}).
    \item \textbf{Respecter l'ordre des traits} : Horizontal avant vertical (\zh{十}) ; Gauche avant droite (\zh{人}) ; Haut avant bas (\zh{三}) ; Extérieur avant intérieur (\zh{国}) ; Traits centraux avant latéraux symétriques (\zh{小}) ; Point final en haut à droite souvent dernier.
    \item \textbf{Comparer les caractères proches} (形近字) pour éviter les confusions.
    \item \textbf{S'exercer} : Écrire 5 à 10 fois en dictant le nom des traits ou l'ordre.
\end{enumerate}
\end{methode}

\begin{propriete}[Analyse des caractères complexes du thème « Budget »]
\begin{center}
\begin{tabularx}{\linewidth}{|X|X|X|X|X|}\hline
\rowcolor{popgreenL} \textbf{Caractère} & \textbf{Pinyin} & \textbf{Radical (Clé)} & \textbf{Nb traits} & \textbf{Caractères proches / Pièges} \\ \hline
\zh{存} & cún & \zh{子} (zǐ, enfant) & 6 & \zh{在} (zài, être à) : \zh{存} a \zh{子} en bas, \zh{在} a \zh{土} (terre). \textit{Sens : déposer, conserver.} \\ \hline
\zh{预} & yù & \zh{页} (yè, tête/page) & 13 & \zh{序} (xù, ordre) : même structure haut/bas, composant phonétique diff. (\zh{予} vs \zh{予}...). \textit{Sens : à l'avance, budget \zh{预算} (yùsuàn).} \\ \hline
\zh{算} & suàn & \zh{竹} (zhú, bambou) & 14 & \zh{管} (guǎn, tube/gérer) : même radical \zh{竹}, phonétique \zh{工} vs \zh{官}. \textit{Sens : calculer, \zh{算账} (suànzhàng).} \\ \hline
\zh{支} & zhī & \zh{又} (yòu, main droite) / \zh{扌} (shǒu) & 7 & \zh{枝} (zhī, branche) : même phonétique, radical \zh{木} (bois). \textit{Sens : payer, brancher ; classifieur pour stylos, factures.} \\ \hline
\zh{出} & chū & \zh{凵} (kǎn, réceptacle) & 5 & \zh{山} (shān, montagne) : traits verticaux différents. \textit{Sens : sortir, dépense \zh{支出} (zhīchū).} \\ \hline
\zh{入} & rù & \zh{入} (rù, entrer) & 2 & \zh{人} (rén, personne) : \zh{入} trait long + court ; \zh{人} traits symétriques. \textit{Sens : entrer, revenu \zh{收入} (shōurù).} \\ \hline
\zh{笔} & bǐ & \zh{竹} (zhú, bambou) & 10 & \zh{签} (qiān, signer) : même radical. \textit{Sens : stylo, trait, somme d'argent (\zh{一笔钱}).} \\ \hline
\zh{款} & kuǎn & \zh{欠} (qiàn, devoir) & 12 & \zh{宽} (kuān, large) : même phonétique \zh{宽}, radical \zh{欠} vs \zh{宀}. \textit{Sens : fonds, article, \zh{款项} (kuǎnxiàng).} \\ \hline
\zh{项} & xiàng & \zh{页} (yè, tête) & 9 & \zh{顺} (shùn,顺利) : même structure. \textit{Sens : article, poste de dépense \zh{项目} (xiàngmù).} \\ \hline
\zh{费} & fèi & \zh{贝} (bèi, coquillage/argent) & 12 & \zh{费} vs \zh{弗} (fú, non) : \zh{费} a \zh{贝} + \zh{弗}. \textit{Sens : frais, dépense \zh{学费} (xuéfèi), \zh{车费} (chēfèi).} \\ \hline
\end{tabularx}
\end{center}
\end{propriete}

\begin{experience}[Atelier écriture : « Mon carnet de budget »]
\begin{enumerate}
    \item Écrire chaque caractère du tableau ci-dessus 5 fois en respectant l'ordre des traits (utiliser papier quadrillé \zh{田字格}).
    \item Former 3 mots composés par caractère (ex: \zh{存} $\rightarrow$ \zh{存钱}, \zh{存在}, \zh{储存}).
    \item Dictée rapide : le professeur dicte le mot composé, l'élève écrit le caractère cible en grand au tableau.
\end{enumerate}
\end{experience}

\courssub{Lexique actif et collocations du thème}

\begin{center}
\begin{tabularx}{\linewidth}{|X|X|X|X|}\hline
\rowcolor{popblueL} \textbf{Caractères} & \textbf{Pinyin} & \textbf{Nature} & \textbf{Traduction / Collocations} \\ \hline
零花钱 & línghuáqián & n. & Argent de poche (\zh{给/花/存/省 零花钱}) \\ \hline
预算 & yùsuàn & n./v. & Budget ; établir un budget (\zh{做预算}, \zh{超预算}) \\ \hline
收入 & shōurù & n. & Revenu, rentrée d'argent (\zh{月收入}, \zh{收入来源}) \\ \hline
支出 & zhīchū & n. & Dépense, sortie d'argent (\zh{控制支出}, \zh{支出明细}) \\ \hline
结余 & jiéyú & n./v. & Solde, reste (\zh{月结余}, \zh{结余多少}) \\ \hline
存钱 / 存款 & cúnqián / cúnkuǎn & v. & Épargner, déposer en banque (\zh{存一千块}, \zh{定期存款}) \\ \hline
取钱 & qǔqián & v. & Retirer de l'argent (\zh{去银行取钱}) \\ \hline
转账 & zhuǎnzhàng & v. & Virement bancaire (\zh{微信转账}, \zh{支付宝转账}) \\ \hline
花费 & huāfèi & v./n. & Dépenser / Coût (\zh{花费时间}, \zh{花费不少钱}) \\ \hline
省钱 & shěngqián & v. & Faire des économies, coûter moins cher (\zh{买打折的省钱}) \\ \hline
乱花钱 & luàn huā qián & v.o. & Dépenser n'importe comment (\zh{不乱花钱}) \\ \hline
够 / 不够 & gòu / bù gòu & v. & Suffire / Ne pas suffire (\zh{钱不够买手机}) \\ \hline
差 & chà & v. & Manquer (quantité) (\zh{差五百块}, \zh{还差一点}) \\ \hline
努力 & nǔlì & adv./v. & Faire des efforts (\zh{努力存钱}, \zh{努力学习}) \\ \hline
继续 & jìxù & v. & Continuer (\zh{继续努力}, \zh{继续存钱}) \\ \hline
智能手机 & zhìnéng shǒujī & n. & Smartphone (Classif. \zh{部} bù / \zh{个} gè) \\ \hline
汉法词典 & Hàn-Fǎ cídiǎn & n. & Dictionnaire chinois-français (Classif. \zh{本} běn) \\ \hline
法郎 & fǎláng & n. & Franc (CFA) \\ \hline
元 / 块 & yuán / kuài & n. & Yuan (unité formelle / orale) \\ \hline
\end{tabularx}
\end{center}

\courssub{Thème : français → chinois}

\begin{methode}[Procédure standard pour le thème FR $\rightarrow$ ZH (4 étapes)]
\begin{enumerate}
    \item \textbf{Analyser la phrase française} : Identifier \cle{sujet}, \cle{verbe} (temps/aspect), \cle{objet}, \cle{compléments} (temps, lieu, manière, cause, but). Souligner \cle{connecteurs} (\emph{d'abord, parce que, donc, mais}).
    \item \textbf{Ordonner selon la syntaxe chinoise} : Appliquer \cle{Sujet + Temps + Lieu + Manière/But + Verbe + Objet}. Placer la cause (\zh{因为}) avant la conséquence (\zh{所以}).
    \item \textbf{Traduire syntagme par syntagme} : Choisir le bon vocabulaire (collocations : \zh{存钱} $\neq$ \zh{省钱}), le bon \cle{classifieur} (\zh{一本} dictionnaire, \zh{一部} téléphone, \zh{一块} argent), la bonne structure (\zh{太...了}, \zh{不乱...}, V+\zh{了}+Quantité).
    \item \textbf{Vérifier} : Caractères (traits, radicaux), Tons (changements \zh{bu}, \zh{yi}), Ponctuation pleine chasse (\zh{， 。}).
\end{enumerate}
\end{methode}

\begin{popfigure}
\begin{tikzpicture}[
    node distance=1.0cm and 1.2cm,
    box/.style={draw=popblue, thick, rounded corners, fill=popblueL, text width=3.0cm, align=center, font=\small},
    arrow/.style={-Stealth, thick, popdark}
]
\node[box, align=center] (a){\textbf{1. ANALYSER}\\\textit{Sujet, Verbe, Objet,}\\\textit{Compléments, Connecteurs}};
\node[box, right=of a, align=center] (b){\textbf{2. ORDONNER}\\\textit{Suj + Temps + Lieu +}\\\textit{Manière + V + Obj}\\\textit{因为...所以...}};
\node[box, right=of b, align=center] (c){\textbf{3. TRADUIRE}\\\textit{Lexique, Classifieurs,}\\\textit{Grammaire (了, 过, 着)}};
\node[box, right=of c, align=center] (d){\textbf{4. VÉRIFIER}\\\textit{Caractères, Tons,}\\\textit{Ponctuation, Pinyin}};
\draw[arrow] (a) -- (b);
\draw[arrow] (b) -- (c);
\draw[arrow] (c) -- (d);
\end{tikzpicture}
\legende{Procédure en 4 étapes pour réussir un thème français → chinois}
\end{popfigure}

\courssub{Analyse et traduction des phrases modèles du programme}

\begin{textebox}[Phrase 1 -- Réception de l'argent (Temps + Verbe + Objet)]
\zh{我每个星期得到零花钱。} \\
Wǒ měi gè xīngqī dédào línghuáqián. \\
« Je reçois de l'argent de poche chaque semaine. »
\end{textebox}

\begin{exemplebox}[Variantes : \zh{收到} (shōudào) « recevoir » (insiste sur la réception effective), \zh{给} (gěi) « donner » (voix active : \zh{爸爸给我零花钱}).]

\begin{textebox}[Phrase 2 -- Séquence d'actions (Connecteurs)]
\zh{首先，我买早饭，然后，我坐车。} \\
Shǒuxiān, wǒ mǎi zǎofàn, ránhòu, wǒ zuò chē. \\
« D'abord, j'achète le petit-déjeuner, ensuite je prends le car. »
\end{textebox}

\end{exemplebox}

\begin{exemplebox}[Collocations « Transport \& Repas »]
\begin{itemize}
    \item \zh{买早饭} (mǎi zǎofàn) : Acheter le petit-déjeuner / manger dehors le matin.
    \item \zh{坐车} (zuò chē) : Prendre un véhicule (bus, car, taxi, train) — générique.
    \item \zh{坐公交车} (zuò gōngjiāo chē) : Prendre le bus urbain.
    \item \zh{打车} (dǎ chē) : Prendre un taxi / VTC (DiDi).
    \item \zh{付车费} (fù chēfèi) : Payer le transport (plus formel).
\end{itemize}
\end{exemplebox}

\begin{textebox}[Phrase 3 -- Cause, But, Quantité, Classifieur]
\zh{因为我想买一本词典，所以我每个星期存一千法郎。} \\
Yīnwèi wǒ xiǎng mǎi yī běn cídiǎn, suǒyǐ wǒ měi gè xīngqī cún yīqiān fǎláng. \\
« Parce que je veux acheter un dictionnaire, j'économise mille francs par semaine. »
\end{textebox}

\begin{propriete}[Classifieurs (mesureurs) obligatoires]
Tout nombre + nom concret exige un classifieur.
\begin{itemize}
    \item \zh{一 \textbf{本} 词典} (yī \textbf{běn} cídiǎn) : \zh{本} pour livres, cahiers, dictionnaires.
    \item \zh{一部 智能手机} (yī \textbf{bù} zhìnéng shǒujī) : \zh{部} pour appareils, téléphones, machines.
    \item \zh{一 \textbf{块} 钱} (yī \textbf{kuài} qián) : \zh{块} (kuài) classifieur oral pour l'argent (\zh{元} yuán formel).
    \item \zh{一千 法郎} (yīqiān fǎláng) : \zh{千} (qiān) = nombre « mille ». Pas de classifieur entre le nombre et la devise \zh{法郎}, ou \zh{一千块法郎} (rare).
\end{itemize}
\end{propriete}

\begin{attention}[Piège n°2 : \zh{存} (cún) vs \zh{省} (shěng) -- « Économiser »]
\begin{itemize}
    \item \zh{\textbf{存}钱} (cún qián) : \textbf{Mettre de côté, déposer, épargner} (action positive d'accumulation). \rightarrow \emph{« J'économise 1000 F par semaine. »} = \zh{我每个星期存一千法郎。}
    \item \zh{\textbf{省}钱} (shěng qián) : \textbf{Faire des économies, dépenser moins, acheter moins cher} (réduction de dépense). \rightarrow \emph{« J'ai économisé 50 F en achetant d'occasion. »} = \zh{我买二手的，省了五十法郎。}
\end{itemize}
\emph{Erreur fréquente :} \zh{*我省一千法郎。} (Signifie « Je réduis mes dépenses de 1000F », pas « Je mets 1000F de côté »).
\end{attention}

\begin{textebox}[Phrase 4 -- Négation + Manière (\zh{乱})]
\zh{我不乱花钱。} \\
Wǒ bù luàn huā qián. \\
« Je ne dépense pas mon argent n'importe comment. »
\end{textebox}

\begin{propriete}[Adverbe de manière \zh{乱} (luàn) + Verbe : « N'importe comment / En désordre »]
\zh{乱} devant un verbe = manière désordonnée, irresponsable. \zh{不} + \zh{乱} + V = discipline, règle de conduite.
\medskip
\begin{center}
\begin{tabularx}{\linewidth}{|X|X|X|}\hline
\rowcolor{popgreenL} \textbf{Structure} & \textbf{Exemple chinois + Pinyin} & \textbf{Traduction} \\ \hline
\zh{不 + 乱 + V} & \zh{不乱说} (bù luàn shuō) & Ne pas parler n'importe comment / dire n'importe quoi \\ \hline
& \zh{不乱跑} (bù luàn pǎo) & Ne pas courir partout / rester sage \\ \hline
& \zh{\textbf{不乱花钱}} (bù luàn huā qián) & \textbf{Ne pas dépenser n'importe comment / Gérer son budget} \\ \hline
\zh{不 + 瞎 + V} & \zh{不瞎担心} (bù xiā dānxīn) & Ne pas s'inquiéter pour rien (oral) \\ \hline
\end{tabularx}
\end{center}
\end{propriete}

\begin{exemplebox}[Structure \zh{太}... \zh{le} : Degré excessif « Trop / Vraiment »]
\zh{坐车太贵了。} (Zuò chē tài guì le.) -- « Le transport coûte trop cher. »
\medskip
\noindent\textbf{Règle :} \zh{太} + Adjectif qualificatif + \zh{了} (particule de réalisation/changement d'état). Exprime un degré excessif subjectif.
\medskip
\begin{center}
\begin{tabularx}{\linewidth}{|X|X|X|}\hline
\rowcolor{poppinkL} \textbf{Français} & \textbf{Chinois} & \textbf{Pinyin} \\ \hline
C'est trop cher. & \zh{太贵了。} & Tài guì le. \\ \hline
C'est trop loin. & \zh{太远了。} & Tài yuǎn le. \\ \hline
Il fait trop chaud. & \zh{太热了。} & Tài rè le. \\ \hline
\end{tabularx}
\end{center}
\end{exemplebox}

\begin{propriete}[Particule \zh{了} (le) : Deux positions, deux fonctions]
\begin{enumerate}
    \item \textbf{Verbe + \zh{了} (Aspect parfaitif / Accompli) :} Action terminée, quantité atteinte.
    \zh{我存\textbf{了}一千法郎。} (Wǒ cún \textbf{le} yīqiān fǎláng.) -- « J'\textbf{ai} économisé 1000 F. »
    \item \textbf{Phrase + \zh{le} (Fin de phrase : Changement d'état / Situation nouvelle) :}
    \zh{我有钱\textbf{了}。} (Wǒ yǒu qián \textbf{le}.) -- « Maintenant j'ai de l'argent / Je ne suis plus sans le sou. »
    \zh{现在我有一万法郎\textbf{了}。} (Xiànzài wǒ yǒu yī wàn fǎláng \textbf{le}.) -- « Maintenant j'ai 10 000 F (alors qu'avant non). »
\end{enumerate}
\end{propriete}

\begin{attention}[Piège n°3 : \zh{了} $\neq$ Passé composé universel]
\begin{itemize}
    \item \textbf{Pas de \zh{了}} si action habituelle : \zh{我每个星期存钱} (pas de \zh{了}).
    \item \textbf{Pas de \zh{了}} si nié par \zh{没} : \zh{我没存钱} (pas de \zh{了}).
    \item \textbf{Position cruciale} : \zh{存了五百块} (Action accomplie, quantité) $\neq$ \zh{存五百块了} (Situation nouvelle : « J'en suis à 500 F économisés »).
\end{itemize}
\end{attention}

\begin{center}
\begin{tabularx}{\linewidth}{|X|X|X|X|}\hline
\rowcolor{popblueL} \textbf{Notion} & \textbf{Structure Chinoise} & \textbf{Exemple Modèle} & \textbf{Piège Francophone} \\ \hline
Temps (fréquence) & Suj + \zh{每个星期} + V + Obj & \zh{我每个星期存钱} & Placer le temps à la fin (*存钱每个星期) \\ \hline
Séquence & \zh{首先}... \zh{然后}... & \zh{首先买早饭，然后坐车} & Oublier la virgule après le connecteur \\ \hline
Cause $\rightarrow$ Conséquence & \zh{因为} Cause, \zh{所以} Conséquence & \zh{因为想买词典，所以存钱} & Inverser l'ordre (Conséquence \zh{因为} Cause) \\ \hline
Classifieur Livre & Num + \zh{本} + Nom & \zh{一本词典} & Omettre \zh{本} (*一词典) \\ \hline
Classifieur Téléphone & Num + \zh{部} + Nom & \zh{一部手机} & Utiliser \zh{个} ou \zh{本} (*一个手机/一本手机) \\ \hline
Économiser (mettre de côté) & \zh{存钱} / \zh{存 + Quantité} & \zh{存一千法郎} & Confondre avec \zh{省钱} (dépenser moins) \\ \hline
Négation manière & \zh{不乱} + V & \zh{不乱花钱} & Dire *不花钱乱 / *不随便花钱 (moins idiomatique) \\ \hline
Degré excessif & \zh{太} + Adj + \zh{了} & \zh{坐车太贵了} & Oublier \zh{了} final (*太贵) \\ \hline
Accompli (quantifié) & V + \zh{了} + Quantité & \zh{存了五百块} & Mettre \zh{了} en fin de phrase (*存五百块了 = changement d'état) \\ \hline
Manquer (quantité) & \zh{差} + Quantité & \zh{还差三千} & Dire *缺三千 (moins naturel pour argent) \\ \hline
\end{tabularx}
\end{center}

\courssub{Exercice résolu : Thème guidé}

\begin{exoresolu}[Thème : « Mon budget de mai »]
\textbf{Énoncé :} Traduire en chinois (caractères + pinyin) le paragraphe suivant :
\emph{« Je m'appelle Marie. Je suis en Terminale au lycée de Buea. Chaque mois, je reçois dix mille francs. Je veux acheter un smartphone, donc j'économise trois mille francs par mois. Je ne dépense pas l'argent n'importe comment. Le mois dernier, j'ai économisé trois mille francs. Maintenant j'ai six mille francs. C'est encore trop peu. »}

\tcblower

\textbf{Solution détaillée (étape par étape) :}

\begin{enumerate}
    \item \textbf{Phrase 1 :} « Je m'appelle Marie. Je suis en Terminale au lycée de Buea. »
    \begin{itemize}
        \item \zh{我叫玛丽。} (Wǒ jiào Mǎlì.)
        \item \zh{我是布埃亚高中的终班学生。} (Wǒ shì Bù'āiyà gāozhōng de zhōngbān xuéshēng.) -- \textit{Structure : Lieu + Nom + \zh{的} + Fonction. Buea = \zh{布埃亚}.}
    \end{itemize}
    \item \textbf{Phrase 2 :} « Chaque mois, je reçois dix mille francs. »
    \begin{itemize}
        \item Temps : \zh{每个月} (měi gè yuè) -- avant le verbe.
        \item Verbe : \zh{收到} (shōudào) « recevoir » (plus précis que \zh{得到} pour argent/courrier).
        \item \zh{我每个月收到一万法郎（零花钱）。} (Wǒ měi gè yuè shōudào yī wàn fǎláng (línghuáqián).)
    \end{itemize}
    \item \textbf{Phrase 3 :} « Je veux acheter un smartphone, donc j'économise trois mille francs par mois. »
    \begin{itemize}
        \item Cause/Conséquence : \zh{因为...所以...}.
        \item Smartphone : \zh{智能手机} (zhìnéng shǒujī). Classifieur \zh{部} (bù).
        \item Économiser (mettre de côté) : \zh{存钱}. Temps : \zh{每个月}.
        \item \zh{因为我想买一部智能手机，所以我每个月存三千法郎。} (Yīnwèi wǒ xiǎng mǎi yī bù zhìnéng shǒujī, suǒyǐ wǒ měi gè yuè cún sānqiān fǎláng.)
    \end{itemize}
    \item \textbf{Phrase 4 :} « Je ne dépense pas l'argent n'importe comment. »
    \begin{itemize}
        \item Structure \zh{不乱 + V + O}.
        \item \zh{我不乱花钱。} (Wǒ bù luàn huā qián.)
    \end{itemize}
    \item \textbf{Phrase 5 :} « Le mois dernier, j'ai économisé trois mille francs. »
    \begin{itemize}
        \item Temps : \zh{上个月} (shàng gè yuè).
        \item Accompli + Quantité : V + \zh{了} + Quantité.
        \item \zh{上个月我存了三千法郎。} (Shàng gè yuè wǒ cún le sānqiān fǎláng.)
    \end{itemize}
    \item \textbf{Phrase 6 :} « Maintenant j'ai six mille francs. »
    \begin{itemize}
        \item Changement d'état / Situation actuelle : \zh{现在} + Suj + V + Obj + \zh{了}.
        \item \zh{现在我有六千法郎了。} (Xiànzài wǒ yǒu liùqiān fǎláng le.)
    \end{itemize}
    \item \textbf{Phrase 7 :} « C'est encore trop peu. »
    \begin{itemize}
        \item « Trop peu » : \zh{太少了} (tài shǎo le). « Encore » : \zh{还} (hái) / \zh{还是} (háishi).
        \item \zh{还是太少了。} (Háishi tài shǎo le.) ou \zh{还太少了。}
    \end{itemize}
\end{enumerate}

\textbf{Texte final assemblé :}
\zh{我叫玛丽。我是布埃亚高中的终班学生。每个月我收到一万法郎零花钱。因为我想买一部智能手机，所以我每个月存三千法郎。我不乱花钱。上个月我存了三千法郎。现在我有六千法郎了，还是太少了。}
\end{exoresolu}

\courssub{Exercices d'application : Thème (Français $\rightarrow$ Chinois)}

\begin{exercice}[Entraînement au thème -- Thème Budget]
Traduire les phrases suivantes en chinois (caractères simplifiés + pinyin avec tons). Respecter l'ordre des mots, les classifieurs, la négation et l'aspect.
\begin{enumerate}
    \item Ma sœur aînée me donne 500 francs chaque jour.
    \item D'abord je paie le transport, ensuite j'achète à manger.
    \item Parce que le dictionnaire est trop cher, je ne l'achète pas maintenant.
    \item J'ai déjà économisé 5000 francs le mois dernier.
    \item Ce portable est trop cher, je n'ai pas assez d'argent.
\end{enumerate}
\end{exercice}

\begin{corrige}[Exercice Thème Budget]
\begin{enumerate}
    \item \zh{我姐姐每天给我五百法郎。} (Wǒ jiějie měitiān gěi wǒ wǔbǎi fǎláng.) -- \textit{Suj + Temps + V + Obj indirect + Obj direct.}
    \item \zh{首先我付车费，然后买东西吃。} (Shǒuxiān wǒ fù chēfèi, ránhòu mǎi dōngxi chī.) -- \textit{\zh{付车费} (payer le transport), \zh{买东西吃} (acheter qqch à manger). Sujet \zh{我} omis 2e fois.}
    \item \zh{因为词典太贵了，所以我现在不买了。} (Yīnwèi cídiǎn tài guì le, suǒyǐ wǒ xiànzài bù mǎi le.) -- \textit{\zh{不买了} = ne plus acheter (changement d'état). Attention : \zh{不买} = ne pas acheter (habitude/futur), \zh{没买} = n'a pas acheté (passé).}
    \item \zh{上个月我已经存了五千法郎。} (Shàng gè yuè wǒ yǐjīng cún le wǔqiān fǎláng.) -- \textit{Structure \zh{已经...了} pour action passée accomplie avec résultat actuel.}
    \item \zh{这个手机太贵了，我钱不够。} (Zhège shǒujī tài guì le, wǒ qián bù gòu.) -- \textit{\zh{不够} (bù gòu) = ne pas suffire. Alternative : \zh{我没钱买。} (Je n'ai pas l'argent pour l'acheter).}
\end{enumerate}
\end{corrige}

\courssub{Version : chinois → français et grille d'évaluation}

\begin{exercice}[Version : « Le budget de Wang Ming »]
Traduire le texte suivant en français. Respecter les temps, la logique des connecteurs et le vocabulaire budgétaire.
\begin{textebox}[Texte à traduire]
\zh{王明是北京的一名高中生。他的父母每个月给他三百元零花钱。因为他想买一辆二手自行车，所以他每个月存一百元。他不乱买饮料，也不乱买衣服。上个学期，他存了一千五百元。现在他有两千元了。离买自行车还差五百元。所以他下学期要打工赚钱。} \\
Wáng Míng shì Běijīng de yī míng gāozhōngshēng. Tā de fùmǔ měi gè yuè gěi tā sānbǎi yuán línghuáqián. Yīnwèi tā xiǎng mǎi yī liàng èrshǒu zìxíngchē, suǒyǐ tā měi gè yuè cún yībǎi yuán. Tā bù luàn mǎi yǐnliào, yě bù luàn mǎi yīfu. Shàng gè xuéqī, tā cún le yīqiān wǔbǎi yuán. Xiànzài tā yǒu liǎngqiān yuán le. Lí mǎi zìxíngchē hái chà wǔbǎi yuán. Suǒyǐ tā xià xuéqī yào dǎgōng zhuànqián.
\end{textebox}
\end{exercice}

\begin{corrige}[Exercice Version -- Grille d'évaluation (20 points)]
\textbf{Traduction de référence :}
« Wang Ming est un lycéen de Pékin. Ses parents lui donnent 300 yuans d'argent de poche chaque mois. Parce qu'il veut acheter un vélo d'occasion, il économise 100 yuans par mois. Il n'achète pas de boissons n'importe comment, ni de vêtements n'importe comment. Le semestre dernier, il a économisé 1500 yuans. Maintenant il a 2000 yuans. Il lui manque encore 500 yuans pour le vélo. Donc le semestre prochain, il doit faire un job pour gagner de l'argent. »

\medskip
\noindent\textbf{Grille d'évaluation détaillée :}
\begin{center}
\begin{tabularx}{\linewidth}{|X|X|X|}\hline
\rowcolor{poporangeL} \textbf{Critère} & \textbf{Points} & \textbf{Indicateurs de réussite} \\ \hline
Vocabulaire clé & 5 & \zh{高中生} (lycéen), \zh{零花钱}, \zh{二手自行车} (vélo d'occasion), \zh{存} (économiser), \zh{不乱买} (n'achète pas n'importe comment), \zh{学期} (semestre), \zh{打工} (job/boulot), \zh{赚钱} (gagner de l'argent), \zh{离...还差...} (il manque... pour...). \\ \hline
Grammaire / Structures & 6 & \zh{因为...所以...} (Cause/conséquence), \zh{每个月...存...} (Temps avant V), \zh{存了} (V+了 accompli), \zh{有...了} (Changement d'état), \zh{还差} (Manque encore), \zh{要...} (Obligation future). \\ \hline
Syntaxe / Ordre des mots & 4 & Respect de l'ordre Suj-Temps-Lieu-V-Obj en français. Gestion de la négation \zh{不乱}. \\ \hline
Connecteurs / Cohérence & 3 & \zh{因为/所以}, \zh{上个学期/现在/下学期}, \zh{所以} (conséquence finale). \\ \hline
Qualité du français & 2 & Français naturel, correct, registres adaptés (pas de calque « il économise 100 yuans \emph{chaque mois} » mais « \emph{chaque mois}, il économise... »). \\ \hline
\textbf{Total} & \textbf{20} & \\ \hline
\end{tabularx}
\end{center}

\emph{Note :} Pénaliser 0,5 pt par erreur de vocabulaire majeure, 1 pt par erreur de structure grammaticale (ordre des mots, \zh{了}, \zh{因为...所以...}), 0,5 pt par maladresse de français.

\end{corrige}

\courssub{Expression orale et débat}

\begin{experience}[Débat guidé : « Argent de poche : économiser ou profiter ? »]
\textbf{Contexte :} Deux élèves discutent de leur gestion budgétaire.
\textbf{Rôles :}
\begin{itemize}
    \item \textbf{Élève A (La Fourmi)} : Épargne rigoureuse, objectif précis (ex: dictionnaire, téléphone), note ses dépenses.
    \item \textbf{Élève B (La Cigale)} : Dépense pour le plaisir (snacks, jeux, sorties), vit au jour le jour, « l'argent est fait pour être dépensé ».
\end{itemize}
\textbf{Structures imposées à utiliser :}
\begin{itemize}
    \item \zh{我认为...因为...所以...} (Je pense que... parce que... donc...)
    \item \zh{每个星期/月我都...} (Chaque semaine/mois je...)
    \item \zh{不乱花钱 / 存钱 / 省钱}
    \item \zh{太...了} (C'est trop...), \zh{不够} (Pas assez), \zh{差...} (Il manque...)
    \item \zh{以后 / 下个学期 我要...} (À l'avenir / Le semestre prochain je vais...)
\end{itemize}
\textbf{Déroulement :} 3 min préparation (notes en pinyin/caractères) $\rightarrow$ 3 min jeu de rôle $\rightarrow$ 2 min feedback classe (correction tons, structures).
\end{experience}

\begin{savaistu}[Le Franc CFA et le Yuan : Repères pour le vocabulaire budgétaire]
Au Cameroun, la monnaie est le \textbf{Franc CFA} (XAF). En chinois : \zh{中非法郎} (Zhōngfēi Fǎláng) ou simplement \zh{法郎} (Fǎláng) dans le contexte local.
En Chine, la monnaie est le \textbf{Yuan} (CNY), unité courante : \zh{元} (yuán) / \zh{块} (kuài - oral). Sous-unités : \zh{角} (jiǎo) / \zh{毛} (máo - oral) = 1/10 yuan ; \zh{分} (fēn) = 1/100 yuan.
\medskip
\noindent\textbf{Repères de conversion approximatifs (variables) :} 1 Yuan $\approx$ 80--85 FCFA. 1000 FCFA $\approx$ 12 Yuan.
\medskip
\noindent\textbf{Exprimer un prix :}
\begin{itemize}
    \item 5000 FCFA : \zh{五千法郎} (Wǔqiān fǎláng).
    \item 50 Yuan : \zh{五十元} (Wǔshí yuán) / \zh{五十块钱} (Wǔshí kuài qián) / \zh{五十块} (Wǔshí kuài).
\end{itemize}
Connaître ces termes permet de situer le budget d'un lycéen camerounais (\zh{零花钱} souvent 2000--5000 F/semaine) par rapport à un lycéen chinois (\zh{零花钱} souvent 100--300 Yuan/semaine).
\end{savaistu}

\coursec{Version : chinois → français et grille d'évaluation}

Dans la section précédente, nous avons appris à traduire du français vers le chinois (le \cle{thème}) : il fallait penser en chinois, respecter l'ordre des mots chinois et tracer correctement les caractères. Nous faisons maintenant le chemin inverse : la \cle{version}, c'est-à-dire la traduction du chinois vers le français. L'exercice semble plus facile — on écrit dans sa langue maternelle — mais il cache ses propres pièges : calquer l'ordre chinois, oublier les nuances des négations \zh{不} et \zh{没}, ou traduire mot à mot au lieu de restituer le sens. Cette section vous donne une méthode complète, une traduction guidée phrase par phrase, puis la grille d'évaluation qui servira à noter vos productions écrites en classe et au Bac blanc.

\courssub{Méthode de la version : du sens vers la langue}

\begin{methode}[Réussir une version chinois → français en 5 étapes]
\begin{enumerate}
\item \textbf{Lecture globale silencieuse.} Lisez tout le texte une première fois sans rien écrire. Identifiez le thème (ici : l'argent de poche), le locuteur (\zh{我} ou un tiers ?) et le ton (récit, conseil, opinion ?).
\item \textbf{Repérage des connecteurs.} Soulignez les mots de liaison : \zh{但是} (mais), \zh{所以} (donc), \zh{因为} (parce que), \zh{我觉得} (je pense que). Ils donnent le squelette logique du texte : opposition → conséquence → avis personnel.
\item \textbf{Découpage en segments de sens.} Traduisez par groupes de mots, jamais caractère par caractère. \zh{零花钱} se traduit d'un bloc par « argent de poche », pas par « zéro dépenser argent ».
\item \textbf{Rédaction en français naturel.} Reformulez en français correct : ajoutez les articles (le, un, des), conjuguez les verbes, placez les compléments à la française. L'ordre chinois \zh{他常常乱花钱} devient « Il dépense souvent n'importe comment », jamais « Il souvent désordonné dépense argent ».
\item \textbf{Vérification finale.} Relisez en vous posant quatre questions : le sens global est-il fidèle ? Les temps sont-ils cohérents ? Les négations \zh{不}/\zh{没} sont-elles bien distinguées ? Les modaux (\zh{应该} = devrait) sont-ils rendus au conditionnel ?
\end{enumerate}
\end{methode}

\begin{attention}[Le piège du mot à mot]
Le chinois n'a ni articles, ni conjugaison, et son ordre des compléments diffère du français. Une version qui « colle » au chinois est une version ratée. Exemple : \zh{我的同学每个月得到一万法郎的零花钱} se traduit « Mon camarade reçoit dix mille francs d'argent de poche chaque mois » — et non « Mon camarade chaque mois obtient dix mille francs de argent de poche ». Le français exige des articles et une place naturelle du complément de temps.
\end{attention}

\courssub{Texte de version et observation}

\begin{textebox}[Texte de version — L'argent de poche de mon camarade]
我的同学每个月得到一万法郎的零花钱。但是他常常乱花钱，所以没有钱买学习用品。我觉得他应该做一个预算。\\
Wǒ de tóngxué měi ge yuè dédào yí wàn fǎláng de línghuāqián. Dànshì tā chángcháng luàn huā qián, suǒyǐ méiyǒu qián mǎi xuéxí yòngpǐn. Wǒ juéde tā yīnggāi zuò yí ge yùsuàn.\\
« Mon camarade de classe reçoit dix mille francs d'argent de poche chaque mois. Mais il dépense souvent n'importe comment, donc il n'a pas d'argent pour acheter des fournitures scolaires. Je pense qu'il devrait établir un budget. »
\end{textebox}

Observons d'abord la phrase clé du texte, celle qui porte le conseil final :

\begin{textebox}[La phrase du conseil]
他应该做一个预算。\\
Tā yīnggāi zuò yí ge yùsuàn.\\
« Il devrait établir un budget. » (\zh{应该} = devoir, exprime le conseil)
\end{textebox}

\begin{propriete}[Le modal 应该 (yīnggāi) : « devrait », le conseil]
Structure : \textbf{Sujet + 应该 + verbe}. \zh{应该} exprime un devoir moral ou un conseil, comme le conditionnel français « devrait ». En version, ne le traduisez jamais par le futur (« il fera ») ni par l'impératif : le chinois dit « il devrait », nuance de conseil, pas d'ordre. La négation se place \textbf{avant} le modal : \zh{不应该} « ne devrait pas ».
\begin{itemize}
\item \zh{我们应该节约用钱。} \emph{Wǒmen yīnggāi jiéyuē yòng qián.} « Nous devrions dépenser avec économie. »
\item \zh{你不应该乱花钱。} \emph{Nǐ bù yīnggāi luàn huā qián.} « Tu ne devrais pas dépenser n'importe comment. »
\end{itemize}
\end{propriete}

\courssub{Traduction guidée phrase par phrase}

\textbf{Phrase 1 :} \zh{我的同学每个月得到一万法郎的零花钱。}

Découpage : \zh{我的同学} (mon camarade) + \zh{每个月} (chaque mois) + \zh{得到} (recevoir/obtenir) + \zh{一万法郎的零花钱} (dix mille francs d'argent de poche). Choix de traduction : \zh{得到} se rend naturellement par « reçoit » (plus idiomatique que « obtient » pour de l'argent de poche). Le complément de temps \zh{每个月}, placé avant le verbe en chinois, passe en fin de phrase en français : « Mon camarade reçoit dix mille francs d'argent de poche chaque mois. » Notez que \zh{一万} = 10 000 : le chinois compte en \zh{万} (dizaines de mille), là où le français compte en milliers. Attention au changement de ton de \zh{一} : devant un quatrième ton (\zh{万}), il se prononce \emph{yí} — on écrit donc \emph{yí wàn}.

\textbf{Phrase 2 :} \zh{但是他常常乱花钱，所以没有钱买学习用品。}

Deux connecteurs : \zh{但是} (mais, opposition) et \zh{所以} (donc, conséquence). \zh{乱花钱} est une collocation figée : « dépenser n'importe comment / gaspiller ». \zh{没有钱} : la négation \zh{没} + \zh{有} signifie « ne pas avoir » ; en français courant, on dira « il n'a pas d'argent » ou « il n'a plus d'argent ». Traduction : « Mais il dépense souvent n'importe comment, donc il n'a pas d'argent pour acheter des fournitures scolaires. » Le verbe \zh{买} (acheter) introduit un complément d'objet que le français relie par « pour ».

\textbf{Phrase 3 :} \zh{我觉得他应该做一个预算。}

\zh{我觉得} introduit une opinion : « je pense que / je trouve que » ; le français exige la subordonnée avec « que ». Le classificateur \zh{个} devant \zh{预算} disparaît en français (« un budget » suffit). Traduction : « Je pense qu'il devrait établir un budget. »

\begin{exemplebox}[Exemples traduits autour de 乱花钱]
\zh{他常常乱花钱。} \emph{Tā chángcháng luàn huā qián.} « Il dépense souvent n'importe comment. »\\
\zh{弟弟乱花钱，妈妈很生气。} \emph{Dìdi luàn huā qián, māma hěn shēngqì.} « Mon petit frère gaspille son argent, maman est très en colère. »\\
\zh{不要乱花钱，要存钱！} \emph{Bú yào luàn huā qián, yào cún qián !} « Ne dépense pas n'importe comment, économise ! »
\end{exemplebox}

\courssub{Vérification : les quatre points de contrôle}

\begin{attention}[没 (méi) contre 不 (bù) : la négation juste]
\zh{没} nie les actions passées ou accomplies et sert obligatoirement devant \zh{有} ; \zh{不} nie le présent, l'habitude ou la volonté. Comparez :
\begin{itemize}
\item \zh{我没有钱。} \emph{Wǒ méiyǒu qián.} « Je n'ai pas d'argent. » (état, possession → toujours \zh{没})
\item \zh{我不花钱。} \emph{Wǒ bù huā qián.} « Je ne dépense pas (d'argent). » (habitude, refus → \zh{不})
\item \zh{他昨天没买学习用品。} \emph{Tā zuótiān méi mǎi xuéxí yòngpǐn.} « Hier, il n'a pas acheté de fournitures scolaires. » (passé → \zh{没})
\end{itemize}
Confondre les deux en version fait perdre la nuance temporelle : « il n'a pas d'argent » (constat) n'est pas « il ne dépense pas » (comportement).
\end{attention}

\begin{center}
\begin{tabularx}{\linewidth}{|X|X|X|X|}
\hline
\rowcolor{popblueL} Caractères & Pinyin & Nature & Traduction \\
\hline
零花钱 & línghuāqián & nom & argent de poche \\
\hline
预算 & yùsuàn & nom & budget \\
\hline
乱花钱 & luàn huā qián & locution verbale & dépenser n'importe comment \\
\hline
存钱 & cún qián & verbe + objet & économiser, épargner \\
\hline
学习用品 & xuéxí yòngpǐn & nom & fournitures scolaires \\
\hline
得到 & dédào & verbe & recevoir, obtenir \\
\hline
应该 & yīnggāi & modal & devoir (conseil) \\
\hline
节约 & jiéyuē & verbe / adjectif & économiser ; économe \\
\hline
\end{tabularx}
\end{center}

\courssub{Exercice résolu de version}

\begin{exoresolu}[Version guidée]
Traduisez en français naturel : \zh{我每个月有五千法郎的零花钱。我不乱花钱，因为我想存钱买一本书。}
\tcblower
\textbf{Étape 1 — Lecture globale :} le locuteur (\zh{我}) parle de son argent de poche et de son projet d'achat. Connecteurs : \zh{不} (négation d'habitude), \zh{因为} (parce que).\\
\textbf{Étape 2 — Segments :} \zh{每个月} chaque mois ; \zh{五千法郎} cinq mille francs ; \zh{不乱花钱} ne dépense pas n'importe comment ; \zh{存钱买一本书} économiser pour acheter un livre.\\
\textbf{Étape 3 — Traduction :} « Chaque mois, je reçois cinq mille francs d'argent de poche. Je ne dépense pas n'importe comment, parce que je veux économiser pour m'acheter un livre. »\\
\textbf{Vérification :} \zh{不} devant \zh{乱花钱} = habitude au présent (« je ne dépense pas ») ; \zh{想} = « vouloir » (je veux) ; le classificateur \zh{本} (livres) devient simplement « un livre ».
\end{exoresolu}

\courssub{Grille d'évaluation de la production écrite}

En classe comme au Bac blanc, votre production écrite sur l'argent de poche est notée sur 20 points selon cinq critères de 4 points chacun. Connaître cette grille, c'est savoir exactement où gagner des points.

\begin{popfigure}
\begin{tikzpicture}[
  crit/.style={draw=popblue, fill=popblueL, rounded corners=2pt, align=left, text width=3.4cm, minimum height=1.15cm, font=\small\bfseries, anchor=west},
  desc/.style={draw=popline, fill=popcream, rounded corners=2pt, align=left, text width=7.6cm, minimum height=1.15cm, font=\small, anchor=west},
  pts/.style={draw=poporange, fill=poporangeL, rounded corners=2pt, align=center, text width=1.2cm, minimum height=1.15cm, font=\small\bfseries, anchor=west}
]
\node[crit] at (0,0) {1. Exactitude des caractères};
\node[desc] at (3.7,0) {Traits corrects, radicaux bien placés ; aucun caractère confondu (买/卖, 没/沿)};
\node[pts] at (11.5,0) {/ 4};
\node[crit] at (0,-1.4) {2. Grammaire et ordre des mots};
\node[desc] at (3.7,-1.4) {Sujet + verbe + objet ; compléments de temps avant le verbe ; négations 不/没 bien choisies};
\node[pts] at (11.5,-1.4) {/ 4};
\node[crit] at (0,-2.8) {3. Connecteurs et structure};
\node[desc] at (3.7,-2.8) {Au moins deux connecteurs (但是, 所以, 因为, 我觉得) ; texte organisé en début--développement--conclusion};
\node[pts] at (11.5,-2.8) {/ 4};
\node[crit] at (0,-4.2) {4. Vocabulaire du thème};
\node[desc] at (3.7,-4.2) {Lexique du budget employé à bon escient : 零花钱, 预算, 存钱, 节约, 乱花钱};
\node[pts] at (11.5,-4.2) {/ 4};
\node[crit] at (0,-5.6) {5. Respect du sujet et longueur};
\node[desc] at (3.7,-5.6) {Le texte répond à la consigne ; longueur demandée respectée (environ 5 à 8 phrases)};
\node[pts] at (11.5,-5.6) {/ 4};
\end{tikzpicture}
\legende{Grille d'évaluation de la production écrite : 5 critères de 4 points, total sur 20.}
\end{popfigure}

\begin{aretenir}[Où se perdent les points ?]
La plupart des copies perdent des points sur les critères 1 et 2 : caractères mal formés et ordre des mots calqué du français. Un texte simple mais correct (critères 1, 2, 5) vaut déjà 12/20. Les connecteurs et le lexique riche (critères 3 et 4) font la différence entre une copie moyenne et une excellente copie.
\end{aretenir}

\begin{savaistu}[Caractères mal tracés ou caractères confondus ?]
Une copie avec des caractères mal tracés mais lisibles perd moins de points qu'une copie avec des caractères confondus qui changent le sens. Écrire \zh{卖} (mài, vendre) au lieu de \zh{买} (mǎi, acheter) inverse complètement la phrase : « il n'a pas d'argent pour vendre des fournitures » n'a plus aucun sens ! La règle d'or : mieux vaut un caractère simple et sûr qu'un caractère savant et faux. La différence entre \zh{买} et \zh{卖} ? Le petit élément supérieur \zh{十} de \zh{卖} : on « vend » ce qu'on met en plus sur l'étal.
\end{savaistu}

\courssub{Auto-évaluation et évaluation croisée entre pairs}

\begin{experience}[Évaluation croisée avec la grille]
\begin{enumerate}
\item \textbf{Production (15 min) :} chaque élève rédige un texte de 5 à 8 phrases sur le sujet : « Comment je gère mon argent de poche » (en utilisant au moins \zh{零花钱}, \zh{预算}, \zh{存钱} et deux connecteurs).
\item \textbf{Auto-évaluation (5 min) :} chacun se note sur 20 avec la grille, en justifiant chaque critère par une ligne de sa copie.
\item \textbf{Échange (10 min) :} les copies tournent par binôme. Le correcteur applique la grille, surligne en vert les réussites et entoure en rouge les caractères douteux et les négations \zh{不}/\zh{没} à vérifier.
\item \textbf{Discussion (5 min) :} le binôme compare auto-évaluation et évaluation croisée ; tout écart supérieur à 2 points sur un critère doit être discuté et justifié.
\item \textbf{Remédiation :} chaque élève réécrit les deux phrases les plus faibles de sa copie.
\end{enumerate}
\end{experience}

\begin{exercice}[À vous de jouer]
\textbf{A. Version.} Traduisez en français naturel : \zh{我的哥哥每个月得到两万法郎，但是他乱花钱，所以没有钱。我觉得他应该存钱。}\\
\textbf{B. Correction de copie.} Voici la traduction d'un élève : « Mon grand frère chaque mois obtient vingt mille francs, mais il dépense, donc pas d'argent. Je pense il doit économiser. » Relevez quatre défauts à l'aide de la grille et proposez une version améliorée.
\end{exercice}

\begin{corrige}[Exercice A et B]
\textbf{A.} « Mon grand frère reçoit vingt mille francs chaque mois, mais il dépense n'importe comment, donc il n'a plus d'argent. Je pense qu'il devrait économiser. » Points de vigilance : \zh{两万} = 20 000 (on dit \zh{两万}, jamais \zh{二万}) ; \zh{没(有)钱} = « n'a pas d'argent » (\zh{没} obligatoire devant \zh{有}) ; \zh{应该} = conditionnel « devrait ».\\
\textbf{B.} Défauts : (1) ordre calqué du chinois « chaque mois obtient » au lieu de « reçoit ... chaque mois » ; (2) \zh{乱花钱} réduit à « dépense » : la nuance « n'importe comment » est perdue (critère 4) ; (3) « donc pas d'argent » : phrase sans sujet, incorrecte en français (critère 2) ; (4) « je pense il doit » : il manque « que » et \zh{应该} vaut « devrait », pas « doit » (critère 2). Version améliorée : « Mon grand frère reçoit vingt mille francs chaque mois, mais il dépense n'importe comment, donc il n'a plus d'argent. Je pense qu'il devrait économiser. »
\end{corrige}

Vous disposez désormais de tous les outils de la leçon : un modèle de texte, des connecteurs, l'écriture des caractères du budget, le lexique actif, et les méthodes de thème et de version. Entraînez-vous régulièrement avec la grille d'évaluation : c'est en corrigeant — ses propres copies et celles des autres — qu'on apprend le plus vite à écrire juste.

\newpage

\coursec{Activité d'intégration}

\coursec{Leçon 27 — Expression écrite : gestion de l'argent de poche et caractères complexes du budget}

\courssub{Objectifs de la leçon}

À l'issue de cette leçon, l'élève sera capable de :
\begin{itemize}
\item rédiger un texte structuré de 8 à 10 phrases en chinois simplifié sur la gestion de son argent de poche (\zh{我的零花钱预算}) ;
\item utiliser correctement les connecteurs logiques \zh{首先}、\zh{然后}、\zh{因为……所以……} et la structure d'opinion \zh{我觉得……} ;
\item écrire sans faute les caractères complexes du thème budgétaire : \zh{零}、\zh{钱}、\zh{预}、\zh{算}、\zh{节}、\zh{省}、\zh{花}、\zh{费} (traits, radicaux, caractères proches) ;
\item présenter un mini-tableau budgétaire (reçu / dépensé / économisé) en chiffres et en caractères chinois ;
\item réaliser un thème (français $\to$ chinois) et une version (chinois $\to$ français) sur le vocabulaire financier quotidien ;
\item corriger la production d'un pair à l'aide d'une grille d'évaluation sur 20 points et lire son texte à voix haute avec des tons corrects.
\end{itemize}

\courssub{Situation de communication}

\begin{textebox}[Contexte : Semaine « Jeunesse et argent » au lycée de Yaoundé]
Tu es élève en Terminale A au lycée de Yaoundé. Chaque semaine, tes parents te donnent de l'argent de poche pour tes dépenses : transport (taxi/moto), goûter (beignets, jus), fournitures scolaires, crédit téléphonique. Comme beaucoup de jeunes Camerounais, tu dépenses parfois trop vite et tu n'arrives plus à économiser pour t'acheter le ballon de football dont tu rêves.

Le club de chinois organise une semaine thématique « Jeunesse et argent ». Le professeur demande à chaque élève de rédiger, en chinois, un petit texte intitulé \zh{我的零花钱预算}, accompagné d'un mini-tableau de budget, et de le présenter sous forme d'affiche. Les meilleures affiches seront exposées au mur de la salle de chinois et lues à voix haute devant la classe.
\end{textebox}

\begin{textebox}[Document de départ : Consignes du professeur]
同学们好！\\
\textit{Tóngxuémen hǎo!}\\
« Bonjour à tous ! »\\
这个星期的作业：写一篇短文，题目是《我的零花钱预算》。\\
\textit{Zhège xīngqī de zuòyè: xiě yì piān duǎnwén, tímù shì « Wǒ de línghuāqián yùsuàn ».}\\
« Devoir de la semaine : rédiger un court texte intitulé "Mon budget d'argent de poche". »\\
文章要有八到十句话，用"首先"、"然后"、"因为……所以……"和"我觉得"。\\
\textit{Wénzhāng yào yǒu bā dào shí jùhuà, yòng « shǒuxiān », « ránhòu », « yīnwèi... suǒyǐ... » hé « wǒ juéde ».}\\
« Le texte doit contenir 8 à 10 phrases, avec "d'abord", "ensuite", "parce que... donc..." et "je pense que". »\\
还要做一个小表：收到的钱、花的钱、省下的钱。\\
\textit{Hái yào zuò yí ge xiǎo biǎo: shōudào de qián, huā de qián, shěngxià de qián.}\\
« Il faut aussi faire un petit tableau : argent reçu, argent dépensé, argent économisé. »
\end{textebox}

\courssub{Modèle de texte commenté}

\begin{exoresolu}[Affiche modèle : 我的零花钱预算]
Texte de 9 phrases respectant le plan en 4 parties, avec connecteurs obligatoires et mini-tableau.
\tcblower
\textbf{Titre :} \zh{我的零花钱预算} \textit{(Wǒ de línghuāqián yùsuàn)}

\textbf{Texte :}
\begin{enumerate}
\item \zh{每个星期，妈妈给我五千块零花钱。} \textit{Měi ge xīngqī, māma gěi wǒ wǔqiān kuài línghuāqián.} « Chaque semaine, maman me donne 5000 F d'argent de poche. »
\item \zh{首先，我花一千块坐车去学校。} \textit{Shǒuxiān, wǒ huā yìqiān kuài zuò chē qù xuéxiào.} « D'abord, je dépense 1000 F pour aller à l'école en voiture. »
\item \zh{然后，我花两千块买吃的和本子。} \textit{Ránhòu, wǒ huā liǎngqiān kuài mǎi chī de hé běnzi.} « Ensuite, je dépense 2000 F pour acheter à manger et des cahiers. »
\item \zh{有时候，我还花五百块充手机话费。} \textit{Yǒushíhou, wǒ hái huā wǔbǎi kuài chōng shǒujī huàfèi.} « Parfois, je dépense encore 500 F pour recharger mon crédit téléphone. »
\item \zh{因为我想买一个足球，所以我每个星期节省一千块。} \textit{Yīnwèi wǒ xiǎng mǎi yí ge zúqiú, suǒyǐ wǒ měi ge xīngqī jiéshěng yìqiān kuài.} « Parce que je veux acheter un ballon, j'économise 1000 F chaque semaine. »
\item \zh{我觉得节省钱很有用。} \textit{Wǒ juéde jiéshěng qián hěn yǒuyòng.} « Je pense qu'économiser est très utile. »
\item \zh{我也觉得管理零花钱不难。} \textit{Wǒ yě juéde guǎnlǐ línghuāqián bù nán.} « Je trouve aussi que gérer son argent de poche n'est pas difficile. »
\item \zh{下个星期，我要继续节省。} \textit{Xià ge xīngqī, wǒ yào jìxù jiéshěng.} « La semaine prochaine, je vais continuer à économiser. »
\item \zh{希望我能早日买到足球。} \textit{Xīwàng wǒ néng zǎorì mǎidào zúqiú.} « J'espère pouvoir acheter le ballon le plus tôt possible. »
\end{enumerate}

\textbf{Mini-tableau de budget :}
\begin{center}
\begin{tabular}{|c|c|c|}
\hline
\rowcolor{popblueL} \zh{项目} & \zh{金额 (chiffres)} & \zh{金额 (caractères)} \\
\hline
\zh{收到的钱} & 5000 & \zh{五千块} \\
\hline
\zh{花的钱} & 3500 & \zh{三千五百块} \\
\hline
\zh{省下的钱} & 1500 & \zh{一千五百块} \\
\hline
\end{tabular}
\end{center}

\textbf{Commentaire pédagogique :} Le plan en 4 parties est respecté (Revenus $\to$ Dépenses \zh{首先/然后} $\to$ Épargne/But \zh{因为……所以……} $\to$ Opinion \zh{我觉得}). Les connecteurs obligatoires sont présents. Le classificateur \zh{块} (kuài) est utilisé pour les montants. L'aspect accompli \zh{了} n'est pas nécessaire ici car il s'agit d'habitudes (\zh{每个星期}), mais il serait requis pour un récit passé ponctuel (ex: \zh{上个星期我花了...}). Les caractères cibles (\zh{零、预、算、节、省、花、费}) apparaissent naturellement.
\end{exoresolu}

\courssub{Structure et connecteurs du texte}

\begin{propriete}[Connecteurs logiques obligatoires (Niveau HSK 3)]
\begin{itemize}
\item \textbf{Énumération séquentielle :} \zh{首先……，然后……} \textit{(shǒuxiān... ránhòu...)} \\
Structure : \zh{首先} + Phrase 1 \zh{，} \zh{然后} + Phrase 2. \\
Exemple : \zh{首先我吃饭，然后我去学校。} \textit{(D'abord je mange, ensuite je vais à l'école.)}
\item \textbf{Relation cause-conséquence :} \zh{因为……，所以……} \textit{(yīnwèi... suǒyǐ...)} \\
Structure : \zh{因为} + Cause \zh{，} \zh{所以} + Conséquence. \zh{所以} ne peut pas être omis. \\
Exemple : \zh{因为下雨，所以我不去了。} \textit{(Parce qu'il pleut, donc je n'y vais pas.)} \\
\emph{Attention :} En français on dit « Parce que..., je... » (pas de « donc » obligatoire). En chinois \zh{所以} est quasi obligatoire à l'écrit formel.
\item \textbf{Expression de l'opinion :} \zh{我觉得……} \textit{(wǒ juéde...)} \\
Structure : \zh{我觉得} + Proposition complète (Sujet + Verbe/Adjectif...). \\
\emph{Interdiction :} Ne pas mettre \zh{是} (shì) devant un adjectif après \zh{觉得}. \\
Correct : \zh{我觉得很重要。} \textit{(Je trouve que c'est important.)} \\
Incorrect : \zh{我觉得是很重要。}
\item \textbf{Marque d'habitude / généralité :} Pas de \zh{了} pour les actions habituelles (\zh{每个星期}). \zh{了} marque l'accompli / le changement d'état pour un événement ponctuel passé.
\end{itemize}
\end{propriete}

\begin{methode}[Comment structurer son texte « 我的零花钱预算 »]
\begin{enumerate}
\item \textbf{Introduction (1 phrase) :} Source + Fréquence + Montant. \zh{每个星期，爸爸/妈妈给我[X]块零花钱。}
\item \textbf{Dépenses (2-3 phrases) :} Utiliser \zh{首先} pour le poste principal (transport), \zh{然后} pour les autres (nourriture, fournitures, téléphone). Verbe \zh{花} + Montant + \zh{块} + Verbe d'action (\zh{坐车}、\zh{买}、\zh{充话费}).
\item \textbf{Épargne et Objectif (2 phrases) :} \zh{因为我想买[Objectif]，所以我每个星期节省[Montant]块。} + \zh{我觉得节省钱很重要/有用。}
\item \textbf{Conclusion / Projection (1-2 phrases) :} \zh{下个星期我要继续节省。} \zh{希望我能早日买到[Objectif].}
\item \textbf{Tableau :} 3 lignes : \zh{收到的钱}、\zh{花的钱}、\zh{省下的钱}. Montants en chiffres ET en caractères (\zh{五千}、\zh{三千五百}).
\end{enumerate}
\end{methode}

\courssub{Écriture correcte des caractères complexes du budget}

\begin{propriete}[Analyse graphémique : Traits, Radicaux, Caractères proches]
Le tableau ci-dessous détaille les 8 caractères cibles. Respectez l'ordre des traits : \textbf{haut $\to$ bas, gauche $\to$ droite, horizontal avant vertical, trait central avant les ailes, cadre avant contenu, fermeture en dernier}.
\end{propriete}

\begin{center}
\begin{tabularx}{\linewidth}{|X|X|X|X|X|}
\hline
\rowcolor{popblueL} Caractère & Pinyin & Radical (Clé) & Nb traits & Caractères proches / Pièges \\
\hline
\zh{零} & líng & \zh{雨} (yǔ, pluie) en haut & 13 & \zh{令} (lìng, ordre - sans pluie), \zh{领} (lǐng, cou - clé 页) \\
\hline
\zh{钱} & qián & \zh{钅} (jīn, métal/or) à gauche & 10 & \zh{线} (xiàn, fil - clé 纟), \zh{鉴} (jiàn, miroir - clé 金) \\
\hline
\zh{预} & yù & \zh{页} (yè, page/tête) à droite & 12 & \zh{序} (xù, ordre - même structure, clé 予 différent), \zh{预} vs \zh{豫} (yù, province Henan) \\
\hline
\zh{算} & suàn & \zh{竹} (zhú, bambou) en haut & 14 & \zh{管} (guǎn, tube/gérer - même clé, haut différent), \zh{箱} (xiāng, caisse) \\
\hline
\zh{节} & jié & \zh{竹} (zhú, bambou) en haut & 13 & \zh{即} (jí, aussitôt - sans bambou), \zh{迎} (yíng, accueillir) \\
\hline
\zh{省} & shěng/xǐng & \zh{目} (mù, œil) en bas & 9 & \zh{看} (kàn, regarder - main 手 au lieu de 少), \zh{眉} (méi, sourcil) \\
\hline
\zh{花} & huā & \zh{艹} (cǎo, herbe) en haut & 7 & \zh{华} (huá, splendeur/Chine - même composants, structure diff.), \zh{哗} (huā, bruit - clé 口) \\
\hline
\zh{费} & fèi & \zh{贝} (bèi, coquillage/argent) & 9 & \zh{费} vs \zh{弗} (fú, non - haut différent), \zh{贝} vs \zh{见} (jiàn, voir) \\
\hline
\end{tabularx}
\end{center}

\begin{experience}[Atelier d'écriture guidée (15 min)]
\begin{enumerate}
\item Sur papier quadrillé (Tianzige), écrire 5 fois chaque caractère en respectant l'ordre des traits affiché au tableau.
\item Exercice de discrimination : Le professeur dicte : \zh{零钱}、\zh{预算}、\zh{节省}、\zh{花费}、\zh{线}、\zh{看}、\zh{华}. L'élève écrit le bon caractère.
\item Recherche de radicaux : Classer les 8 caractères par clé : \zh{钅} (钱), \zh{雨} (零), \zh{页} (预), \zh{竹} (算、节), \zh{目} (省), \zh{艹} (花), \zh{贝} (费).
\end{enumerate}
\end{experience}

\courssub{Lexique actif et collocations du thème}

\begin{center}
\begin{tabularx}{\linewidth}{|X|X|X|X|}
\hline
\rowcolor{popblueL} Caractères & Pinyin & Nature & Traduction \\
\hline
零花钱 & línghuāqián & n. & argent de poche \\
\hline
预算 & yùsuàn & n. / v. & budget / établir un budget \\
\hline
收到 & shōudào & v. & recevoir (colis, argent, message) \\
\hline
收入 & shōurù & n. & revenus (plus formel que 收到) \\
\hline
花 & huā & v. & dépenser (argent, temps) \\
\hline
花费 & huāfèi & n. / v. & frais, dépenses / coûter \\
\hline
节省 / 省 & jiéshěng / shěng & v. & économiser (argent, temps, énergie) \\
\hline
存钱 & cúnqián & v. & mettre de l'argent de côté (banque/tirelire) \\
\hline
贵 / 便宜 & guì / piányi & adj. & cher / bon marché \\
\hline
需要 & xūyào & v. & avoir besoin de \\
\hline
想买 & xiǎng mǎi & v. & vouloir acheter \\
\hline
管理 & guǎnlǐ & v. & gérer, administrer \\
\hline
重要 & zhòngyào & adj. & important \\
\hline
块 & kuài & m. & classificateur pour l'argent (Yuan / F CFA familier) \\
\hline
充话费 & chōng huàfèi & v-o & recharger le crédit téléphonique \\
\hline
坐车 & zuò chē & v-o & prendre le transport (taxi, bus, moto) \\
\hline
\end{tabularx}
\end{center}

\begin{aretenir}[Collocations utiles pour l'expression écrite]
\begin{itemize}
\item \zh{制定预算} (zhìdìng yùsuàn) : établir un budget.
\item \zh{控制花费} (kòngzhì huāfèi) : contrôler ses dépenses.
\item \zh{开源节流} (kāiyuán jiéliú) \emph{(chengyu)} : augmenter les revenus et réduire les dépenses (vivre dans ses moyens).
\item \zh{月光族} (yuèguāngzú) \emph{(argot)} : « clan du mois nu » = personnes qui dépensent tout leur salaire avant la fin du mois.
\item \zh{存钱罐} (cúnqián guàn) : tirelire (lit. « pot à économies »).
\end{itemize}
\end{aretenir}

\courssub{Thème : Français $\to$ Chinois}

\begin{exercice}[Thème guidé : Mon budget hebdomadaire]
Traduire les phrases suivantes en chinois simplifié. Utilisez le vocabulaire de la leçon. Respectez l'ordre des mots (Sujet + Temps + Lieu + Verbe + Complément) et les classificateurs.
\begin{enumerate}
\item Chaque semaine, mon père me donne trois mille francs d'argent de poche.
\item D'abord, je dépense mille francs pour le transport. Ensuite, je dépense mille cinq cents francs pour manger.
\item Parce que je veux acheter un nouveau téléphone portable, je dois économiser de l'argent.
\item Je pense que gérer l'argent de poche est très important pour les lycéens.
\item Reçu : 3000 F. Dépensé : 2500 F. Économisé : 500 F.
\end{enumerate}
\end{exercice}

\begin{corrige}[Thème guidé]
\begin{enumerate}
\item \zh{每个星期，爸爸给我三千块零花钱。} \textit{(Měi ge xīngqī, bàba gěi wǒ sānqiān kuài línghuāqián.)}
\item \zh{首先，我花一千块坐车。然后，我花一千五百块买吃的。} \textit{(Shǒuxiān, wǒ huā yìqiān kuài zuò chē. Ránhòu, wǒ huā yìqiān wǔbǎi kuài mǎi chī de.)}
\item \zh{因为我想买一个新手机，所以我必须节省钱。} \textit{(Yīnwèi wǒ xiǎng mǎi yí ge xīn shǒujī, suǒyǐ wǒ bìxū jiéshěng qián.)} \emph{Note : \zh{必须} (bìxū) « devoir » niveau HSK 3, sinon \zh{要} (yào).}
\item \zh{我觉得管理零花钱对高中生很重要。} \textit{(Wǒ juéde guǎnlǐ línghuāqián duì gāozhōngshēng hěn zhòngyào.)}
\item \zh{收到的钱：三千块。花的钱：二千五百块。省下的钱：五百块。} \textit{(Shōudào de qián: sānqiān kuài. Huā de qián: èrqiān wǔbǎi kuài. Shěngxià de qián: wǔbǎi kuài.)}
\end{enumerate}
\end{corrige}

\courssub{Version : Chinois $\to$ Français et Grille d'évaluation}

\begin{exercice}[Version : Lettre d'un correspondant chinois]
Traduire le texte suivant en français naturel. Attention aux connecteurs et aux montants.
\begin{textebox}[Texte à traduire : 李明的预算]
我是李明，今年十八岁，在北京上高三。\\
\textit{Wǒ shì Lǐ Míng, jīnnián shíbā suì, zài Běijīng shàng gāosān.}\\
每个月，父母给我两千块零花钱。\\
\textit{Měi ge yuè, fùmǔ gěi wǒ liǎngqiān kuài línghuāqián.}\\
首先，我花八百块坐地铁去学校。\\
\textit{Shǒuxiān, wǒ huā bābǎi kuài zuò dìtiě qù xuéxiào.}\\
然后，我花六百块吃午饭和买文具。\\
\textit{Ránhòu, wǒ huā liùbǎi kuài chī wǔfàn hé mǎi wénjù.}\\
因为我想给妹妹买生日礼物，所以我每个月节省六百块。\\
\textit{Yīnwèi wǒ xiǎng gěi mèimei mǎi shēngrì lǐwù, suǒyǐ wǒ měi ge yuè jiéshěng liùbǎi kuài.}\\
我觉得存钱是一个好习惯。\\
\textit{Wǒ juéde cúnqián shì yí ge hǎo xíguàn.}
\end{textebox}
\end{exercice}

\begin{corrige}[Version : Li Ming's Budget]
« Je m'appelle Li Ming, j'ai 18 ans cette année, je suis en Terminale (Gāosān) à Pékin.
Chaque mois, mes parents me donnent 2000 Yuan d'argent de poche.
D'abord, je dépense 800 Yuan pour prendre le métro pour aller à l'école.
Ensuite, je dépense 600 Yuan pour le déjeuner et acheter de la papeterie.
Parce que je veux acheter un cadeau d'anniversaire pour ma petite sœur, j'économise 600 Yuan chaque mois.
Je pense qu'économiser est une bonne habitude. »

\textbf{Points de vigilance :}
\begin{itemize}
\item \zh{高三} = Terminale (dernière année du lycée). \zh{上} = suivre des cours / aller à l'école.
\item \zh{每个月} (chaque mois) vs \zh{每个星期} (chaque semaine). Unité : \zh{块} (Yuan en Chine).
\item \zh{坐地铁} (prendre le métro) : transport typique pékinois.
\item \zh{妹妹} (mèimei) : petite sœur. \zh{生日礼物} : cadeau d'anniversaire.
\item \zh{存钱} (cúnqián) : mettre en épargne (banque/tirelire), synonyme de \zh{节省} mais plus concret pour « mettre de côté ».
\item \zh{习惯} (xíguàn) : habitude. Structure \zh{是一个好习惯}.
\end{itemize}
\end{corrige}

\courssub{Grille d'évaluation de l'affiche « 我的零花钱预算 » (Sur 20 points)}

\begin{center}
\begin{tabularx}{\linewidth}{|X|c|X|}
\hline
\rowcolor{popblueL} \textbf{Critères} & \textbf{Points} & \textbf{Indicateurs de réussite} \\
\hline
Respect des consignes (titre, 8-10 phrases, tableau) & 3 & Titre exact, nombre de phrases dans la fourchette, tableau présent avec 3 lignes. \\
\hline
Structure et connecteurs (\zh{首先}、\zh{然后}、\zh{因为……所以……}、\zh{我觉得}) & 4 & Les 4 connecteurs utilisés correctement au moins une fois, ordre logique respecté. \\
\hline
Correction grammaticale (ordre mots, \zh{了} aspectuel, classificateur \zh{块}) & 5 & Phrases SVO correctes, \zh{块} systématique pour l'argent, \zh{了} utilisé si action passée ponctuelle (pas pour habitude). \\
\hline
Vocabulaire du thème (caractères cibles : 零钱预算节省花费收到省下) & 4 & Au moins 6 des 8 caractères cibles employés correctement dans le texte ou le tableau. \\
\hline
Exactitude de l'écriture des caractères (traits, radicaux) & 2 & Caractères lisibles, proportions respectées, pas de confusion (ex: \zh{钱}/\zh{线}, \zh{省}/\zh{看}). \\
\hline
Présentation de l'affiche (lisibilité, mise en page, dessin) & 2 & Texte recopié au propre, tableau aligné, titre visible, illustration pertinente. \\
\hline
Expression orale (lecture à voix haute : tons, fluidité) & \textbf{0 / Bonus} & \emph{Non noté sur 20 écrit, mais évalué à part pour l'oral. Bonus +1 si excellent.} \\
\hline
\textbf{TOTAL} & \textbf{20} &  \\
\hline
\end{tabularx}
\end{center}

\courssub{Bilan et ouverture socioculturelle}

\begin{savaistu}[Argent de poche : Cameroun vs Chine]
\begin{itemize}
\item \textbf{Fréquence :} Au Cameroun, l'argent de poche (\zh{零花钱}) est souvent \textbf{hebdomadaire} (semaine de cours). En Chine, il est plutôt \textbf{mensuel} (\zh{每个月}) pour les lycéens, parfois versé sur une carte bancaire ou WeChat/Alipay (\zh{微信/支付宝}).
\item \textbf{Montants :} Au Cameroun (Yaoundé/Douala), 3000--5000 F CFA/semaine est courant pour le transport + goûter. En Chine (Pékin/Shanghai), 500--2000 Yuan/mois selon les familles (1 Yuan $\approx$ 80 F CFA).
\item \textbf{Paiement mobile :} Les lycéens chinois utilisent massivement le QR code (\zh{扫码} sǎomǎ) pour payer le bus, la cantine, les achats. L'argent liquide (\zh{现金} xiànjīn) disparaît. Au Cameroun, le Mobile Money (Orange Money, MTN MoMo) se développe mais le cash reste roi pour les petits achats (beignets, taxi-moto).
\item \textbf{Éducation financière :} En Chine, des cours de « gestion financière pour jeunes » (\zh{理财教育} lǐcái jiàoyù) apparaissent dans les écoles. Au Cameroun, l'éducation financière est souvent informelle (conseils parents, tontines familiales).
\end{itemize}
\end{savaistu}

\begin{experience}[Débat oral : « Faut-il donner de l'argent de poche aux lycéens ? »]
\textbf{Consigne :} En groupes de 4. Préparez 3 arguments POUR et 3 arguments CONTRE en chinois (niveau HSK 3). Utilisez \zh{我觉得}、\zh{因为……所以……}、\zh{另外} (de plus).
\begin{itemize}
\item \textbf{Pour :} \zh{学会管理钱} (apprendre à gérer l'argent), \zh{买需要的东西} (acheter ce qui est nécessaire), \zh{培养责任感} (cultiver le sens des responsabilités).
\item \textbf{Contre :} \zh{乱花钱} (dépenser n'importe comment), \zh{买零食不好} (acheter des snacks c'est mauvais pour la santé), \zh{比较攀比} (comparer/concurrence entre camarades).
\end{itemize}
Un porte-parole par groupe présente la synthèse en 1 minute.
\end{experience}

\begin{attention}[Pièges de l'examen (Baccalauréat / Évaluation continue)]
\begin{itemize}
\item \textbf{Ordre des mots :} Ne pas écrire \zh{我花钱一千块} (SVO incorrect). Correct : \zh{我花一千块钱} (S + V + Quantité + Objet) ou \zh{我花了一千块钱} (avec \zh{了}).
\item \textbf{Classificateur \zh{块} vs \zh{元} :} \zh{块} (kuài) = oral / familier (utilisé au Cameroun pour les F CFA). \zh{元} (yuán) = écrit / formel / officiel (Chine). À l'écrit d'examen, préférez \zh{元} ou \zh{块钱}.
\item \textbf{\zh{省} vs \zh{看} :} \zh{省} (économiser) a \zh{少} (peu) en haut à droite. \zh{看} (regarder) a \zh{手} (main) en haut à droite. \emph{Mnémotechnique :} Pour économiser, on a « peu » (少) dans les yeux (目). Pour regarder, on met la main (手) sur les yeux (visière).
\item \textbf{\zh{预算} :} Ne pas confondre \zh{预} (yù, à l'avance) et \zh{序} (xù, ordre). \zh{预} a \zh{予} (donner) à droite ; \zh{序} a \zh{予} aussi mais structure différente. Rappel : \zh{预习} (réviser à l'avance), \zh{预报} (prévoir/annoncer).
\item \textbf{Traduction de « argent » :} \zh{钱} (qian) = argent, monnaie. \zh{零花钱} = argent de poche. Ne pas dire \zh{金钱} (jīnqián, argent/finances - trop formel) pour l'argent de poche.
\end{itemize}
\end{attention}

\newpage

\coursec{Méthodes et exercices résolus}

\coursec{Leçon 27 — Expression écrite : gestion de l'argent de poche et caractères complexes du budget}

\courssub{Méthode 1 — Rédiger un texte simple sur la gestion de l'argent de poche}

\begin{methode}[Produire un texte sur mon budget]
Pour rédiger un texte de 5 à 8 phrases sur la gestion de l'argent de poche, suis cette démarche :
\begin{enumerate}
\item \textbf{Présenter la situation} : qui reçoit de l'argent, combien, de qui. Structure : \zh{Sujet + 给 + Bénéficiaire + Objet} (donner quelque chose à quelqu'un). Exemple : \zh{我爸爸妈妈每个月给我五千法郎零花钱。} \emph{Wǒ bàba māma měi ge yuè gěi wǒ wǔ qiān fǎ láng línghuāqián.}
\item \textbf{Dire comment on dépense} : verbes \zh{花} (dépenser de l'argent), \zh{买} (acheter), \zh{用} (utiliser). Structure : \zh{花 + somme + (在) + objet} ou \zh{买 + objet}. Exemple : \zh{我花两千法郎买学习用品。}
\item \textbf{Dire comment on économise} : \zh{存} (mettre de côté, déposer), \zh{省} (économiser, éviter le gaspillage). Structures : \zh{存 + somme} ; \zh{把 + 钱 + 存起来} (construction \zh{把} pour marquer l'action sur l'objet).
\item \textbf{Donner une opinion} : \zh{我觉得……} (je pense que…), \zh{我认为……} (j'estime que…), \zh{……很重要/很有必要} (… est important/nécessaire).
\item \textbf{Conclure par un projet ou un conseil} : \zh{以后我想……} (plus tard, je voudrais…), \zh{我打算……} (je prévois de…), \zh{我们应该……} (nous devrions…).
\item \textbf{Relire et vérifier} : les tons (pinyin), l'ordre \textbf{Sujet – Temps – Verbe – Complément}, la place du complément de temps \textbf{avant le verbe}, les classificateurs (\zh{块钱}, \zh{本子}, \zh{辆车}), et les mots-outils \zh{的、得、地}.
\end{enumerate}
\cle{Connecteurs à retenir} : \zh{先……然后……} (d'abord… ensuite…), \zh{但是 / 可是} (mais), \zh{所以} (donc, conséquence), \zh{因为……所以……} (parce que… donc…), \zh{除了……以外，还……} (à part…, aussi…), \zh{而且} (de plus).
\end{methode}

\begin{exoresolu}[Rédaction guidée : mon argent de poche]
\textbf{Énoncé (type Bac).} Rédigez en chinois un texte de 6 à 8 phrases intitulé « Mon argent de poche » : dites qui vous en donne et combien, comment vous le dépensez, ce que vous économisez, et votre opinion. Utilisez au moins trois connecteurs.
\tcblower
\textbf{Raisonnement.} On suit le plan de la méthode : situation → dépenses → économies → opinion → projet. On place les compléments de temps (\zh{每个月}, \zh{先}) avant le verbe, on emploie \zh{花} pour dépenser de l'argent, \zh{存} pour mettre de côté, et le classificateur \zh{块} (ou \zh{块钱}) pour les sommes.

\textbf{Réponse rédigée.}

\zh{我的零花钱}

\zh{Wǒ de línghuāqián}

\zh{我爸爸妈妈每个月给我五千法郎零花钱。我先买学习用品，比如本子和笔，然后买一点吃的东西。除了买东西以外，我还每个月存一千法郎。但是有时候钱不够，所以我要小心地花钱。我觉得管理零花钱很重要。以后我想用存的钱买一辆自行车。}

\emph{Wǒ bàba māma měi ge yuè gěi wǒ wǔ qiān fǎ láng línghuāqián. Wǒ xiān mǎi xuéxí yòngpǐn, bǐrú běnzi hé bǐ, ránhòu mǎi yìdiǎn chī de dōngxi. Chúle mǎi dōngxi yǐwài, wǒ hái měi ge yuè cún yì qiān fǎ láng. Dànshì yǒu shíhou qián bù gòu, suǒyǐ wǒ yào xiǎoxīn de huā qián. Wǒ juéde guǎnlǐ línghuāqián hěn zhòngyào. Yǐhòu wǒ xiǎng yòng cún de qián mǎi yí liàng zìxíngchē.}

« Mon père et ma mère me donnent chaque mois cinq mille francs d'argent de poche. J'achète d'abord des fournitures scolaires, par exemple des cahiers et des stylos, puis un peu de nourriture. À part les achats, je mets aussi mille francs de côté chaque mois. Mais parfois l'argent ne suffit pas, donc je dois dépenser avec prudence. Je pense que gérer son argent de poche est très important. Plus tard, je voudrais acheter un vélo avec l'argent économisé. »

\textbf{Justifications grammaticales.} \zh{每个月} (complément de temps) est placé avant le verbe \zh{给}. \zh{先……然后……} ordonne les actions chronologiquement. \zh{除了……以外，还……} ajoute une information (épargne) à la liste des dépenses. \zh{一辆自行车} : classificateur \zh{辆} (liàng) pour les véhicules. \zh{小心地花钱} : \zh{地} (de) marque le complément de manière devant le verbe. \zh{存的钱} : \zh{的} (de) relie le verbe \zh{存} au nom \zh{钱} (l'argent \emph{économisé}).

\textbf{Conclusion.} Le texte compte 7 phrases, utilise trois connecteurs (\zh{先……然后……}, \zh{除了……以外，还……}, \zh{所以}) et exprime une opinion personnelle : toutes les consignes sont respectées.
\end{exoresolu}

\courssub{Méthode 2 — Écrire correctement les caractères complexes du budget}

\begin{methode}[Maîtriser les caractères du thème : radicaux, traits, confusions]
\begin{enumerate}
\item \textbf{Identifier le radical (la clé)} : il donne l'indice sémantique. Pour le budget : \zh{钅} (métal/argent) dans \zh{钱} ; \zh{贝} (coquillage = monnaie ancienne) dans \zh{费}, \zh{贵}, \zh{赚}, \zh{账} ; \zh{子} (enfant/objet) dans \zh{存} ; \zh{目} (œil) dans \zh{省}.
\item \textbf{Respecter l'ordre des traits} : de haut en bas (上→下), de gauche à droite (左→右), horizontal avant vertical (横竖), extérieur avant intérieur (外→内), fermer la boîte en dernier (封口最后).
\item \textbf{Distinguer les caractères proches} (paires visuelles ou sémantiques) en les comparant deux par deux (voir tableau).
\item \textbf{S'entraîner} : écrire chaque caractère 5 fois minimum en prononçant le pinyin avec le ton à voix haute.
\item \textbf{Mémoriser par mots entiers (collocations)} et non par caractères isolés : \zh{零花钱} (línghuāqián), \zh{花钱} (huā qián), \zh{存钱} (cún qián), \zh{太贵了} (tài guì le), \zh{记账} (jì zhàng).
\end{enumerate}
\begin{center}
\begin{tabularx}{\linewidth}{|X|X|X|X|}
\hline
\rowcolor{popblueL} Caractère & Pinyin & Radical (Clé) & Point de vigilance / Piège fréquent \\
\hline
钱 & qián & \zh{钅} (métal) & Ne pas confondre avec \zh{线} (xiàn, fil) : le trait du haut diffère (point vs trait horizontal). \\
\hline
费 & fèi & \zh{贝} (monnaie) & Ne pas confondre avec \zh{废} (fèi, inutile/déchet) : \zh{费} a \zh{弗} en haut (phonétique), \zh{废} a \zh{广} (toit). \\
\hline
贵 & guì & \zh{贝} (monnaie) & Visuellement proche de \zh{遗} (yí, perdre) ou \zh{避} (bì, éviter). Sens opposé : \zh{便宜} (piányi, bon marché). \\
\hline
存 & cún & \zh{子} (enfant) & Ne pas confondre avec \zh{在} (zài, être à/préposition) : \zh{存} a \zh{宁} (níng) en haut, \zh{在} a \zh{土} (terre). \\
\hline
省 & shěng & \zh{目} (œil) & Deux sens : 1. économiser (\zh{省钱}) 2. province (\zh{省份}). Ne pas oublier le trait horizontal sous \zh{目}. \\
\hline
账 & zhàng & \zh{贝} (monnaie) & Caractère simplifié standard pour « compte/facture ». Ne pas écrire \zh{帐} (zhàng, tente/toile), forme traditionnelle/alternative. \\
\hline
赚 & zhuàn & \zh{贝} (monnaie) & Phonétique \zh{赟} (yǔn). Ne pas confondre avec \zh{赚} (zuàn, profit/bénéfice — même caractère, ton différent). \\
\hline
\end{tabularx}
\end{center}
\end{methode}

\begin{exoresolu}[Caractères proches : choisir le bon]
\textbf{Énoncé.} Complétez les phrases avec le bon caractère parmi les paires proposées, puis traduisez en français.
\begin{enumerate}
\item 我妈妈给我零花（钱 / 线）。
\item 这本书太（贵 / 遗）了，我不买。
\item 我每个月（存 / 在）五百法郎。
\item 他喜欢（记账 / 记帐），每天写一遍。
\end{enumerate}
\tcblower
\textbf{Raisonnement.} On s'appuie sur le sens (radical \zh{钅} ou \zh{贝} pour l'argent) et sur la forme graphique précise.

\textbf{Solutions.}
\begin{enumerate}
\item \zh{钱} : \zh{我妈妈给我零花钱。} \emph{Wǒ māma gěi wǒ línghuāqián.} « Ma mère me donne de l'argent de poche. » \zh{钱} (qián, argent) a le radical \zh{钅} (métal) ; \zh{线} (xiàn, fil) a le radical \zh{纟} (soie).
\item \zh{贵} : \zh{这本书太贵了，我不买。} \emph{Zhè běn shū tài guì le, wǒ bù mǎi.} « Ce livre est trop cher, je ne l'achète pas. » \zh{贵} (guì, cher) porte le radical \zh{贝} ; \zh{遗} (yí, perdre) a le radical \zh{辶} (marche). L'antonyme est \zh{便宜} (piányi).
\item \zh{存} : \zh{我每个月存五百法郎。} \emph{Wǒ měi ge yuè cún wǔ bǎi fǎ láng.} « Je mets cinq cents francs de côté chaque mois. » \zh{存} (cún, épargner) a \zh{宁} en haut ; \zh{在} (zài, situer) a \zh{土}.
\item \zh{记账} : \zh{他喜欢记账，每天写一遍。} \emph{Tā xǐhuan jì zhàng, měi tiān xiě yí biàn.} « Il aime tenir ses comptes, il écrit chaque jour. » \zh{账} (compte) a le radical \zh{贝} ; \zh{帐} (tente) a le radical \zh{巾} (tissu). En chinois simplifié standard, on utilise \zh{账} pour les comptes.
\end{enumerate}
\textbf{Conclusion.} Le radical \zh{贝} (ancienne monnaie-coquillage) est la clé de presque tous les caractères du budget : le repérer systématiquement aide à ne plus confondre \zh{费、贵、赚、账}.
\end{exoresolu}

\courssub{Méthode 3 — Thème : traduire du français vers le chinois}

\begin{methode}[Réussir le thème sur le budget : 5 étapes]
\begin{enumerate}
\item \textbf{Analyser la phrase française} : identifier le Sujet, le Verbe, les Compléments (temps, lieu, objet, manière).
\item \textbf{Poser le squelette chinois} : \textbf{Sujet + Complément de temps + Verbe + Complément d'objet / Manière / Lieu}. Règle d'or : \cle{le temps précède le verbe}.
\item \textbf{Choisir le verbe chinois précis} :
\begin{itemize}
\item « Donner (à qqn) » : \zh{给 + qqn + objet}.
\item « Dépenser (argent) » : \zh{花钱} / \zh{花 + somme + (在) + objet}.
\item « Acheter » : \zh{买}.
\item « Économiser / Mettre de côté » : \zh{存钱} / \zh{存 + somme}.
\item « Coûter » : \zh{要 + somme} (ex: \zh{要十块钱}) ou adjectif \zh{贵} (cher) / \zh{便宜} (bon marché).
\item « Gagner (par le travail) » : \zh{赚钱}.
\end{itemize}
\item \textbf{Insérer les classificateurs obligatoires} :
\begin{itemize}
\item Somme d'argent : \zh{块} (kuài) ou \zh{块钱} — \zh{十块钱} (10 unités).
\item Objets plats (livres, cahiers) : \zh{本} (běn) — \zh{一本书}.
\item Véhicules (vélo, voiture) : \zh{辆} (liàng) — \zh{一辆自行车}.
\item Objets généraux (stylo, ballon, téléphone) : \zh{个} (gè) / \zh{支} (zhī, pour stylo) — \zh{一个足球}, \zh{一支笔}.
\end{itemize}
\item \textbf{Traiter les structures de degré} : « Trop… » = \zh{太 + adj. + 了} ; « Pas assez » = \zh{不够} (placé après le nom : \zh{钱不够}) ; « De plus en plus » = \zh{越来越 + adj. + 了}.
\item \textbf{Vérifier} : tons (pinyin), ponctuation chinoise pleine chasse （， 。 ？ ！）， mots-outils \zh{的、得、地}.
\end{enumerate}
\end{methode}

\begin{exoresolu}[Thème : trois phrases types Bac]
\textbf{Énoncé.} Traduisez en chinois (caractères + pinyin) :
\begin{enumerate}
\item Mon père me donne de l'argent de poche chaque mois.
\item Ce cahier est trop cher, je n'ai pas assez d'argent.
\item J'économise mille francs pour acheter un ballon de football.
\end{enumerate}
\tcblower
\textbf{Raisonnement.} Pour chaque phrase : 1) Squelette S + Temps + V + O. 2) Choix du verbe. 3) Classificateur. 4) Structure de degré.

\textbf{Solutions.}
\begin{enumerate}
\item \zh{我爸爸每个月给我零花钱。} \emph{Wǒ bàba měi ge yuè gěi wǒ línghuāqián.}
(Sujet \zh{我爸爸} + Temps \zh{每个月} + Verbe \zh{给} + Bénéficiaire \zh{我} + Objet \zh{零花钱}).
\item \zh{这个本子太贵了，我的钱不够。} \emph{Zhè ge běnzi tài guì le, wǒ de qián bù gòu.}
« Ce cahier » = \zh{这个本子} (classif. \zh{本}) ; « trop cher » = \zh{太贵了} ; « mon argent » = \zh{我的钱} ; « ne suffit pas » = \zh{不够} (postposé au sujet \zh{钱}).
\item \zh{我存一千法郎买一个足球。} \emph{Wǒ cún yì qiān fǎ láng mǎi yí ge zúqiú.}
« J'économise » = \zh{存} (verbe direct) ; « mille francs » = \zh{一千法郎} ; « pour acheter » = but exprimé par simple juxtaposition verbale \zh{存……买……} (ou \zh{为了买……}) ; « un ballon » = \zh{一个足球} (classif. \zh{个}).
\end{enumerate}
\textbf{Conclusion.} Trois réflexes conditionnent la réussite : \textbf{temps avant verbe}, \textbf{classificateur devant le nom}, \zh{\textbf{太……了}} pour l'excès.
\end{exoresolu}

\courssub{Méthode 4 — Version : traduire du chinois vers le français}

\begin{methode}[Réussir la version : du chinois au français naturel]
\begin{enumerate}
\item \textbf{Lecture globale} : lire le texte une fois sans traduire pour identifier le thème (ici : budget d'un élève) et le ton.
\item \textbf{Découpage en blocs syntaxiques} : isoler [Sujet] [Temps] [Verbe] [Compléments] ; repérer les \cle{mots-outils} (\zh{的、了、得、地、把、被、比、越来越}).
\item \textbf{Traduction bloc par bloc} : respecter l'ordre français (SVO). Le complément de temps chinois (avant le verbe) se place souvent en début ou fin de phrase en français.
\item \textbf{Rendre les structures figées} :
\begin{itemize}
\item \zh{除了 A 以外，还 B} → « À part A, (il) B aussi / de plus il B ».
\item \zh{越来越 + Adj. + 了} → « De plus en plus + adj. » / « De moins en moins + adj. » (si adj. négatif comme \zh{不够}).
\item \zh{比} → Comparatif « plus… que… ».
\item \zh{把 + Objet + Verbe + 起来/结果} → Action sur l'objet (souvent voix active française).
\item \zh{被} → Passif français.
\end{itemize}
\item \textbf{Reformulation} : éviter le calque mot à mot ; accorder les verbes français (temps, mode) ; choisir le vocabulaire précis du budget (argent de poche, économiser, dépenser, coûter, bon marché).
\item \textbf{Relecture finale} : cohérence, fluidité, orthographe française.
\end{enumerate}
\end{methode}

\begin{exoresolu}[Version : le budget de Xiao Ming (contexte Chine)]
\textbf{Énoncé.} Traduisez en français : \zh{小明是一个中学生。他每个月有五千法郎零花钱。除了买书和笔以外，他还存钱。他觉得钱越来越不够，所以他想买便宜的东西。}
\tcblower
\textbf{Raisonnement.} Phrase 1 : présentation (identité). Phrase 2 : \zh{有} = « disposer de / recevoir » ; complément de temps \zh{每个月} → « chaque mois » en fin de phrase. Phrase 3 : structure additive \zh{除了……以外，还……}. Phrase 4 : \zh{越来越不够} = évolution négative ; \zh{所以} = conséquence ; \zh{便宜} = bon marché.

\textbf{Solution détaillée.}

\zh{小明是一个中学生。} \emph{Xiǎo Míng shì yí ge zhōngxuéshēng.} « Xiao Ming est un collégien (élève du secondaire). » — \zh{个} classificateur générique → article indéfini « un ».

\zh{他每个月有五千法郎零花钱。} \emph{Tā měi ge yuè yǒu wǔ qiān fǎ láng línghuāqián.} « Il dispose de (reçoit) cinq mille francs d'argent de poche chaque mois. » — \zh{有} (yǒu) = avoir/disposer ; \zh{法郎} (fǎ láng) unité monétaire CFA (contexte Cameroun/Zone Franc) ; \zh{零花钱} (línghuāqián) = argent de poche.

\zh{除了买书和笔以外，他还存钱。} \emph{Chúle mǎi shū hé bǐ yǐwài, tā hái cún qián.} « À part l'achat de livres et de stylos, il met aussi de l'argent de côté. » — \zh{除了……以外} (chúle… yǐwài) = à part / hormis ; \zh{还} (hái) = aussi / encore ; \zh{存钱} (cún qián) = épargner.

\zh{他觉得钱越来越不够，所以他想买便宜的东西。} \emph{Tā juéde qián yuè lái yuè bù gòu, suǒyǐ tā xiǎng mǎi piányi de dōngxi.} « Il trouve que l'argent suffit de moins en moins, c'est pourquoi il voudrait acheter des choses bon marché. » — \zh{觉得} (juéde) = trouver/penser ; \zh{越来越不够} (yuè lái yuè bù gòu) = de moins en moins suffisant ; \zh{所以} (suǒyǐ) = donc / c'est pourquoi ; \zh{便宜} (piányi) = bon marché (adj.) ; \zh{东西} (dōngxi) = choses / objets.

\textbf{Conclusion.} La version exige de \textbf{repérer les structures grammaticales figées} (\zh{除了……}, \zh{越来越……}, \zh{所以}) avant de traduire le lexique : elles portent la logique de la phrase.
\end{exoresolu}

\courssub{Méthode 5 — Employer le lexique actif et les collocations du budget}

\begin{methode}[Utiliser les bonnes collocations (associations verbo-nominales)]
En chinois, un verbe se combine avec des noms précis : apprendre les \cle{collocations} évite les calques du français (ex: *\zh{用钱} pour « dépenser », *\zh{得到钱} pour « gagner »).
\begin{center}
\begin{tabularx}{\linewidth}{|X|X|X|X|}
\hline
\rowcolor{popblueL} Collocation & Pinyin & Nature & Traduction \\
\hline
零花钱 & línghuāqián & Nom (composé) & Argent de poche \\
\hline
花钱 & huā qián & Verbe + Compl. & Dépenser de l'argent \\
\hline
存钱 & cún qián & Verbe + Compl. & Mettre de l'argent de côté / Épargner \\
\hline
省钱 & shěng qián & Verbe + Compl. & Économiser (faire des économies, éviter frais) \\
\hline
赚钱 & zhuàn qián & Verbe + Compl. & Gagner de l'argent (par le travail) \\
\hline
借钱 & jiè qián & Verbe + Compl. & Emprunter (\zh{向…借}) / Prêter (\zh{给…借}) de l'argent \\
\hline
太贵了 & tài guì le & Structure degré & C'est trop cher (\zh{了} final obligatoire) \\
\hline
钱不够 & qián bù gòu & Sujet + Verbe & L'argent ne suffit pas / Manquer d'argent \\
\hline
记账 & jì zhàng & Verbe + Compl. & Tenir ses comptes / Noter les dépenses \\
\hline
买单 & mǎi dān & Verbe + Compl. & Payer l'addition (au restaurant) \\
\hline
\end{tabularx}
\end{center}
\cle{Réflexes indispensables} :
\begin{itemize}
\item « Dépenser » = \zh{花钱} (jamais \zh{用钱} seul qui signifie « utiliser de l'argent »).
\item « Gagner » = \zh{赚钱} (par le travail), \zh{中奖} (zhòng jiǎng, gagner à la loterie).
\item « Cher » = \zh{贵} (guì, 4\up{e} ton) ; « Bon marché » = \zh{便宜} (piányi, 2\up{e} + 2\up{e} ton → 2\up{e} + 3\up{e} par sandhi).
\item « Coûter » = \zh{要} + somme (\zh{这要五十块钱}) ou \zh{值} (zhí, valoir).
\end{itemize}
\end{methode}

\begin{exoresolu}[Collocations : compléter et transformer]
\textbf{Énoncé.}
\begin{enumerate}
\item Complétez avec le bon verbe (\zh{花}, \zh{存}, \zh{赚}, \zh{借}) :
\begin{enumerate}
\item 我不想（ ）太多钱买零食。
\item 我哥哥暑假打工，每个月（ ）不少钱。
\item 我没带钱，可以向同学（ ）一点儿钱吗？
\end{enumerate}
\item Transformez la phrase avec la structure \zh{越来越……} : « Les fournitures scolaires deviennent de plus en plus chères. »
\end{enumerate}
\tcblower
\textbf{Raisonnement.} Choix du verbe selon le sens précis : dépenser (\zh{花}), gagner par le travail (\zh{赚}), emprunter (\zh{向…借}), épargner (\zh{存}). Pour la transformation : \zh{越来越 + adjectif + 了}.

\textbf{Solutions.}
\begin{enumerate}
\item a) \zh{花} : \zh{我不想花太多钱买零食。} \emph{Wǒ bù xiǎng huā tài duō qián mǎi língshí.} « Je ne veux pas dépenser trop d'argent pour acheter des en-cas. »
\item b) \zh{赚} : \zh{我哥哥暑假打工，每个月赚不少钱。} \emph{Wǒ gēge shǔjià dǎgōng, měi ge yuè zhuàn bù shǎo qián.} « Mon grand frère fait un job d'été, il gagne pas mal d'argent chaque mois. » — \zh{打工} (dǎgōng) = job étudiant / travail manuel.
\item c) \zh{借} : \zh{我没带钱，可以向同学借一点儿钱吗？} \emph{Wǒ méi dài qián, kěyǐ xiàng tóngxué jiè yìdiǎnr qián ma?} « Je n'ai pas apporté d'argent, puis-je en emprunter un peu à un camarade ? » — \zh{向 + personne + 借} = emprunter à qqn ; \zh{给 + personne + 借} = prêter à qqn.
\item \zh{学习用品越来越贵了。} \emph{Xuéxí yòngpǐn yuè lái yuè guì le.} « Les fournitures scolaires deviennent de plus en plus chères. » — \zh{学习用品} (xuéxí yòngpǐn) = fournitures scolaires ; structure \zh{越来越 + 贵 + 了}.
\end{enumerate}
\textbf{Conclusion.} Les collocations \zh{花钱}, \zh{赚钱}, \zh{借钱}, \zh{存钱}, \zh{省钱} forment le noyau dur du lexique « budget » : les connaître par cœur (caractères + pinyin + sens) garantit un texte naturel et correct.
\end{exoresolu}

\begin{attention}[Les 5 erreurs qui coûtent des points au Bac — Check-list de relecture]
\begin{enumerate}
\item \textbf{Complément de temps mal placé (Erreur n°1).} Le français place « chaque mois » à la fin ; le chinois l'exige \textbf{avant le verbe}.
\begin{quote}
\zh{Correct : 我每个月存钱。} \emph{Wǒ měi ge yuè cún qián.} \\
\zh{Faux : *我存钱每个月。}
\end{quote}
\item \textbf{Classificateur oublié ou incorrect (Erreur n°2).} Tout nom dénombré (somme, objet) requiert son classificateur.
\begin{quote}
Somme : \zh{十块钱} (shí kuài qián) — \textbf{pas} *\zh{十钱}. \\
Véhicule : \zh{一辆自行车} (yí liàng zìxíngchē) — \textbf{pas} *\zh{一自行车}. \\
Livre/Cahier : \zh{一本书} (yí běn shū) / \zh{一个本子} (yí ge běnzi).
\end{quote}
\item \textbf{Caractères confondus par négligence du radical (Erreur n°3).}
\begin{quote}
\zh{钱} (qián, argent, radical \zh{钅}) ≠ \zh{线} (xiàn, fil, radical \zh{纟}). \\
\zh{存} (cún, épargner, haut \zh{宁}) ≠ \zh{在} (zài, être à, haut \zh{土}). \\
\zh{账} (zhàng, compte, radical \zh{贝}) ≠ \zh{帐} (zhàng, tente, radical \zh{巾}). \\
\zh{费} (fèi, frais, radical \zh{贝}) ≠ \zh{废} (fèi, déchet, radical \zh{广}).
\end{quote}
\textbf{Astuce :} surligner le radical \zh{贝} ou \zh{钅} à la relecture.
\item \textbf{Calque lexical sur « dépenser » et « gagner » (Erreur n°4).}
\begin{quote}
\zh{Correct : 花钱} (huā qián), \zh{赚钱} (zhuàn qián). \\
\zh{Faux : *用钱} (utiliser de l'argent), \zh{*得到钱} (obtenir de l'argent), \zh{*挣钱} (zhèng qián — oral/familier, éviter à l'écrit formel).
\end{quote}
Structure « trop + adj. » : \zh{太贵了} (tài guì le) — \textbf{le \zh{了} final est obligatoire} pour marquer l'excès réalisé.
\item \textbf{Tons négligés dans le pinyin (Erreur n°5).} L'inversion de tons change le sens radicalement.
\begin{quote}
\zh{买} (mǎi, 3\up{e} ton, acheter) ≠ \zh{卖} (mài, 4\up{e} ton, vendre). \\
\zh{贵} (guì, 4\up{e} ton, cher) ≠ \zh{跪} (guì, 4\up{e} ton, s'agenouiller — même son, sens diff.). \\
\zh{省} (shěng, 3\up{e} ton, économiser) ≠ \zh{胜} (shèng, 4\up{e} ton, vaincre).
\end{quote}
\textbf{Consigne :} Marquer \textbf{chaque ton} sur \textbf{chaque voyelle} dans toute réponse exigeant le pinyin (ā á ǎ à, ē é ě è, ī í ǐ ì, ō ó ǒ ò, ū ú ǔ ù, ǖ ǘ ǚ ǜ). Pas de chiffres (1.2,3.4).
\end{enumerate}
\end{attention}

\begin{exercice}[Entraînement Bac — Expression écrite et traduction]
\textbf{Exercice 1 (Rédaction — 10 pts).} Rédigez en chinois (caractères + pinyin) un texte de 7 à 9 phrases intitulé « \zh{我的理财小计划} » (Mon petit plan de gestion financière). Contraintes :
\begin{itemize}
\item Présentez votre source de revenus (parents, job d'été, Nouvel An…).
\item Détaillez deux postes de dépenses (ex: \zh{学习用品}, \zh{吃饭}, \zh{娱乐}).
\item Indiquez votre montant d'épargne mensuel et où vous le mettez (ex: \zh{存银行}, \zh{存钱罐}).
\item Exprimez votre opinion sur l'importance de \zh{记账} (tenir ses comptes).
\item Utilisez \textbf{au moins 4 connecteurs} parmi : \zh{先……然后……}, \zh{除了……以外……}, \zh{但是}, \zh{所以}, \zh{因为……所以……}, \zh{而且}, \zh{如果……就……}.
\end{itemize}

\textbf{Exercice 2 (Thème — 5 pts).} Traduisez en chinois (caractères + pinyin) :
\begin{enumerate}
\item Ma mère pense que je dépense trop d'argent en jeux vidéo.
\item Ce téléphone portable est très bon marché, il ne coûte que mille francs.
\item Si j'économise chaque mois, l'année prochaine je pourrai m'acheter un ordinateur.
\end{enumerate}

\textbf{Exercice 3 (Version — 5 pts).} Traduisez en français :
\zh{王红是一个高中生。她每个星期六在超市打工，赚两百块钱。她不花这些钱，全存起来。她说：“存钱很有意思，将来我想去北京旅游。”}
\end{exercice}

\begin{corrige}[Exercice 1 — Proposition de rédaction type]
\textbf{Texte :}
\zh{我的理财小计划}

\zh{Wǒ de lǐcái xiǎo jìhuà}

\zh{我每个月的零花钱是爸爸妈妈给的，暑假我也打工赚一点钱。先买学习用品，然后请同学吃饭。除了这些花费以外，我每个月存两千法郎在存钱罐里。但是有时候我想买新衣服，钱就不够了。所以我觉得记账很有必要，能知道钱花哪儿了。而且如果我坚持存钱，明年就能买一台电脑。以后我想用自己的钱去北京旅游。}

\emph{Wǒ měi ge yuè de línghuāqián shì bàba māma gěi de, shǔjià wǒ yě dǎgōng zhuàn yìdiǎn qián. Xiān mǎi xuéxí yòngpǐn, ránhòu qǐng tóngxué chī fàn. Chúle zhè xiē huàfèi yǐwài, wǒ měi ge yuè cún liǎng qiān fǎ láng zài cúnqiánguàn lǐ. Dànshì yǒu shíhou wǒ xiǎng mǎi xīn yīfu, qián jiù bù gòu le. Suǒyǐ wǒ juéde jì zhàng hěn yǒu bìyào, néng zhīdào qián huā nǎr le. Érqiě rúguǒ wǒ jiānchí cún qián, míngnián jiù néng mǎi yí tái diànnǎo. Yǐhòu wǒ xiǎng yòng zìjǐ de qián qù Běijīng lǚyóu.}

« Mon argent de poche mensuel vient de mes parents, et pendant les vacances d'été je fais aussi des petits boulots pour gagner un peu d'argent. D'abord j'achète des fournitures scolaires, puis j'invite des camarades à manger. À part ces dépenses, je mets deux mille francs de côté chaque mois dans ma tirelire. Mais parfois je veux acheter de nouveaux vêtements, et l'argent ne suffit plus. Donc je trouve que tenir ses comptes est très nécessaire, ça permet de savoir où l'argent est parti. De plus, si je persévère à économiser, l'année prochaine je pourrai m'acheter un ordinateur. Plus tard, je voudrais aller voyager à Pékin avec mon propre argent. »
\textbf{Connecteurs utilisés :} \zh{先……然后……}, \zh{除了……以外}, \zh{但是}, \zh{所以}, \zh{而且}, \zh{如果……就……} (6 connecteurs). \textbf{Structures :} \zh{把钱存起来} (omise ici pour simplicité mais \zh{存……在……} correct), \zh{记账}, \zh{太……了} (implicite dans \zh{不够了}), \zh{越来越} (non utilisé mais \zh{坚持} utilisé).
\end{corrige}

\begin{corrige}[Exercice 2 — Thème]
\begin{enumerate}
\item \zh{我妈妈觉得我在电子游戏上花太多钱了。} \emph{Wǒ māma juéde wǒ zài diànzǐ yóuxì shàng huā tài duō qián le.} (Structure \zh{在……上花钱} / \zh{觉得} + proposition).
\item \zh{这个手机很便宜，只要一千法郎。} \emph{Zhè ge shǒujī hěn piányi, zhǐ yào yì qiān fǎ láng.} (\zh{手机} téléphone portable ; \zh{只要} « ne coûter que / seulement coûter »).
\item \zh{如果我每个月存钱，明年我就能买一台电脑。} \emph{Rúguǒ wǒ měi ge yuè cún qián, míngnián wǒ jiù néng mǎi yí tái diànnǎo.} (\zh{如果……就……} structure conditionnelle ; \zh{台} classif. pour ordinateur/machine ; \zh{能} capacité/possibilité).
\end{enumerate}
\end{corrige}

\begin{corrige}[Exercice 3 — Version]
\zh{王红是一个高中生。} \emph{Wáng Hóng shì yí ge gāozhōngshēng.} « Wang Hong est une lycéenne. »

\zh{她每个星期六在超市打工，赚两百块钱。} \emph{Tā měi ge xīngqīliù zài chāoshì dǎgōng, zhuàn liǎng bǎi kuài qián.} « Elle fait un job étudiant au supermarché chaque samedi, et gagne deux cents yuans (unités monétaires). » — \zh{星期六} (xīngqīliù) samedi ; \zh{超市} (chāoshì) supermarché ; \zh{打工} (dǎgōng) job/boulot ; \zh{块钱} (kuài qián) unité de base (yuan/kuai).

\zh{她不花这些钱，全存起来。} \emph{Tā bù huā zhè xiē qián, quán cún qǐlai.} « Elle ne dépense pas cet argent, elle le met tout de côté. » — \zh{这些} (zhè xiē) ces ; \zh{全} (quán) tout/entièrement ; \zh{存起来} (cún qǐlai) mettre de côté (complément de résultat \zh{起来}).

\zh{她说：“存钱很有意思，将来我想去北京旅游。”} \emph{Tā shuō: “Cún qián hěn yǒu yìsi, jiānglái wǒ xiǎng qù Běijīng lǚyóu.”} « Elle dit : “Économiser est très intéressant / amusant, plus tard je veux aller voyager à Pékin.” » — \zh{有意思} (yǒu yìsi) intéressant/amusant ; \zh{将来} (jiānglái) dans le futur/plus tard ; \zh{北京} (Běijīng) Pékin ; \zh{旅游} (lǚyóu) voyager.
\end{corrige}

\newpage

\coursec{Exercices}

\courssub{Je vérifie mes connaissances}

\begin{exercice}[QCM : le vocabulaire du budget]
Pour chaque question, entoure la bonne réponse.
\begin{enumerate}
\item \zh{零花钱} (línghuāqián) signifie :
\begin{itemize}
\item[A.] le salaire
\item[B.] l'argent de poche
\item[C.] la monnaie rendue
\end{itemize}
\item Pour dire « économiser de l'argent », on utilise :
\begin{itemize}
\item[A.] \zh{省钱} (shěng qián)
\item[B.] \zh{买钱} (mǎi qián)
\item[C.] \zh{问钱} (wèn qián)
\end{itemize}
\item \zh{预算} (yùsuàn) signifie :
\begin{itemize}
\item[A.] le prix
\item[B.] le cadeau
\item[C.] le budget
\end{itemize}
\item Quel connecteur signifie « d'abord » ?
\begin{itemize}
\item[A.] \zh{然后} (ránhòu)
\item[B.] \zh{首先} (shǒuxiān)
\item[C.] \zh{但是} (dànshì)
\end{itemize}
\end{enumerate}
\end{exercice}

\begin{exercice}[Vrai ou faux justifié]
Indique si chaque affirmation est vraie (V) ou fausse (F), puis justifie ta réponse en français.
\begin{enumerate}
\item \zh{贵} (guì) signifie « bon marché ».
\item Dans \zh{因为……所以……} (yīnwèi... suǒyǐ...), les deux mots peuvent s'utiliser ensemble dans la même phrase en chinois.
\item \zh{储蓄} (chǔxù) signifie « l'épargne ».
\item Le caractère \zh{钱} (qián) contient le radical de la soie \zh{纟}.
\end{enumerate}
\end{exercice}

\begin{exercice}[Vocabulaire : complète la traduction]
Complète le tableau en écrivant la traduction française de chaque mot.
\begin{center}
\begin{tabularx}{\linewidth}{|X|X|X|}\hline
\rowcolor{popblueL} Caractères & Pinyin & Traduction \\ \hline
\zh{管理} & guǎnlǐ & \ldots \\ \hline
\zh{花费} & huāfèi & \ldots \\ \hline
\zh{便宜} & piányi & \ldots \\ \hline
\zh{打算} & dǎsuàn & \ldots \\ \hline
\zh{浪费} & làngfèi & \ldots \\ \hline
\end{tabularx}
\end{center}
\end{exercice}

\begin{exercice}[Associe les caractères au pinyin]
Relie chaque caractère de la colonne A à son pinyin de la colonne B.
\begin{center}
\begin{tabularx}{\linewidth}{|X|X|}\hline
\rowcolor{popblueL} A — Caractère & B — Pinyin \\ \hline
1. \zh{钱} & a. suàn \\ \hline
2. \zh{费} & b. chǔ \\ \hline
3. \zh{算} & c. qián \\ \hline
4. \zh{储} & d. fèi \\ \hline
5. \zh{贵} & e. guì \\ \hline
\end{tabularx}
\end{center}
\end{exercice}

\courssub{Je m'entraîne}

\begin{exercice}[Traduction : français vers chinois]
Traduis les phrases suivantes en chinois (caractères et pinyin).
\begin{enumerate}
\item Je reçois cent yuans d'argent de poche chaque mois.
\item Parce que ce cahier est trop cher, je ne l'achète pas.
\item D'abord je paie mes fournitures scolaires, ensuite je mets le reste de côté.
\item Je pense que gérer son argent de poche est très important.
\end{enumerate}
\end{exercice}

\begin{exercice}[Texte à trous : mots-outils et connecteurs]
Complète le texte suivant avec les mots de la liste : \zh{首先}、\zh{然后}、\zh{因为}、\zh{所以}、\zh{但是}、\zh{的}、\zh{了} (le mot \zh{了} sert deux fois).

\begin{textebox}[Mon budget]
我每个月有一百块零花钱。\_\_\_\_，我先买学习用品，\_\_\_\_ 再存钱。\_\_\_\_ 我想买一本汉语词典，\_\_\_\_ 我不买零食。\_\_\_\_ 有时候东西太贵\_\_\_\_，我就不买那些贵\_\_\_\_ 东西\_\_\_\_。这样，我每个月都能存一点儿钱。\\
(Wǒ měi ge yuè yǒu yìbǎi kuài línghuāqián. \_\_\_\_, wǒ xiān mǎi xuéxí yòngpǐn, \_\_\_\_ zài cún qián. \_\_\_\_ wǒ xiǎng mǎi yì běn Hànyǔ cídiǎn, \_\_\_\_ wǒ bù mǎi língshí. \_\_\_\_ yǒu shíhou dōngxi tài guì \_\_\_\_, wǒ jiù bù mǎi nàxiē guì \_\_\_\_ dōngxi \_\_\_\_. Zhèyàng, wǒ měi ge yuè dōu néng cún yìdiǎnr qián.)\\
« Chaque mois, je reçois cent yuans d'argent de poche. \ldots, j'achète d'abord mes fournitures scolaires, \ldots je mets de l'argent de côté. \ldots je veux acheter un dictionnaire de chinois, \ldots je n'achète pas de friandises. \ldots parfois les choses sont trop chères \ldots, alors je n'achète plus ces choses chères \ldots. Ainsi, chaque mois, j'arrive à mettre un peu d'argent de côté. »
\end{textebox}
\end{exercice}

\begin{exercice}[Remise en ordre de mots]
Remets les mots dans le bon ordre pour former des phrases correctes, puis traduis-les en français.
\begin{enumerate}
\item \zh{零花钱 / 我 / 管理 / 会 / 自己的}
\item \zh{比 / 这个 / 那个 / 贵 / 一点儿}
\item \zh{因为 / 所以 / 太贵 / 我 / 不买 / 它 / ，}
\item \zh{存 / 想 / 钱 / 我 / 买 / 一个 / 足球 / 新 / 的}
\end{enumerate}
\end{exercice}

\begin{exercice}[买 ou 卖 ?]
Complète chaque phrase avec \zh{买} (mǎi, acheter) ou \zh{卖} (mài, vendre), puis traduis.
\begin{enumerate}
\item 我想\_\_\_\_一个书包，我的旧书包坏了。
\item 市场上的阿姨\_\_\_\_水果和蔬菜。
\item 这本书太贵了，我不\_\_\_\_了。
\item 他把自己的旧手机\_\_\_\_了，然后用这些钱\_\_\_\_了一本词典。
\end{enumerate}
\end{exercice}

\courssub{Je me prépare au Bac}

\begin{exercice}[Compréhension de l'écrit et questions de langue]
Lis le texte suivant, puis réponds aux questions.

\begin{textebox}[李明的零花钱]
李明是成都的一个高中生。他每个月有二百块零花钱。首先，他用这些钱买学习用品和早饭。然后，他把剩下的钱存起来，因为他想买一台电脑。他的同学王华每个月都把零花钱花完，但是李明觉得这样不好。他说：« 我们应该学会管理自己的钱。 » 所以，李明每个月都写一个小预算：多少钱用来吃饭，多少钱用来买东西，多少钱存起来。他的妈妈很高兴，她说李明是一个会计划的孩子。\\
(Lǐ Míng shì Chéngdū de yí ge gāozhōngshēng. Tā měi ge yuè yǒu èrbǎi kuài línghuāqián. Shǒuxiān, tā yòng zhèxiē qián mǎi xuéxí yòngpǐn hé zǎofàn. Ránhòu, tā bǎ shèngxià de qián cún qǐlai, yīnwèi tā xiǎng mǎi yì tái diànnǎo. Tā de tóngxué Wáng Huá měi ge yuè dōu bǎ línghuāqián huā wán, dànshì Lǐ Míng juéde zhèyàng bù hǎo. Tā shuō : « Wǒmen yīnggāi xuéhuì guǎnlǐ zìjǐ de qián. » Suǒyǐ, Lǐ Míng měi ge yuè dōu xiě yí ge xiǎo yùsuàn : duōshao qián yònglái chīfàn, duōshao qián yònglái mǎi dōngxi, duōshao qián cún qǐlai. Tā de māma hěn gāoxìng, tā shuō Lǐ Míng shì yí ge huì jìhuà de háizi.)\\
« Li Ming est un lycéen de Chengdu. Chaque mois, il reçoit deux cents yuans d'argent de poche. D'abord, il utilise cet argent pour acheter ses fournitures scolaires et son petit-déjeuner. Ensuite, il met de côté l'argent qui reste, car il veut acheter un ordinateur. Son camarade Wang Hua, lui, dépense tout son argent de poche chaque mois, mais Li Ming pense que ce n'est pas bien. Il dit : "Nous devons apprendre à gérer notre argent." C'est pourquoi Li Ming rédige chaque mois un petit budget : combien pour manger, combien pour acheter des choses, combien à épargner. Sa mère est très contente, elle dit que Li Ming est un enfant qui sait planifier. »
\end{textebox}

\textbf{Questions de compréhension} (réponds en français) :
\begin{enumerate}
\item Combien d'argent de poche Li Ming reçoit-il chaque mois ?
\item Pourquoi met-il de l'argent de côté ?
\item Que fait Li Ming chaque mois pour bien gérer son argent ?
\item Pourquoi sa mère est-elle contente ?
\end{enumerate}

\textbf{Questions de langue} :
\begin{enumerate}
\item Releve dans le texte deux connecteurs et explique leur emploi.
\item Explique la structure de la phrase \zh{他把剩下的钱存起来} (construction avec \zh{把}).
\item Trouve dans le texte un mot de la même famille que \zh{花} et donne son sens.
\end{enumerate}
\end{exercice}

\begin{exercice}[Thème et version]
\textbf{A. Thème.} Traduis en chinois (caractères et pinyin) les quatre phrases suivantes. Tu utiliseras obligatoirement les connecteurs indiqués entre parenthèses.
\begin{enumerate}
\item D'abord, je fais mon budget, ensuite j'achète ce dont j'ai besoin. (\zh{首先}……\zh{然后}……)
\item Parce que je veux économiser de l'argent, je n'achète pas de choses inutiles. (\zh{因为}……\zh{所以}……)
\item Je pense que l'argent de poche des lycéens camerounais n'est pas beaucoup. (\zh{我觉得})
\item Mon petit frère dépense toujours tout son argent de poche, moi, je mets la moitié de côté.
\end{enumerate}

\textbf{B. Version.} Traduis en français le texte suivant.

\begin{textebox}[小芳的储蓄]
小芳是北京的一个学生。她的爸爸妈妈每个星期给她五十块零花钱。小芳不乱花钱：她先买需要的东西，然后把剩下的钱存进银行。半年以后，她存了五百块。她用这些钱给妈妈买了一个生日礼物。妈妈非常感动。\\
(Xiǎo Fāng shì Běijīng de yí ge xuésheng. Tā de bàba māma měi ge xīngqī gěi tā wǔshí kuài línghuāqián. Xiǎo Fāng bù luàn huā qián : tā xiān mǎi xūyào de dōngxi, ránhòu bǎ shèngxià de qián cún jìn yínháng. Bàn nián yǐhòu, tā cún le wǔbǎi kuài. Tā yòng zhèxiē qián gěi māma mǎi le yí ge shēngrì lǐwù. Māma fēicháng gǎndòng.)
\end{textebox}
\end{exercice}

\begin{exercice}[Expression écrite guidée]
\textbf{Sujet :} \zh{我怎么管理我的零花钱} (Wǒ zěnme guǎnlǐ wǒ de línghuāqián) — « Comment je gère mon argent de poche ».

Rédige un texte en chinois d'\textbf{au moins 8 phrases} (environ 80 à 100 caractères) dans lequel tu expliques :
\begin{itemize}
\item combien d'argent de poche tu reçois et de qui ;
\item comment tu dépenses cet argent (fournitures, nourriture, transport) ;
\item si tu économises et pourquoi ;
\item ton opinion sur la gestion de l'argent de poche.
\end{itemize}

\textbf{Contraintes obligatoires :}
\begin{enumerate}
\item Utilise au moins \textbf{4 connecteurs} différents parmi : \zh{首先}、\zh{然后}、\zh{因为}……\zh{所以}……、\zh{但是}、\zh{我觉得}、\zh{最后}。
\item Utilise au moins \textbf{5 mots du glossaire} de la leçon : \zh{零花钱}、\zh{预算}、\zh{储蓄}、\zh{便宜}、\zh{贵}、\zh{管理}、\zh{浪费}、\zh{打算}。
\item Écris les caractères avec soin, en respectant l'ordre des traits ; souligne les caractères \zh{钱}、\zh{费}、\zh{算}、\zh{储} quand tu les emploies.
\end{enumerate}

\textbf{Grille d'évaluation indicative :} respect du sujet et de la longueur (4 pts), exactitude des caractères (4 pts), grammaire et ordre des mots (4 pts), connecteurs et vocabulaire du thème (4 pts), présentation générale (4 pts).
\end{exercice}



\newpage

\coursec{Corrigés des exercices}

\courssub{Corrigés de « Je vérifie mes connaissances »}

\begin{corrige}[Exercice 1 — QCM : le vocabulaire du budget]
\begin{enumerate}
\item \textbf{B} — \zh{零花钱} (línghuāqián) = l'argent de poche. Le salaire se dit \zh{工资} (gōngzī) ; la monnaie rendue se dit \zh{找钱} (zhǎoqián).
\item \textbf{A} — \zh{省钱} (shěngqián) = économiser de l'argent. \zh{买钱} et \zh{问钱} n'existent pas : \zh{买} (mǎi) = acheter, \zh{问} (wèn) = demander.
\item \textbf{C} — \zh{预算} (yùsuàn) = le budget. Le prix se dit \zh{价格} (jiàgé) ; le cadeau se dit \zh{礼物} (lǐwù).
\item \textbf{B} — \zh{首先} (shǒuxiān) = d'abord ; \zh{然后} (ránhòu) = ensuite ; \zh{但是} (dànshì) = mais.
\end{enumerate}
\end{corrige}

\begin{corrige}[Exercice 2 — Vrai ou faux justifié]
\begin{enumerate}
\item \textbf{F} — \zh{贵} (guì) signifie « cher ». « Bon marché » se dit \zh{便宜} (piányi). Attention à ce faux ami fréquent à l'examen.
\item \textbf{V} — Contrairement au français (« parce que » et « donc » ne se combinent pas), le chinois emploie volontiers les deux ensemble : \zh{因为太贵，所以我不买。} (Yīnwèi tài guì, suǒyǐ wǒ bù mǎi.) « Parce que c'est trop cher, je n'achète pas. »
\item \textbf{V} — \zh{储蓄} (chǔxù) = l'épargne, le fait de mettre de l'argent de côté.
\item \textbf{F} — \zh{钱} contient le radical du \textbf{métal} \zh{钅} (jīn), logique : la monnaie était autrefois en métal. Le radical de la soie \zh{纟} (sī) se trouve par exemple dans \zh{红} (hóng, rouge).
\end{enumerate}
\end{corrige}

\begin{corrige}[Exercice 3 — Vocabulaire : complète la traduction]
\zh{管理} (guǎnlǐ) = gérer, administrer ; \zh{花费} (huāfèi) = dépenser / la dépense ; \zh{便宜} (piányi) = bon marché ; \zh{打算} (dǎsuàn) = avoir l'intention de, prévoir ; \zh{浪费} (làngfèi) = gaspiller.\\
Point de vigilance : \zh{便宜} se prononce \emph{piányi} (et non biànyí) dans le sens « bon marché » ; \zh{花费} peut être verbe ou nom selon le contexte.
\end{corrige}

\begin{corrige}[Exercice 4 — Associe les caractères au pinyin]
1-c (\zh{钱} qián) ; 2-d (\zh{费} fèi) ; 3-a (\zh{算} suàn) ; 4-b (\zh{储} chǔ) ; 5-e (\zh{贵} guì).\\
Astuce mémoire : \zh{钱} a le radical du métal \zh{钅} ; \zh{费} et \zh{贵} contiennent tous deux la clé \zh{贝} (bèi, le coquillage-monnaie), comme beaucoup de caractères liés à l'argent ; \zh{算} porte la clé du bambou \zh{⺮} (zhú, on comptait avec des baguettes de bambou).
\end{corrige}

\courssub{Corrigés de « Je m'entraîne »}

\begin{corrige}[Exercice 5 — Traduction : français vers chinois]
\begin{enumerate}
\item \zh{我每个月有一百块零花钱。} (Wǒ měige yuè yǒu yìbǎi kuài línghuāqián.) — Le complément de temps \zh{每个月} se place juste après le sujet ; \zh{块} (kuài) est le classificateur de la monnaie courante.
\item \zh{因为这个本子太贵，所以我不买。} (Yīnwèi zhège běnzi tài guì, suǒyǐ wǒ bù mǎi.) — « Cahier » = \zh{本子} (běnzi) ; on peut aussi dire \zh{练习本} (liànxíběn).
\item \zh{首先我买学习用品，然后把剩下的钱存起来。} (Shǒuxiān wǒ mǎi xuéxí yòngpǐn, ránhòu bǎ shèngxià de qián cúnqǐlai.) — « Fournitures scolaires » = \zh{学习用品} (xuéxí yòngpǐn) ; structure \zh{把} + objet + \zh{存起来} (cúnqǐlai, mettre de côté).
\item \zh{我觉得管理零花钱很重要。} (Wǒ juéde guǎnlǐ línghuāqián hěn zhòngyào.) — Le verbe \zh{管理} (guǎnlǐ) sert ici de sujet de la subordonnée, sans aucun changement de forme.
\end{enumerate}
\end{corrige}

\begin{corrige}[Exercice 6 — Texte à trous : mots-outils et connecteurs]
\zh{我每个月有一百块零花钱。首先，我先买学习用品，然后再存钱。因为我想买一本汉语词典，所以我不买零食。但是有时候东西太贵了，我就不买那些贵的东西了。}\\
Points de langue : \zh{太贵了} (tài guì le) exprime l'excès ; \zh{贵的} (guì de) est l'adjectif \zh{贵} suivi du mot-outil \zh{的} devant le nom \zh{东西} ; le second \zh{了} (le) marque le changement de situation (« je ne les achète plus »).
\end{corrige}

\begin{corrige}[Exercice 7 — Remise en ordre de mots]
\begin{enumerate}
\item \zh{我会管理自己的零花钱。} (Wǒ huì guǎnlǐ zìjǐ de línghuāqián.) « Je sais gérer mon argent de poche. »
\item \zh{这个比那个贵一点儿。} (Zhège bǐ nàge guì yìdiǎnr.) « Celui-ci est un peu plus cher que celui-là. » Structure : A + \zh{比} (bǐ) + B + adjectif (+ \zh{一点儿}).
\item \zh{因为它太贵，所以我不买。} (Yīnwèi tā tài guì, suǒyǐ wǒ bù mǎi.) « Parce que c'est trop cher, je ne l'achète pas. »
\item \zh{我想存钱买一个新的足球。} (Wǒ xiǎng cúnqián mǎi yí ge xīn de zúqiú.) « Je veux économiser pour acheter un ballon de football neuf. »
\end{enumerate}
\end{corrige}

\begin{corrige}[Exercice 8 — 买 ou 卖 ?]
\begin{enumerate}
\item \zh{买} (mǎi) — « Je veux acheter un cartable, mon vieux cartable est cassé. »
\item \zh{卖} (mài) — « La dame du marché vend des fruits et des légumes. »
\item \zh{买} (mǎi) — « Ce livre est trop cher, je ne l'achète plus. »
\item \zh{卖} (mài) puis \zh{买} (mǎi) — « Il a vendu son vieux téléphone, puis a acheté un dictionnaire avec cet argent. »
\end{enumerate}
Moyen mnémotechnique : \zh{卖} (mài, 4ᵉ ton) porte en haut le composant \zh{十} (shí) ; on vend « à dix personnes », on achète pour soi. Attention aux tons : \zh{买} mǎi (3ᵉ ton) = acheter, \zh{卖} mài (4ᵉ ton) = vendre.
\end{corrige}

\courssub{Corrigés de « Je me prépare au Bac »}

\begin{corrige}[Exercice 9 — Compréhension de l'écrit et questions de langue]
\textbf{Compréhension :}
\begin{enumerate}
\item Il reçoit deux cents yuans (\zh{二百块} èrbǎi kuài) par mois.
\item Il met de l'argent de côté parce qu'il veut acheter un ordinateur (\zh{他想买一台电脑} Tā xiǎng mǎi yí tái diànnǎo).
\item Chaque mois, il rédige un petit budget : il décide combien d'argent sert à manger, combien sert à acheter des choses et combien est épargné.
\item Sa mère est contente parce que Li Ming est un enfant qui sait planifier (\zh{会计划的孩子} huì jìhuà de háizi).
\end{enumerate}
\textbf{Langue :}
\begin{enumerate}
\item Exemples : \zh{首先} (shǒuxiān, d'abord) et \zh{然后} (ránhòu, ensuite) ordonnent les actions ; \zh{因为} (yīnwèi) introduit la cause ; \zh{但是} (dànshì, mais) introduit l'opposition ; \zh{所以} (suǒyǐ) introduit la conséquence.
\item Structure avec \zh{把} (bǎ) : Sujet + \zh{把} + complément d'objet + verbe + complément de résultat/direction. Ici : \zh{他} (tā, sujet) + \zh{把} + \zh{剩下的钱} (shèngxià de qián, objet : l'argent restant) + \zh{存起来} (cúnqǐlai, mettre de côté). Le verbe ne peut jamais rester seul après l'objet dans une phrase en \zh{把}.
\item \zh{花完} (huāwán) = finir de dépenser, tout dépenser (complément de résultat) ; \zh{花} (huā) seul signifie « dépenser » (et aussi « fleur »).
\end{enumerate}
Barème indicatif : compréhension 4 x 1 pt ; questions de langue 2 + 2 + 2 pts.
\end{corrige}

\begin{corrige}[Exercice 10 — Thème et version]
\textbf{A. Thème (propositions de corrigé) :}
\begin{enumerate}
\item \zh{首先，我做预算，然后买我需要的东西。} (Shǒuxiān, wǒ zuò yùsuàn, ránhòu mǎi wǒ xūyào de dōngxi.)
\item \zh{因为我想省钱，所以我不买没有用的东西。} (Yīnwèi wǒ xiǎng shěngqián, suǒyǐ wǒ bù mǎi méiyǒu yòng de dōngxi.)
\item \zh{我觉得喀麦隆高中生的零花钱不多。} (Wǒ juéde Kāmàilóng gāozhōngshēng de línghuāqián bù duō.)
\item \zh{我弟弟总是把零花钱花完，我存一半。} (Wǒ dìdi zǒngshì bǎ línghuāqián huāwán, wǒ cún yíbàn.)
\end{enumerate}
\textbf{B. Version (corrigé) :}\\
« Xiao Fang est une élève de Pékin. Ses papa et maman lui donnent cinquante yuans d'argent de poche chaque semaine. Xiao Fang ne dépense pas son argent n'importe comment : elle achète d'abord les choses dont elle a besoin, puis dépose l'argent restant à la banque. Six mois plus tard, elle avait économisé cinq cents yuans. Avec cet argent, elle a acheté un cadeau d'anniversaire pour sa maman. Sa maman a été très émue. »\\
Points de vigilance : \zh{乱花钱} (luàn huāqián) = dépenser n'importe comment, gaspiller ; \zh{存进银行} (cún jìn yínháng) = déposer à la banque ; \zh{半年以后} (bàn nián yǐhòu) = six mois plus tard ; \zh{感动} (gǎndòng) = ému, touché. Barème indicatif : thème 4 x 2 pts ; version 8 pts (fidélité 4, qualité du français 2, vocabulaire 2).
\end{corrige}

\begin{corrige}[Exercice 11 — Expression écrite guidée : modèle de production]
\begin{textebox}[Modèle de rédaction]
我的爸爸妈妈每个月给我一百块零花钱。首先，我做一个小预算。然后，我用这些钱买学习用品和早饭。因为我想买一本汉语词典，所以我不浪费钱，也不买太贵的东西。最后，我把剩下的钱存起来。我觉得管理零花钱很重要：会管理钱的孩子，长大以后也会计划自己的生活。\\
(Wǒ de bàba māma měige yuè gěi wǒ yìbǎi kuài línghuāqián. Shǒuxiān, wǒ zuò yíge xiǎo yùsuàn. Ránhòu, wǒ yòng zhèxiē qián mǎi xuéxí yòngpǐn hé zǎofàn. Yīnwèi wǒ xiǎng mǎi yì běn Hànyǔ cídiǎn, suǒyǐ wǒ bù làngfèi qián, yě bù mǎi tài guì de dōngxi. Zuìhòu, wǒ bǎ shèngxià de qián cúnqǐlai. Wǒ juéde guǎnlǐ línghuāqián hěn zhòngyào: huì guǎnlǐ qián de háizi, zhǎngdà yǐhòu yě huì jìhuà zìjǐ de shēnghuó.)\\
« Mes papa et maman me donnent cent yuans d'argent de poche chaque mois. D'abord, je fais un petit budget. Ensuite, j'utilise cet argent pour acheter mes fournitures scolaires et mon petit-déjeuner. Parce que je veux acheter un dictionnaire de chinois, je ne gaspille pas mon argent et je n'achète pas de choses trop chères. Enfin, je mets de côté l'argent qui reste. Je pense que gérer son argent de poche est très important : un enfant qui sait gérer son argent saura aussi planifier sa vie une fois adulte. »
\end{textebox}
Ce modèle respecte les contraintes : 5 connecteurs (\zh{首先}、\zh{然后}、\zh{因为}……\zh{所以}……、\zh{最后}、\zh{我觉得}) et 6 mots du glossaire (\zh{零花钱}、\zh{预算}、\zh{浪费}、\zh{贵}、\zh{管理}、\zh{存}). Adapte-le à ta situation : montant réel, contexte de Yaoundé ou de Douala, projet d'achat personnel (ballon de football, téléphone, livre).\\
Barème indicatif : respect du sujet et de la longueur (4 pts), exactitude des caractères (4 pts), grammaire et ordre des mots (4 pts), connecteurs et vocabulaire du thème (4 pts), présentation générale (4 pts). Points de vigilance : place du complément de temps après le sujet, \zh{的} (de) obligatoire entre adjectif et nom, tons de \zh{买} (mǎi) / \zh{卖} (mài), ne pas oublier \zh{了} (le) dans \zh{太……了} (tài… le).
\end{corrige}

\newpage

\coursec{Fiche bilan}

\begin{popfigure}
\begin{center}\tcbox[colback=popblue,colframe=popblue,coltext=white,arc=9pt,boxsep=4pt,left=6pt,right=6pt]{\bfseries\small Gérer son argent de poche \zh{零花钱管理}}\end{center}
\vspace{-2pt}
\begin{tcbraster}[raster columns=3,raster equal height=rows,raster column skip=6pt,raster row skip=6pt,enhanced,blankest]
\begin{tcolorbox}[enhanced,colback=poporange,colframe=poporange,coltext=white,boxrule=0.9pt,arc=3pt,left=4pt,right=4pt,top=3pt,bottom=3pt,boxsep=1pt,halign=left]\bfseries\scriptsize L'exemple \zh{例子} un texte modèle\end{tcolorbox}
\begin{tcolorbox}[enhanced,colback=popgreen,colframe=popgreen,coltext=white,boxrule=0.9pt,arc=3pt,left=4pt,right=4pt,top=3pt,bottom=3pt,boxsep=1pt,halign=left]\bfseries\scriptsize Structure début -- corps -- fin\end{tcolorbox}
\begin{tcolorbox}[enhanced,colback=poppurple,colframe=poppurple,coltext=white,boxrule=0.9pt,arc=3pt,left=4pt,right=4pt,top=3pt,bottom=3pt,boxsep=1pt,halign=left]\bfseries\scriptsize Connecteurs \zh{首先、然后、} \zh{但是、所以}\end{tcolorbox}
\begin{tcolorbox}[enhanced,colback=poppink,colframe=poppink,coltext=white,boxrule=0.9pt,arc=3pt,left=4pt,right=4pt,top=3pt,bottom=3pt,boxsep=1pt,halign=left]\bfseries\scriptsize Lexique \zh{钱、花、省、} \zh{存、预算}\end{tcolorbox}
\begin{tcolorbox}[enhanced,colback=popgold,colframe=popgold,coltext=white,boxrule=0.9pt,arc=3pt,left=4pt,right=4pt,top=3pt,bottom=3pt,boxsep=1pt,halign=left]\bfseries\scriptsize Caractères traits, radicaux \zh{贝} = argent\end{tcolorbox}
\begin{tcolorbox}[enhanced,colback=popblue!60,colframe=popblue!60,coltext=white,boxrule=0.9pt,arc=3pt,left=4pt,right=4pt,top=3pt,bottom=3pt,boxsep=1pt,halign=left]\bfseries\scriptsize Thème français $\rightarrow$ chinois\end{tcolorbox}
\begin{tcolorbox}[enhanced,colback=popgreen!60,colframe=popgreen!60,coltext=white,boxrule=0.9pt,arc=3pt,left=4pt,right=4pt,top=3pt,bottom=3pt,boxsep=1pt,halign=left]\bfseries\scriptsize Version chinois $\rightarrow$ français\end{tcolorbox}
\end{tcbraster}

\legende{Carte mentale de la leçon : écrire sur la gestion de l'argent de poche}
\end{popfigure}

\begin{bilan}
\textbf{1. Lexique clé à connaître par cœur}
\begin{itemize}
\item \zh{零花钱} \emph{línghuāqián} : argent de poche ; \zh{钱} \emph{qián} : argent ; \zh{块} \emph{kuài} : yuan (classificateur de la monnaie).
\item \zh{花} \emph{huā} : dépenser ; \zh{省} \emph{shěng} : économiser ; \zh{存} \emph{cún} : épargner, déposer ; \zh{借} \emph{jiè} : prêter / emprunter.
\item \zh{预算} \emph{yùsuàn} : budget ; \zh{需要} \emph{xūyào} : avoir besoin de ; \zh{贵 / 便宜} \emph{guì / piányi} : cher / bon marché.
\item Collocations : \zh{花钱买东西} (dépenser pour acheter), \zh{存钱} (épargner), \zh{省钱} (faire des économies), \zh{做预算} (faire un budget).
\end{itemize}

\textbf{2. Règles de grammaire et structures}

\begin{center}
\begin{tabularx}{\linewidth}{|X|X|X|}
\hline
\rowcolor{popblueL} Structure & Exemple (caractères + pinyin) & Traduction \\
\hline
Sujet + 把 + complément + verbe & \zh{我把钱存起来。} \emph{Wǒ bǎ qián cún qǐlai.} & Je mets l'argent de côté. \\
\hline
verbe + 得 + adverbe (manière) & \zh{他钱花得太多。} \emph{Tā qián huā de tài duō.} & Il dépense trop d'argent. \\
\hline
A + 比 + B + adjectif & \zh{书比衣服便宜。} \emph{Shū bǐ yīfu piányi.} & Les livres sont moins chers que les vêtements. \\
\hline
先 … 然后 … & \zh{我先做预算，然后花钱。} \emph{Wǒ xiān zuò yùsuàn, ránhòu huā qián.} & Je fais d'abord un budget, puis je dépense. \\
\hline
因为 … 所以 … & \zh{因为东西很贵，所以我要省钱。} \emph{Yīnwèi dōngxi hěn guì, suǒyǐ wǒ yào shěng qián.} & Parce que les choses sont chères, je dois économiser. \\
\hline
\end{tabularx}
\end{center}

\textbf{3. Écriture des caractères du thème}
\begin{itemize}
\item Le radical \zh{贝} \emph{bèi} (coquillage, ancienne monnaie) signale le sens « argent » : \zh{贵、赚、贴、账}. On l'écrit : vertical, horizontal-crochet, puis les deux points intérieurs.
\item \zh{钱} : radical métal \zh{钅} à gauche (5 traits), \zh{戋} à droite ; gauche avant droite.
\item \zh{省} : \zh{少} en haut, \zh{目} en bas ; du haut vers le bas. Ne pas confondre avec \zh{雀}.
\item Caractères proches à ne pas confondre : \zh{花} (dépenser / fleur) et \zh{华} ; \zh{存} et \zh{在} ; \zh{借} et \zh{错}.
\end{itemize}

\textbf{4. Méthodes express}
\begin{itemize}
\item \textbf{Rédiger} : 3 paragraphes — situation (d'où vient l'argent), gestion (budget, dépenses, économies), conclusion (avis personnel). Relier avec \zh{首先、然后、但是、所以、最后}.
\item \textbf{Thème} : repérer sujet + temps + verbe ; placer le complément de temps après le sujet ; vérifier les classificateurs (\zh{一百块钱}).
\item \textbf{Version} : traduire d'abord le sens global, puis ajuster en français naturel ; attention à \zh{得} (manière) et \zh{把} (disposition).
\end{itemize}
\end{bilan}

\begin{aretenir}[Checklist avant le Bac]
\begin{itemize}
\item $\square$ Je connais le vocabulaire du budget : \zh{零花钱、预算、花、省、存、借、贵、便宜}.
\item $\square$ J'écris \zh{钱} avec le radical \zh{钅} sans oublier de traits.
\item $\square$ Je reconnais le radical \zh{贝} et son lien avec l'argent.
\item $\square$ Je distingue \zh{存 / 在} et \zh{借 / 错} à l'écrit.
\item $\square$ J'utilise correctement \zh{把} : Sujet + 把 + complément + verbe.
\item $\square$ Je place le complément de temps après le sujet : \zh{我每个月存钱。}
\item $\square$ J'emploie les connecteurs \zh{首先、然后、但是、所以、最后} dans mes textes.
\item $\square$ Je fais la comparaison avec \zh{比} : A + 比 + B + adjectif.
\item $\square$ Je mets le bon classificateur : \zh{块钱}、\zh{本书}、\zh{件衣服}.
\item $\square$ Mon texte a trois parties : introduction, développement, conclusion.
\item $\square$ En version, je traduis le sens avant de polir le français.
\item $\square$ Je relis mes tons et ma ponctuation chinoise （，。）avant de rendre la copie.
\end{itemize}
\end{aretenir}

\findecours

\end{document}
